% Dissertation rédiger une dissertation — Français Tle Tech — Cours signé OnBuch+
% Programme officiel MINESEC. Compiler avec : tectonic N08-dissertation-rediger-dissertation.tex
\newcommand{\DOCMATIERE}{Français}
\newcommand{\DOCNIVEAU}{Tle Tech}
\newcommand{\DOCMODULE}{Français – classe de Terminale technique}
\newcommand{\DOCLECON}{N08}
\newcommand{\DOCTITRE}{Dissertation rédiger une dissertation}
% OnBuch+ — préambule « cours signés » ENRICHI (compilé avec tectonic / XeLaTeX)
% Système visuel « Pop » d'OnBuch+ : fond crème (jamais blanc), bordures encre
% épaisses, ombres « dures » décalées, titres Archivo Black en capitales.
% Version MATHÉMATIQUES : graphes (pgfplots/tikz), boîtes pédagogiques
% (définition, propriété, méthode, attention, le savais-tu, exercice résolu,
% exercice, TP, bilan), page de garde et signature OnBuch+ ancrée en bas.
%
% Macros attendues dans main.tex AVANT \input{preamble} :
%   \DOCMATIERE  \DOCNIVEAU  \DOCMODULE  \DOCLECON (numéro)  \DOCTITRE
\documentclass[11pt]{article}

\usepackage{fontspec}
\usepackage[french]{babel}
\usepackage[a4paper,margin=2cm,top=3.4cm,bottom=2.8cm,headheight=1.4cm,footskip=1.3cm]{geometry}
\usepackage{amsmath,amssymb,amsfonts}
\usepackage{mathtools} % \xmapsto, \coloneqq... souvent utilisés par les modèles
\usepackage{cancel} % simplifications barrées (bilans de forces)
% \abs{z} (module, valeur absolue) et \norm{u} (norme), utilisés par les modèles
\providecommand{\abs}{}\let\abs\relax\DeclarePairedDelimiter{\abs}{\lvert}{\rvert}
\providecommand{\norm}{}\let\norm\relax\DeclarePairedDelimiter{\norm}{\lVert}{\rVert}
\usepackage[table]{xcolor} % [table] : fournit aussi \rowcolors (alternance de lignes)
\usepackage{pagecolor}
\usepackage{tikz}
\usetikzlibrary{arrows.meta,positioning,calc,shapes.geometric,shapes.misc,shapes.arrows,decorations.pathmorphing,decorations.pathreplacing,decorations.markings,patterns,shadows,mindmap,trees,fit,backgrounds,matrix,intersections,angles,quotes}
\usepackage{pgfplots}
\usepackage[version=4]{mhchem} % équations nucléaires \ce{^{235}_{92}U}
\usepackage[european,siunitx]{circuitikz} % schémas électriques
\pgfplotsset{compat=1.18}
\usepgfplotslibrary{fillbetween,groupplots} % aires entre courbes (calcul intégral)
\usepackage{enumitem}
\usepackage{fancyhdr}
% [most] charge toutes les bibliothèques tcolorbox (listings, minted, raster,
% documentation...) et gonfle inutilement la mémoire TeX sur un document long
% (~300 boîtes) : on ne charge que ce qui est réellement utilisé.
\usepackage[skins,breakable]{tcolorbox}
\usepackage{graphicx}
\usepackage{colortbl}
\usepackage{booktabs}
\usepackage{tabularx}
\usepackage{multirow}
\usepackage{diagbox} % case d'en-tête diagonale des tableaux à double entrée
% Colonnes extensibles centrée / gauche / droite (C, L, R), souvent utilisées par les modèles
\newcolumntype{C}{>{\centering\arraybackslash}X}
\newcolumntype{L}{>{\raggedright\arraybackslash}X}
\newcolumntype{R}{>{\raggedleft\arraybackslash}X}
\usepackage{array}
\usepackage{multicol}
\usepackage{siunitx}
% parse-numbers=false : accepte aussi \SI{2,5 \times 10^{-3}}{...}
\sisetup{locale=FR,output-decimal-marker={,},group-minimum-digits=4,parse-numbers=false}
\DeclareSIUnit\ohm{\ensuremath{\symup{\Omega}}} % U+2126 absent de la police
\DeclareSIUnit\division{div} % réglages d'oscilloscope (ms/div, V/div)
\DeclareSIUnit\tr{tr} % tours (vitesse de rotation, ex. \SI{1500}{\tr\per\minute})
\DeclareSIUnit\molar{M} % concentration molaire (ex. \SI{3}{\molar}, \SI{1}{\micro\molar})
\DeclareSIUnit\an{an} % année (le modèle confond parfois avec \year, un compteur TeX) — ex. \SI{5}{\kilo\meter\per\an}
\DeclareSIUnit\FCFA{FCFA} % franc CFA (ex. \SI{500}{\FCFA\per\kilo\gram})
\DeclareSIUnit\degreeBrix{\textdegree Brix} % teneur en sucre d'un sirop/jus (réfractométrie)
\DeclareSIUnit\month{mois} % le modèle utilise parfois \month, un compteur TeX, comme unité de durée
\DeclareSIUnit\microliter{\micro\liter} % le modèle écrit parfois \microliter en un seul token
\DeclareSIUnit\million{millions} % dénombrement (ex. \SI{15}{\million\per\mL} spermatozoïdes)
\usepackage{qrcode}
% \degree : commande fréquemment utilisée par les modèles pour un angle ou une
% température en dehors de \SI{}{} (non fournie par les paquets chargés).
\providecommand{\degree}{^{\circ}}
% \qed (fin de démonstration), utilisé par les modèles sans amsthm
\providecommand{\qed}{\hfill$\square$}
% \barre : double barre des valeurs interdites dans un tableau de variations (tkz-tab)
\providecommand{\barre}{\ensuremath{\|}}
% environnement proof (amsthm non chargé) utilisé par les modèles
\ifdefined\proof\else\newenvironment{proof}[1][Démonstration]{\par\noindent\textit{#1.}\ }{\qed\par}\fi
\usepackage{lastpage}
\usepackage{hyperref}
\hypersetup{hidelinks}

% ============================================================
% Palette « Pop » OnBuch+ — mêmes hex que la classe P (plus_ui.dart)
% ============================================================
\definecolor{poporange}{HTML}{F59321}
\definecolor{popdark}{HTML}{E07A0C}
\definecolor{popcream}{HTML}{FFF6E8}
\definecolor{popink}{HTML}{1C1714}
\definecolor{poppurple}{HTML}{7A5AE0}
\definecolor{popgreen}{HTML}{2E9E6A}
\definecolor{poppink}{HTML}{E0457B}
\definecolor{popblue}{HTML}{2F7DE1}
\definecolor{popgold}{HTML}{C98A12}
\definecolor{popmuted}{HTML}{8A7F76}
\definecolor{popline}{HTML}{EADFD2}
\definecolor{popgray}{HTML}{8A7F76}
\colorlet{popgrey}{popgray}
% Alias tolérants (le modèle nomme parfois les couleurs différemment)
\colorlet{popwhite}{white}
\colorlet{popblack}{popink}
\colorlet{popred}{poppink}
\colorlet{popyellow}{popgold}
% Teintes claires (fonds de tableaux/encadrés) — le modèle nomme parfois
% les couleurs avec un suffixe L (ex. popblueL) sans les avoir définies.
\colorlet{poporangeL}{poporange!20}
\colorlet{popdarkL}{popdark!20}
\colorlet{poppurpleL}{poppurple!20}
\colorlet{popgreenL}{popgreen!20}
\colorlet{poppinkL}{poppink!20}
\colorlet{popblueL}{popblue!20}
\colorlet{popgoldL}{popgold!20}
\colorlet{popredL}{popred!20}
\colorlet{popgrayL}{popgray!20}
% Teintes claires pour fonds de tableaux / zones
\colorlet{popblueL}{popblue!10!white}
\colorlet{poporangeL}{poporange!14!white}
\colorlet{popgreenL}{popgreen!12!white}
\colorlet{poppurpleL}{poppurple!12!white}
\colorlet{poppinkL}{poppink!12!white}
\colorlet{popgoldL}{popgold!14!white}

\pagecolor{popcream}

% ============================================================
% Typographie
% ============================================================
\defaultfontfeatures{Ligatures=TeX}
\setmainfont{PlusJakartaSans-Medium.ttf}[Path=../../../fonts/,BoldFont=PlusJakartaSans-Bold.ttf]
% Maths en sans-serif (Fira Math) avec chiffres et lettres droites de Plus
% Jakarta, pour que les formules se fondent dans le texte.
\usepackage[math-style=ISO,bold-style=ISO]{unicode-math}
\setmathfont{FiraMath-Regular.otf}[Path=../../../fonts/]
\setmathfont{PlusJakartaSans-Medium.ttf}[Path=../../../fonts/,range={up/num,up/{Latin,latin},\mathup}]
\usepackage{icomma} % virgule décimale sans espace en mode math
% Caractères mathématiques Unicode absents des polices de texte (Plus Jakarta,
% Archivo Black) : rendus par leur équivalent mathématique, en texte comme en
% formule — sinon ils s'affichent en carré vide (ex. « Arithmétique dans ℤ »).
\usepackage{newunicodechar}
\newunicodechar{ℝ}{\ensuremath{\mathbb{R}}}
\newunicodechar{ℤ}{\ensuremath{\mathbb{Z}}}
\newunicodechar{ℂ}{\ensuremath{\mathbb{C}}}
\newunicodechar{ℕ}{\ensuremath{\mathbb{N}}}
\newunicodechar{ℚ}{\ensuremath{\mathbb{Q}}}
\newunicodechar{𝔻}{\ensuremath{\mathbb{D}}}
\newunicodechar{α}{\ensuremath{\alpha}}
\newunicodechar{β}{\ensuremath{\beta}}
\newunicodechar{γ}{\ensuremath{\gamma}}
\newunicodechar{δ}{\ensuremath{\delta}}
\newunicodechar{ε}{\ensuremath{\varepsilon}}
\newunicodechar{θ}{\ensuremath{\theta}}
\newunicodechar{λ}{\ensuremath{\lambda}}
\newunicodechar{μ}{\ensuremath{\mu}}
\newunicodechar{σ}{\ensuremath{\sigma}}
\newunicodechar{τ}{\ensuremath{\tau}}
\newunicodechar{φ}{\ensuremath{\varphi}}
\newunicodechar{ψ}{\ensuremath{\psi}}
\newunicodechar{ω}{\ensuremath{\omega}}
\newunicodechar{Σ}{\ensuremath{\Sigma}}
% majuscules produites par \MakeUppercase dans les titres (α-aminés -> Α)
\newunicodechar{Α}{\ensuremath{\alpha}}
\newunicodechar{Β}{\ensuremath{\beta}}
\newunicodechar{Γ}{\ensuremath{\Gamma}}
\newunicodechar{Δ}{\ensuremath{\Delta}}
\newunicodechar{Ε}{\ensuremath{\varepsilon}}
\newunicodechar{Θ}{\ensuremath{\Theta}}
\newunicodechar{Λ}{\ensuremath{\Lambda}}
\newunicodechar{Μ}{\ensuremath{\mu}}
\newunicodechar{Τ}{\ensuremath{\tau}}
\newunicodechar{Φ}{\ensuremath{\Phi}}
\newunicodechar{Ψ}{\ensuremath{\Psi}}
\newunicodechar{Ω}{\ensuremath{\Omega}}
\newunicodechar{↦}{\ensuremath{\mapsto}}
\newunicodechar{∈}{\ensuremath{\in}}
\newunicodechar{∉}{\ensuremath{\notin}}
\newunicodechar{∧}{\ensuremath{\wedge}}
\newunicodechar{⇔}{\ensuremath{\Leftrightarrow}}
\newunicodechar{⟺}{\ensuremath{\Longleftrightarrow}}
\newunicodechar{⇒}{\ensuremath{\Rightarrow}}
\newunicodechar{⟹}{\ensuremath{\Longrightarrow}}
\newunicodechar{∘}{\ensuremath{\circ}}
\newunicodechar{∪}{\ensuremath{\cup}}
\newunicodechar{∩}{\ensuremath{\cap}}
\newunicodechar{≡}{\ensuremath{\equiv}}
\newunicodechar{⊂}{\ensuremath{\subset}}
\newunicodechar{⊥}{\ensuremath{\perp}}
\newunicodechar{∃}{\ensuremath{\exists}}
\newunicodechar{∀}{\ensuremath{\forall}}
\newunicodechar{∛}{\ensuremath{\sqrt[3]{\,}}}
\newunicodechar{∜}{\ensuremath{\sqrt[4]{\,}}}
\newunicodechar{⃗}{\ensuremath{{}^{\to}}}
\newunicodechar{ⁿ}{\ifmmode^{n}\else\textsuperscript{\ensuremath{n}}\fi}
\newunicodechar{ˣ}{\ifmmode^{x}\else\textsuperscript{\ensuremath{x}}\fi}
\newunicodechar{ᵇ}{\ifmmode^{b}\else\textsuperscript{\ensuremath{b}}\fi}
\newunicodechar{ʳ}{\ifmmode^{r}\else\textsuperscript{\ensuremath{r}}\fi}
\newunicodechar{ᵘ}{\ifmmode^{u}\else\textsuperscript{\ensuremath{u}}\fi}
\newunicodechar{ʸ}{\ifmmode^{y}\else\textsuperscript{\ensuremath{y}}\fi}
\newunicodechar{ᵃ}{\ifmmode^{a}\else\textsuperscript{\ensuremath{a}}\fi}
\newunicodechar{ᵉ}{\ifmmode^{e}\else\textsuperscript{\ensuremath{e}}\fi}
\newunicodechar{ᵏ}{\ifmmode^{k}\else\textsuperscript{\ensuremath{k}}\fi}
\newunicodechar{ᵖ}{\ifmmode^{p}\else\textsuperscript{\ensuremath{p}}\fi}
\newunicodechar{ᴺ}{\ifmmode^{N}\else\textsuperscript{\ensuremath{N}}\fi}
\newunicodechar{ᶜ}{\ifmmode^{c}\else\textsuperscript{\ensuremath{c}}\fi}
\newunicodechar{ʰ}{\ifmmode^{h}\else\textsuperscript{\ensuremath{h}}\fi}
\newunicodechar{ᵠ}{\ifmmode^{q}\else\textsuperscript{\ensuremath{q}}\fi}
\newunicodechar{ₙ}{\ifmmode_{n}\else\textsubscript{\ensuremath{n}}\fi}
\newunicodechar{ₐ}{\ifmmode_{a}\else\textsubscript{\ensuremath{a}}\fi}
\newunicodechar{ᵢ}{\ifmmode_{i}\else\textsubscript{\ensuremath{i}}\fi}
\newunicodechar{ᵣ}{\ifmmode_{r}\else\textsubscript{\ensuremath{r}}\fi}
\newunicodechar{ₖ}{\ifmmode_{k}\else\textsubscript{\ensuremath{k}}\fi}
\newunicodechar{ᵦ}{\ifmmode_{\beta}\else\textsubscript{\ensuremath{\beta}}\fi}
\newunicodechar{ₓ}{\ifmmode_{x}\else\textsubscript{\ensuremath{x}}\fi}
\newfontfamily\popdisplay{ArchivoBlack-Regular.ttf}[Path=../../../fonts/]
\newfontfamily\poplogo{SpaceGrotesk-Bold.ttf}[Path=../../../fonts/]
\newfontfamily\popsemi{PlusJakartaSans-SemiBold.ttf}[Path=../../../fonts/]
\newfontfamily\popxbold{PlusJakartaSans-ExtraBold.ttf}[Path=../../../fonts/]
\newfontfamily\popxboldls{PlusJakartaSans-ExtraBold.ttf}[Path=../../../fonts/,LetterSpace=3.0]

\setlist[enumerate]{leftmargin=1.7em,itemsep=5pt,topsep=4pt}
\setlist[itemize]{leftmargin=1.7em,itemsep=4pt,topsep=4pt,label={\color{poporange}\textbullet}}
\setlength{\parindent}{0pt}
\setlength{\parskip}{5pt}
\renewcommand{\arraystretch}{1.25}

% pgfplots : style Pop par défaut
\pgfplotsset{
  every axis/.append style={axis line style={popink,line width=1pt},tick style={popink},
    grid style={popline,line width=0.5pt},label style={font=\small},tick label style={font=\footnotesize}},
  popcurve/.style={poporange,line width=1.8pt,no markers,smooth},
}
% Même style disponible dans un \draw TikZ simple (hors pgfplots)
\tikzset{popcurve/.style={poporange,line width=1.8pt}}
\tikzset{popgrid/.style={popline,very thin}}
\colorlet{popcurve}{poporange} % le nom du style est parfois utilisé comme couleur

% ============================================================
% Pill / squiggle
% ============================================================
\newcommand{\poppill}[3]{%
  \tikz[baseline=-0.55ex]\node[draw=popink,line width=1.2pt,rounded corners=2.6pt,%
    fill=#2,inner xsep=6.5pt,inner ysep=3pt]{%
    \popxboldls\fontsize{7}{7}\selectfont\color{#3}\MakeUppercase{#1}};%
}
\newcommand{\popsquiggle}[1]{%
  \tikz[baseline=-0.4ex]{
    \draw[#1,line width=2.6pt,line cap=round]
      plot [smooth] coordinates {(0,0.09) (0.34,0.01) (0.68,0.21) (1.02,0.01) (1.36,0.10)};
  }%
}
% Mot-clé surligné (marqueur orange sous le texte)
\newcommand{\cle}[1]{{\popxbold\color{popdark}#1}}

% ============================================================
% En-tête / pied
% ============================================================
\pagestyle{fancy}
\fancyhf{}
\renewcommand{\headrulewidth}{0pt}
\renewcommand{\footrulewidth}{0pt}
\lhead{\raisebox{-0.15ex}{\poplogo\fontsize{13}{13}\selectfont\color{popink}OnBuch\color{poporange}+}}
\rhead{\poppill{\DOCMATIERE\ \textbullet\ \DOCNIVEAU\ \textbullet\ Leçon \DOCLECON}{white}{popink}}
\lfoot{\popxbold\fontsize{7.6}{7.6}\selectfont\color{popmuted}\MakeUppercase{Cours signé OnBuch+}\ \textbullet\ {\popsemi\DOCTITRE}}
\rfoot{\tikz[baseline=-0.6ex]\node[rounded corners=8.5pt,fill=popink,minimum height=17pt,minimum width=17pt,inner xsep=5pt,inner sep=0]{\popxbold\fontsize{8}{8}\selectfont\color{popcream}\thepage\,/\,\pageref*{LastPage}};}

% ============================================================
% Titres
% ============================================================
\newcommand{\chaptitre}[2]{%
  \begin{tcolorbox}[enhanced,colback=poporange,colframe=popink,boxrule=2.6pt,arc=11pt,
      drop shadow=popink,left=15pt,right=15pt,top=12pt,bottom=11pt,before skip=3pt,after skip=18pt]
    \poppill{#2}{popink}{popcream}\par\vspace{9pt}
    {\raggedright\hyphenpenalty=10000\exhyphenpenalty=10000\popdisplay\fontsize{22}{24}\selectfont\color{popcream}\MakeUppercase{#1}\par}
  \end{tcolorbox}
}
% Section numérotée : pastille orange avec le numéro + titre Archivo
\newcounter{popsec}
\newcommand{\coursec}[1]{%
  \stepcounter{popsec}\setcounter{popsub}{0}%
  \par\vspace{16pt}\noindent
  \begin{minipage}{\linewidth}
  \tikz[baseline=-0.7ex]\node[draw=popink,line width=1.6pt,fill=poporange,rounded corners=4pt,minimum size=22pt,inner sep=0]{\popdisplay\fontsize{12}{12}\selectfont\color{popcream}\thepopsec};%
  \hspace{8pt}{\popdisplay\fontsize{15}{17}\selectfont\color{popink}\MakeUppercase{#1}}\par
  \vspace{4pt}\noindent{\color{poporange}\rule{36pt}{3.6pt}}
  \end{minipage}\par\nopagebreak\vspace{8pt}
}
\newcounter{popsub}
\newcommand{\courssub}[1]{%
  \stepcounter{popsub}%
  \par\vspace{9pt}\noindent{\popxbold\fontsize{11.5}{13}\selectfont\color{popink}{\color{poporange}\thepopsec.\thepopsub}\hspace{6pt}#1}\par\nopagebreak\vspace{4pt}
}

% ============================================================
% Boîtes pédagogiques — même recette : fond blanc, bordure épaisse colorée,
% ombre dure encre, badge plein. Titre optionnel [#1].
% ============================================================
\tcbset{popbase/.style={breakable,enhanced,colback=white,boxrule=2.3pt,arc=9pt,
  shadow={3pt}{-3pt}{0pt}{popink},left=14pt,right=14pt,top=11pt,bottom=11pt,before skip=11pt,after skip=15pt,
  fontupper=\popsemi\fontsize{10.6}{14.5}\selectfont\color{popink},
  fontlower=\popsemi\fontsize{10.6}{14.5}\selectfont\color{popink}}}
% Coche dessinée (le glyphe ✓ n'existe pas dans les polices du document)
\newcommand{\popcheck}[1]{\tikz[baseline=-0.5ex]\draw[#1,line width=1.6pt,line cap=round,line join=round](-0.13,0.0)--(-0.04,-0.09)--(0.13,0.1);}
\providecommand{\coche}{\popcheck{popgreen}}
\providecommand{\resultat}[1]{\par\smallskip\noindent\fcolorbox{popgreen}{popgreenL}{\parbox{\dimexpr\linewidth-2\fboxsep-2\fboxrule\relax}{\textbf{Résultat.} #1}}\par\smallskip}
\newcommand{\popboxhead}[3]{\poppill{#1}{#2}{white}\ifx\relax#3\relax\else\hspace{6pt}{\popxbold\fontsize{10.6}{12}\selectfont\color{popink}#3}\fi\par\vspace{7pt}}

\newtcolorbox{aretenir}[1][]{popbase,colframe=popgreen,before upper={\popboxhead{À retenir}{popgreen}{#1}}}
\newtcolorbox{exemplebox}[1][]{popbase,colframe=poporange,before upper={\popboxhead{Exemple}{poporange}{#1}}}
\newtcolorbox{definition}[1][]{popbase,colframe=popblue,before upper={\popboxhead{Définition}{popblue}{#1}}}
\newtcolorbox{propriete}[1][]{popbase,colframe=poppurple,before upper={\popboxhead{Propriété}{poppurple}{#1}}}
\newtcolorbox{methode}[1][]{popbase,colframe=popgold,before upper={\popboxhead{Méthode}{popgold}{#1}}}
\newtcolorbox{attention}[1][]{popbase,colframe=poppink,before upper={\popboxhead{Attention}{poppink}{#1}}}
\newtcolorbox{experience}[1][]{popbase,colframe=popblue,colback=popblueL,before upper={\popboxhead{Expérience}{popblue}{#1}}}
% « Le savais-tu ? » : carte violette pleine (contexte, culture, Cameroun)
\newtcolorbox{savaistu}[1][]{popbase,colback=poppurple,colframe=popink,
  fontupper=\popsemi\fontsize{10.4}{14.2}\selectfont\color{white},
  before upper={\poppill{Le savais-tu\,?}{white}{poppurple}\ifx\relax#1\relax\else\hspace{6pt}{\popxbold\color{white}#1}\fi\par\vspace{7pt}}}
% Exercice résolu : énoncé en haut, \tcblower puis la solution sur fond crème
\newcounter{popexo}\newcounter{popexr}
\newtcolorbox{exoresolu}[1][]{popbase,colframe=poporange,colbacklower=popcream,
  segmentation style={popink,line width=1.2pt,dashed},
  before upper={\stepcounter{popexr}\popboxhead{Exercice résolu \thepopexr}{poporange}{#1}},
  before lower={\poppill{Solution}{popink}{popcream}\par\vspace{6pt}}}
% Exercice (non résolu en place) — la correction est dans la partie Corrigés
\newtcolorbox{exercice}[1][]{popbase,colframe=popink,
  before upper={\stepcounter{popexo}\popboxhead{Exercice \thepopexo}{popink}{#1}}}
% Corrigé
\newtcolorbox{corrige}[1][]{popbase,colframe=popgreen,colback=popgreenL,
  before upper={\popboxhead{Corrigé}{popgreen}{#1}}}
% Objectifs / prérequis / bilan
\newtcolorbox{objectifs}{popbase,colframe=popink,colback=poporangeL,
  before upper={\popboxhead{Objectifs de la leçon}{popink}{}}}
\newtcolorbox{prerequis}{popbase,colframe=popink,colback=popblueL,
  before upper={\popboxhead{Prérequis}{popblue}{}}}
\newtcolorbox{bilan}{popbase,colframe=popink,colback=white,boxrule=2.8pt,
  before upper={\popboxhead{Fiche bilan}{poporange}{L'essentiel à connaître pour l'examen}}}
% Figure encadrée avec légende
\newtcolorbox{popfigure}[1][]{enhanced,breakable=false,colback=white,colframe=popline,boxrule=1.4pt,arc=8pt,
  left=10pt,right=10pt,top=10pt,bottom=8pt,before skip=10pt,after skip=12pt,halign=center,
  lower separated=false,halign lower=center,
  fontlower=\popsemi\fontsize{9}{11.5}\selectfont\color{popmuted},
  #1}
\newcommand{\legende}[1]{\par\vspace{6pt}{\popsemi\fontsize{9}{11.5}\selectfont\color{popmuted}#1}}

% ============================================================
% Signature OnBuch+
% ============================================================
\newcommand{\poplogoinline}[1]{{\poplogo\fontsize{#1}{#1}\selectfont\color{popink}OnBuch\color{poporange}+}}
\newcommand{\signatureonbuch}{%
  \begin{tcolorbox}[enhanced,colback=popink,colframe=popink,boxrule=2.2pt,arc=10pt,
      drop shadow=poporange,left=14pt,right=14pt,top=10pt,bottom=10pt,before skip=0pt,after skip=0pt]
    \tikz[baseline=-0.7ex]\node[circle,fill=popgreen,draw=popcream,line width=1.3pt,minimum size=22pt,inner sep=0]{\popcheck{white}};%
    \hspace{9pt}\begin{minipage}[c]{0.62\linewidth}
      {\popxbold\fontsize{12}{13}\selectfont\color{popcream}Signé\ \textbullet\ {\poplogo OnBuch\color{poporange}+}}\par\vspace{2pt}
      {\raggedright\popsemi\fontsize{8}{10.5}\selectfont\color{popline}\MakeUppercase{Équipe pédagogique OnBuch+}\\\MakeUppercase{Conforme au programme officiel MINESEC}\par}
    \end{minipage}\hfill
    {\poplogo\fontsize{22}{22}\selectfont\color{popcream}OnBuch\color{poporange}+}
  \end{tcolorbox}%
}

% ============================================================
% Page de garde
% #1 titre · #2 matière · #3 niveau · #4 module · #5 n° leçon · #6 durée ·
% #7 description / objectifs (texte ou itemize)
% ============================================================
\newcommand{\pagegarde}[7]{%
  \thispagestyle{empty}
  \newgeometry{a4paper,margin=1.7cm,top=1.6cm,bottom=1.4cm}%
  \begin{tikzpicture}[remember picture,overlay]
    \fill[poporange!18] ([xshift=-3.2cm,yshift=-3.6cm]current page.north east) circle (4.6cm);
    \fill[poppurple!14] ([xshift=2.4cm,yshift=7.6cm]current page.south west) circle (3.8cm);
  \end{tikzpicture}%
  {\poplogo\fontsize{30}{30}\selectfont\color{popink}OnBuch\color{poporange}+}\hfill
  \raisebox{0.6ex}{\poppill{Cours signé \textbullet\ Premium}{popink}{popcream}}\par
  \vspace{1pt}\popsquiggle{poppurple}\par
  \vspace{8pt}
  \begin{tcolorbox}[enhanced,colback=poporange,colframe=popink,boxrule=2.8pt,arc=13pt,
      drop shadow=popink,left=18pt,right=18pt,top=12pt,bottom=13pt,after skip=10pt]
    \poppill{#2 \textbullet\ #3}{popink}{popcream}\hspace{5pt}\poppill{#4}{white}{popink}\par\vspace{8pt}
    {\popdisplay\fontsize{14}{14}\selectfont\color{popink}LEÇON #5}\par\vspace{4pt}
    {\raggedright\hyphenpenalty=10000\exhyphenpenalty=10000\popdisplay\fontsize{25}{27}\selectfont\color{popcream}\MakeUppercase{#1}\par}
    \vspace{8pt}
    {\popxbold\fontsize{9.5}{9.5}\selectfont\color{popink}\MakeUppercase{Durée conseillée : #6}}
  \end{tcolorbox}
  \begin{tcolorbox}[enhanced,colback=white,colframe=popink,boxrule=2.2pt,arc=10pt,
      drop shadow=popink,left=16pt,right=16pt,top=10pt,bottom=10pt,after skip=8pt]
    \poppill{Dans cette leçon}{poppurple}{white}\par\vspace{5pt}
    \popsemi\fontsize{9.3}{12.6}\selectfont\color{popink}#7
  \end{tcolorbox}
  \noindent\begin{minipage}[t]{0.487\linewidth}
    \begin{tcolorbox}[enhanced,colback=popgreen,colframe=popink,boxrule=2.2pt,arc=10pt,
        drop shadow=popink,left=11pt,right=11pt,top=10pt,bottom=10pt,height=3.1cm,valign=center]
      \tcbox[colback=white,colframe=popink,boxrule=1.3pt,arc=5pt,boxsep=0pt,nobeforeafter,
        left=3pt,right=3pt,top=3pt,bottom=3pt,valign=center]{\qrcode[height=1.9cm]{https://play.google.com/store/apps/details?id=com.luvvix.onbuch&hl=fr}}%
      \hspace{8pt}\begin{minipage}[c]{0.5\linewidth}
        {\popdisplay\fontsize{11}{12.5}\selectfont\color{white}\MakeUppercase{Télécharge OnBuch}}\par\vspace{4pt}
        {\raggedright\popsemi\fontsize{8.4}{11}\selectfont\color{white}Gratuit \textbullet\ Google Play\\Tuteur IA, annales, concours, résultats\par}
      \end{minipage}
    \end{tcolorbox}
  \end{minipage}\hfill
  \begin{minipage}[t]{0.487\linewidth}
    \begin{tcolorbox}[enhanced,colback=poppurple,colframe=popink,boxrule=2.2pt,arc=10pt,
        drop shadow=popink,left=11pt,right=11pt,top=10pt,bottom=10pt,height=3.1cm,valign=center]
      \tikz[baseline=-0.7ex]\node[circle,fill=white,draw=popink,line width=1.3pt,minimum size=1.6cm,inner sep=0]{\popdisplay\fontsize{24}{24}\selectfont\color{poppurple}+};%
      \hspace{8pt}\begin{minipage}[c]{0.6\linewidth}
        {\popdisplay\fontsize{11}{12.5}\selectfont\color{white}\MakeUppercase{Passe à OnBuch+}}\par\vspace{4pt}
        {\raggedright\popsemi\fontsize{8.4}{11}\selectfont\color{white}Tous les cours signés, Tuteur IA illimité, fiches PDF — onglet OnBuch+ de l'app\par}
      \end{minipage}
    \end{tcolorbox}
  \end{minipage}
  \vspace{8pt}
  \popcontenu
  \vfill
  \signatureonbuch
  \restoregeometry
  \newpage
}

% Rangée « Ce cours contient » (page de garde)
\newcommand{\poptile}[3]{% #1 couleur #2 gros texte #3 légende
  \begin{tcolorbox}[enhanced,colback=#1,colframe=popink,boxrule=2pt,arc=9pt,shadow={2.5pt}{-2.5pt}{0pt}{popink},
      left=6pt,right=6pt,top=9pt,bottom=9pt,halign=center,valign=center,height=1.75cm,width=\linewidth,nobeforeafter]
    {\popdisplay\fontsize{12.5}{13}\selectfont\color{white}\MakeUppercase{#2}}\par\vspace{3pt}
    {\centering\popsemi\fontsize{7.8}{9.5}\selectfont\color{white}#3\par}
  \end{tcolorbox}}
\newcommand{\popcontenu}{%
  \par\noindent{\popxboldls\fontsize{7.5}{7.5}\selectfont\color{popmuted}CE COURS SIGNÉ CONTIENT}\par\vspace{7pt}
  \noindent\begin{minipage}[t]{0.235\linewidth}\poptile{popblue}{Cours}{complet, illustré, pas à pas}\end{minipage}\hfill
  \begin{minipage}[t]{0.235\linewidth}\poptile{poporange}{Méthodes}{et exercices résolus}\end{minipage}\hfill
  \begin{minipage}[t]{0.235\linewidth}\poptile{poppink}{Exercices}{type examen + corrigés}\end{minipage}\hfill
  \begin{minipage}[t]{0.235\linewidth}\poptile{popgreen}{Fiche bilan}{l'essentiel à retenir}\end{minipage}\par}

% Sommaire
\newtcolorbox{sommaire}{enhanced,colback=white,colframe=popink,boxrule=2.2pt,arc=10pt,
  drop shadow=popink,left=18pt,right=18pt,top=13pt,bottom=13pt,before skip=6pt,after skip=20pt,
  before upper={\poppill{Sommaire}{poppurple}{white}\par\vspace{9pt}\popsemi\fontsize{10.6}{14.5}\selectfont\color{popink}}}

% Signature de fin de cours — poussée tout en bas de la dernière page
\newcommand{\findecours}{%
  \par\vfill\nopagebreak
  \begin{center}\popsquiggle{poporange}\end{center}\vspace{4pt}
  \signatureonbuch
}
\AtBeginDocument{\def\llbracket{\mathopen{[\mkern-3.5mu[}}\def\rrbracket{\mathclose{]\mkern-3.5mu]}}} % intervalles d'entiers (glyphe ⟦ absent de la police)
% Glyphes absents de Fira Math : redéfinis à partir de glyphes disponibles
\AtBeginDocument{%
  \def\setminus{\mathbin{\backslash}}\let\smallsetminus\setminus
  \def\vdots{\mathord{\vbox{\baselineskip3.2pt\lineskiplimit0pt\kern4pt\hbox{.}\hbox{.}\hbox{.}}}}%
  \def\ddots{\mathinner{\mkern1mu\raise6.5pt\hbox{.}\mkern2mu\raise3.6pt\hbox{.}\mkern2mu\raise0.7pt\hbox{.}\mkern1mu}}%
  \def\triangle{\mathord{\scalebox{0.9}{$\symup{\Delta}$}}}%
  \def\nabla{\mathord{\raisebox{\depth}{\scalebox{1}[-1]{$\symup{\Delta}$}}}}%
  \def\checkmark{\popcheck{popdark}}%
  \def\star{\mathbin{\ast}}\def\bigstar{\mathord{\ast}}%
  \def\diamond{\mathbin{\scalebox{0.6}{$\Diamond$}}}%
  \def\bigcup{\mathop{\mathchoice{\raisebox{-0.3em}{\scalebox{1.7}{$\cup$}}}{\scalebox{1.25}{$\cup$}}{\cup}{\cup}}\displaylimits}%
  \def\bigcap{\mathop{\mathchoice{\raisebox{-0.3em}{\scalebox{1.7}{$\cap$}}}{\scalebox{1.25}{$\cap$}}{\cap}{\cap}}\displaylimits}%
  \def\lesseqqgtr{\mathrel{\lessgtr}}\def\bigoplus{\mathop{\oplus}\displaylimits}%
}
\providecommand{\ssi}{si et seulement si} % abréviation fréquente
\AtBeginDocument{\providecommand{\wideparen}{\overparen}} % arc de cercle

\begin{document}

\pagegarde{Dissertation rédiger une dissertation}{Français}{Tle Tech}{Français – classe de Terminale technique}{N08}{14 séances (fiche de progression harmonisée)}{Cette leçon outille l'élève de Terminale technique pour analyser un sujet de dissertation, élaborer un plan détaillé et rédiger introduction, paragraphes et conclusion, en s'appuyant sur la lecture méthodique du Vieux nègre et la médaille et sur la maîtrise des outils de liaison et des stratégies argumentatives.
\vspace{4pt}
\begin{multicols}{2}\small
\begin{itemize}
\item À la fin de cette leçon, je sais analyser un sujet de dissertation (mots-clés, type de sujet, problématique).
\item À la fin de cette leçon, je sais mobiliser les stratégies argumentatives d'un texte (arguments, exemples, procédés rhétoriques).
\item À la fin de cette leçon, je sais lire méthodiquement un extrait du Vieux nègre et la médaille et dégager sa thèse.
\item À la fin de cette leçon, je sais rédiger une introduction complète (amorce, présentation du sujet, problématique, annonce du plan).
\item À la fin de cette leçon, je sais construire un paragraphe argumentatif (idée directrice, argument, exemple, transition).
\item À la fin de cette leçon, je sais employer correctement les outils de liaison (conjonctions, adverbes, pronoms relatifs) pour assurer la cohérence du texte.
\end{itemize}\end{multicols}}

\begin{sommaire}
\begin{multicols}{2}
\begin{enumerate}
\item Lecture méthodique du Vieux nègre et la médaille (1re partie)
\item Le texte argumentatif : stratégies argumentatives
\item Les outils de liaison dans la phrase
\item Rédaction de l'introduction et de la conclusion
\item Rédaction du paragraphe argumentatif
\item Analyse du sujet, plan détaillé et production complète
\item Activité d'intégration
\item Méthodes et exercices résolus
\item Exercices
\item Corrigés des exercices
\item Fiche bilan
\end{enumerate}
\end{multicols}
\end{sommaire}

\begin{objectifs}
\begin{itemize}
\item À la fin de cette leçon, je sais analyser un sujet de dissertation (mots-clés, type de sujet, problématique).
\item À la fin de cette leçon, je sais mobiliser les stratégies argumentatives d'un texte (arguments, exemples, procédés rhétoriques).
\item À la fin de cette leçon, je sais lire méthodiquement un extrait du Vieux nègre et la médaille et dégager sa thèse.
\item À la fin de cette leçon, je sais rédiger une introduction complète (amorce, présentation du sujet, problématique, annonce du plan).
\item À la fin de cette leçon, je sais construire un paragraphe argumentatif (idée directrice, argument, exemple, transition).
\item À la fin de cette leçon, je sais employer correctement les outils de liaison (conjonctions, adverbes, pronoms relatifs) pour assurer la cohérence du texte.
\item À la fin de cette leçon, je sais rédiger une conclusion (bilan, ouverture).
\item À la fin de cette leçon, je sais élaborer un plan détaillé et rédiger une dissertation entière, puis l'améliorer après confrontation.
\item À la fin de cette leçon, je sais justifier ma démarche argumentative et appliquer ces compétences à des situations de la vie courante (débats, prises de position).
\end{itemize}
\end{objectifs}

\begin{prerequis}
\begin{itemize}
\item Notions de base sur le texte argumentatif (thèse, arguments, exemples) vues en classe de Première
\item Les types de phrases et la ponctuation (classes de 4e-3e)
\item Les connecteurs logiques simples (mais, donc, car, cependant)
\item Connaissance générale du contexte colonial camerounais (utile pour Le Vieux nègre et la médaille)
\item Techniques de lecture cursive d'une œuvre intégrale
\end{itemize}
\end{prerequis}

\newpage

\coursec{Lecture méthodique du \emph{Vieux nègre et la médaille} (1\textsuperscript{re} partie)}

\textbf{Situation d'accroche.} À l'approche du Baccalauréat technique, de nombreux candidats de Yaoundé et de Douala perdent des points non pas par manque d'idées, mais parce qu'ils ne savent pas organiser leur pensée dans une dissertation. Lors d'un débat en classe sur la question « L'école garantit-elle encore la réussite sociale au Cameroun ? », les élèves ont des opinions tranchées mais peinent à les structurer en introduction, développement et conclusion. Comment transformer une opinion personnelle en un raisonnement argumenté, cohérent et convaincant ? C'est tout l'enjeu de la dissertation, exercice roi de l'examen. À travers l'étude du \emph{Vieux nègre et la médaille} de Ferdinand Oyono, nous apprendrons à lire méthodiquement un texte argumentatif et à en tirer les outils pour construire notre propre argumentation.

\textbf{Problématique.} Comment un roman peut-il, sous le récit d'une simple remise de médaille, dénoncer l'aliénation coloniale ? Et comment la lecture méthodique de ce texte peut-elle nous donner les clés pour structurer une dissertation rigoureuse ?

\courssub{Présentation de l'œuvre et de son contexte}

\emph{Le Vieux nègre et la médaille} est un roman de \cle{Ferdinand Oyono}, écrivain camerounais né en 1929 à Ngoulemakong (région du Sud), publié en \cle{1956}, quelques années avant l'indépendance du Cameroun (1960). L'œuvre s'inscrit dans le contexte de la \cle{colonisation} : le Cameroun est alors administré par la France, et la société est divisée entre les colons européens, qui détiennent le pouvoir, et les populations africaines, soumises et méprisées.

Le roman raconte l'histoire de \cle{Meka}, un vieux paysan du village de Doum, qui a tout donné aux Blancs : sa terre pour la construction d'une mission, ses deux fils morts « pour la France » pendant la guerre, et sa foi, puisqu'il s'est converti au christianisme. Pour le récompenser, l'administration coloniale décide de lui remettre une \cle{médaille} le jour de la fête du 14 juillet. Mais cette « récompense » tourne au cauchemar : ivre et épuisé après la cérémonie, Meka se perd dans la nuit, est arrêté par les forces de l'ordre coloniales et brutalisé. Cette humiliation provoque chez lui une véritable \cle{prise de conscience} : il comprend que l'amitié des Blancs est un leurre et que la médaille n'est qu'un instrument de domination.

\begin{savaistu}[Ferdinand Oyono, écrivain et homme d'État]
Ferdinand Oyono (1929--2010) n'a pas été seulement un écrivain : après l'indépendance, il a mené une longue carrière de diplomate (ambassadeur du Cameroun dans plusieurs pays) puis de ministre. Avec \emph{Une vie de boy} (1956) et \emph{Le Vieux nègre et la médaille}, il est l'un des grands romanciers de la littérature anticoloniale africaine. Son roman figure au programme du Probatoire et du Baccalauréat technique : le connaître, c'est déjà préparer l'examen.
\end{savaistu}

\courssub{Qu'est-ce que la lecture méthodique ?}

\begin{definition}[La lecture méthodique]
La \cle{lecture méthodique} est une lecture guidée et ordonnée qui consiste à dégager d'un texte : sa \cle{situation d'énonciation} (qui parle ? à qui ? quand ? où ? de quoi ?), son \cle{thème}, sa \cle{thèse} (l'idée que l'auteur défend, souvent implicite dans un texte littéraire) et ses \cle{procédés} d'écriture (ironie, satire, point de vue, champs lexicaux). Elle se déroule en trois temps : \emph{avant} la lecture (titre, auteur, genre), \emph{pendant} la lecture (repérage des indices), \emph{après} la lecture (synthèse et interprétation).
\end{definition}

L'intuition est simple : lire méthodiquement, c'est lire comme un détective. On ne se contente pas de suivre l'histoire ; on relève des indices (mots, tons, regards des personnages) pour reconstituer ce que l'auteur veut réellement nous dire.

\courssub{Lecture méthodique n\textsuperscript{o} 1 : situation d'énonciation et personnages}

\textbf{Situation d'énonciation.} Le roman est raconté par un \cle{narrateur} extérieur à l'histoire (narration à la troisième personne), mais qui adopte souvent le point de vue des villageois africains. Le récit se situe dans le village fictif de \cle{Doum}, au Cameroun colonial, autour de la fête du 14 juillet. Le narrateur s'adresse à un lecteur qu'il invite à juger par lui-même le spectacle de la « cérémonie ».

\textbf{Les personnages principaux.}

\begin{itemize}
\item \cle{Meka} : le « vieux nègre » du titre. Vieillard naïf et généreux, il a sacrifié ses biens, ses fils et sa religion ancestrale pour les Blancs. Il incarne l'Africain aliéné par la colonisation.
\item \cle{Kelara}, sa femme : plus lucide que lui, elle se méfie des Blancs et pressent que la médaille n'apportera rien de bon.
\item \cle{Le commandant} : représentant de l'administration coloniale, il organise la cérémonie de remise de la médaille. Il incarne le pouvoir colonial, paternaliste et méprisant.
\item \cle{Le père Vandermayer} : missionnaire, directeur de la mission construite sur les terres de Meka. Il représente la colonisation religieuse, celle qui convertit et « civilise ».
\item \cle{Les villageois de Doum} : témoins ébahis de la cérémonie, ils représentent la communauté africaine.
\end{itemize}

\begin{popfigure}
\begin{center}
\begin{tikzpicture}[
  every node/.style={font=\small},
  perso/.style={draw=popblue, fill=popblueL, rounded corners=2pt, inner sep=4pt, align=center},
  blanc/.style={draw=poporange, fill=poporangeL, rounded corners=2pt, inner sep=4pt, align=center},
  rel/.style={-{Stealth[length=2mm]}, thick, popdark},
  lbl/.style={font=\scriptsize\itshape, fill=white, inner sep=1pt}]
  \node[perso, draw=poppurple, fill=poppurple!20, align=center] (meka) at (0,0){\textbf{Meka}\\ (le vieux nègre)};
  \node[perso, align=center] (kelara) at (-4.5,0){\textbf{Kelara}\\ (sa femme)};
  \node[blanc, align=center] (commandant) at (0,2.8){\textbf{Le commandant}\\ (administration)};
  \node[blanc, align=center] (pere) at (4.5,2.8){\textbf{Père Vandermayer}\\ (mission)};
  \node[perso, align=center] (village) at (4.5,0){\textbf{Villageois}\\ de Doum};
  \node[perso, align=center] (fils) at (-4.5,-2.6){\textbf{Deux fils}\\ (morts à la guerre)};
  \draw[rel] (kelara) -- node[lbl, above]{méfiance, lucidité} (meka);
  \draw[rel] (commandant) -- node[lbl, right]{remet la médaille, méprise} (meka);
  \draw[rel] (pere) -- node[lbl, right]{a pris sa terre, l'a converti} (meka);
  \draw[rel] (village) -- node[lbl, above]{admiration puis pitié} (meka);
  \draw[rel] (meka) -- node[lbl, left]{a donné « pour la France »} (fils);
  \draw[rel, poporange] (commandant) -- node[lbl, above]{alliance coloniale} (pere);
\end{tikzpicture}
\end{center}
\legende{Arbre des relations entre les personnages du \emph{Vieux nègre et la médaille}. En violet : le héros ; en bleu : les Africains ; en orange : les représentants de la colonisation.}
\end{popfigure}

\courssub{Lecture méthodique n\textsuperscript{o} 2 : dégager la thèse implicite}

Un roman ne dit jamais « je vais vous prouver que... ». Sa thèse est \cle{implicite} : elle se dégage de l'ensemble du récit, du sort réservé aux personnages et du ton du narrateur.

\begin{methode}[Dégager la thèse d'un texte littéraire à visée argumentative]
\begin{enumerate}
\item \textbf{Identifier le thème} : de quoi parle le texte ? (Ici : la remise d'une médaille à un vieil Africain sous la colonisation.)
\item \textbf{Observer le sort des personnages} : qui est récompensé ? qui est puni ? qui a raison à la fin ? (Meka, « honoré », finit humilié et battu ; Kelara, la méfiante, avait raison.)
\item \textbf{Repérer le ton du narrateur} : admire-t-il, plaint-il, ridiculise-t-il ? (Le narrateur ridiculise la cérémonie et plaint Meka.)
\item \textbf{Formuler la thèse en une phrase complète} : « L'auteur montre que... » ou « L'auteur dénonce que... »
\end{enumerate}
\end{methode}

Appliquons la méthode. Le thème est la médaille offerte à Meka. Or cette médaille, symbole de reconnaissance, n'apporte à Meka ni terre, ni argent, ni respect : elle lui attire au contraire l'arrestation et les coups. Le narrateur souligne sans cesse le contraste entre la solennité de la cérémonie et le néant de la récompense. On peut donc formuler la \cle{thèse implicite} ainsi :

\begin{aretenir}[La thèse du roman]
Dans \emph{Le Vieux nègre et la médaille}, Ferdinand Oyono dénonce l'\cle{aliénation coloniale} : la colonisation endort les Africains avec des récompenses illusoires (médailles, religion, « amitié » des Blancs) pour mieux les déposséder de leurs terres, de leurs enfants et de leur dignité. La médaille est le symbole de cette tromperie.
\end{aretenir}

\courssub{Repérage des arguments et des exemples dans le texte}

Pour convaincre le lecteur, Oyono aligne des \cle{arguments} (raisons qui soutiennent la thèse) illustrés par des \cle{exemples} (faits précis du récit). C'est exactement la structure que vous devrez reproduire dans vos paragraphes de dissertation (schéma \cle{Idée directrice -- Argument -- Exemple}).

\begin{center}
\begin{tabularx}{\linewidth}{|l|X|X|}
\hline
\rowcolor{popblueL} \textbf{N\textsuperscript{o}} & \textbf{Argument (ce que l'auteur veut prouver)} & \textbf{Exemple (ce qui le prouve dans le récit)} \\
\hline
1 & La colonisation dépossède les Africains de leurs biens & Meka a cédé sa terre pour la construction de la mission du père Vandermayer \\
\hline
2 & Elle sacrifie les vies africaines à des causes étrangères & Les deux fils de Meka sont morts à la guerre « pour la France » \\
\hline
3 & Elle détruit les croyances et l'identité africaines & Meka a abandonné la religion de ses ancêtres pour le christianisme \\
\hline
4 & Les récompenses coloniales sont des leurres & La médaille ne donne à Meka aucun droit : il est arrêté et battu le soir même de la fête \\
\hline
5 & L'humiliation ouvre les yeux : la prise de conscience & Après sa nuit au poste, Meka comprend que « les Blancs » ne sont pas ses amis \\
\hline
\end{tabularx}
\end{center}

\begin{methode}[Rédiger un paragraphe argumentatif (structure I.A.E.)]
Chaque paragraphe de développement suit le même squelette :
\begin{enumerate}
\item \textbf{Idée directrice (phrase thèse)} : une phrase qui annonce le point précis défendu dans le paragraphe (ex. : « La colonisation spolie les Africains de leurs biens fonciers. »).
\item \textbf{Argument(s) / Explication} : on développe le raisonnement, on explique le mécanisme (ex. : « Sous couvert de développement, l'administration s'approprie les terres ancestrales sans contrepartie réelle. »).
\item \textbf{Exemple(s) / Illustration} : on cite un fait précis, une citation, une référence du texte étudié ou de sa culture générale (ex. : « Dans le roman, Meka donne sa terre pour la mission du père Vandermayer et ne reçoit qu'une médaille en retour. »).
\item \textbf{Mini-conclusion / Transition} : on résume l'apport du paragraphe et on ouvre sur le suivant.
\end{enumerate}
\end{methode}

\begin{popfigure}
\begin{center}
\begin{tikzpicture}[font=\small]
  \draw[-{Stealth[length=3mm]}, very thick, popdark] (0,0) -- (13,0);
  \foreach \x/\date/\evt in {0.6/{Sacrifices de Meka}/{terre, fils, religion donnés aux Blancs},
                             4.2/{Annonce de la médaille}/{joie du village, fierté de Meka},
                             7.8/{14 juillet : remise de la médaille}/{cérémonie, fête, ivresse},
                             11.2/{Nuit : arrestation}/{coups, humiliation, prise de conscience}}{
    \draw[thick, popblue] (\x,0.15) -- (\x,-0.15);
    \node[above, align=center, text width=2.9cm, font=\scriptsize] at (\x,0.2) {\textbf{\date}\\ \evt};
  }
\end{tikzpicture}
\end{center}
\legende{Frise chronologique simplifiée des événements du roman : de l'attribution de la médaille à la prise de conscience de Meka.}
\end{popfigure}

\courssub{Analyse des procédés : ironie, satire, point de vue narratif}

\begin{definition}[Ironie et satire]
L'\cle{ironie} est un procédé qui consiste à dire le contraire de ce que l'on pense, ou à souligner l'écart entre les apparences et la réalité, pour faire réfléchir ou rire le lecteur. La \cle{satire} est la dénonciation des défauts d'une société par la moquerie et le ridicule.
\end{definition}

Le \cle{point de vue narratif} renforce ce procédé : le narrateur adopte souvent le regard naïf de Meka et des villageois, ébahis devant les discours officiels. Ce décalage entre ce que croit Meka (« les Blancs sont mes amis ») et ce que le lecteur voit (le mépris des Blancs) crée une ironie permanente.

\begin{exemplebox}[L'ironie dans la scène de la remise de la médaille]
Au moment même où Meka est « honoré », il est en réalité humilié : on le fait défiler comme une bête curieuse devant la foule ; on le gave de vin jusqu'à l'ivresse ; les Blancs rient de lui et non avec lui. La médaille, posée sur sa poitrine, ne change rien à sa condition : quelques heures plus tard, les mêmes autorités qui l'ont décoré le font battre au poste de police. Dire que Meka est « honoré » alors qu'il est ridiculisé et bientôt brutalisé, c'est le sommet de l'ironie : l'écart entre l'apparence (la fête) et la réalité (la domination) dénonce l'hypocrisie coloniale.
\end{exemplebox}

\courssub{Étude lexicale : médaille, honneur et humiliation}

L'étude du \cle{lexique} (sémantique et lexicologie) permet de vérifier la thèse : les mots du texte s'organisent en \cle{champs lexicaux} qui s'opposent.

\begin{center}
\begin{tabularx}{\linewidth}{|l|X|X|}
\hline
\rowcolor{popblueL} \textbf{Champ lexical} & \textbf{Mots et expressions repérés} & \textbf{Effet produit} \\
\hline
La médaille & médaille, ruban, décoration, croix, récompense, distinction & Apparence de gloire ; objet matériel sans valeur réelle \\
\hline
L'honneur & honorer, fête, cérémonie, discours, solennité, fierté, ami de la France & Langue officielle des Blancs ; illusion entretenue \\
\hline
L'humiliation & vieux nègre, rire, ivresse, gifle, coup de pied, cellule, honte, nuit & Réalité vécue par Meka ; vérité de la colonisation \\
\hline
\end{tabularx}
\end{center}

Remarquons le titre lui-même : « le vieux \emph{nègre} » est un terme méprisant, celui du regard colonial, tandis que « la médaille » est le signe de l'honneur. Le titre condense donc à lui seul l'ironie du roman : l'honneur et le mépris accolés.

\begin{attention}[Erreur fréquente à l'examen]
Ne confondez jamais \textbf{résumer} et \textbf{analyser}. Raconter la scène (« Meka reçoit une médaille, puis il est arrêté ») ne rapporte presque aucun point. Il faut dégager la \cle{visée argumentative} de l'extrait : « En montrant Meka humilié au moment même où il est honoré, Oyono dénonce l'hypocrisie de la colonisation ». À chaque fait du récit, posez-vous la question : \emph{que veut prouver l'auteur ici ?}
\end{attention}

\begin{exoresolu}[Dégager la visée argumentative d'un passage]
Énoncé. Après la cérémonie, Meka, ivre, s'égare dans la nuit ; il est arrêté par les gardes, insulté et enfermé. Dégagez en trois phrases la visée argumentative de ce passage (procédé, sens, thèse).
\tcblower
Solution détaillée.
\begin{enumerate}
\item \textbf{Procédé} : Oyono utilise l'ironie et le contraste brutal entre la fête du 14 juillet et la nuit de l'arrestation.
\item \textbf{Sens} : le même État colonial qui a décoré Meka le matin le fait frapper le soir ; la médaille ne protège de rien.
\item \textbf{Thèse} : l'auteur dénonce ainsi l'aliénation coloniale : les récompenses offertes aux Africains sont des leurres qui masquent la violence réelle de la domination.
\end{enumerate}
\end{exoresolu}

\begin{exercice}[À vous de jouer]
\begin{enumerate}
\item Relevez dans le roman trois mots ou expressions appartenant au champ lexical de l'humiliation et expliquez leur effet.
\item En vous aidant de la méthode étudiée (I.A.E.), rédigez un court paragraphe (5 à 8 lignes) qui dégage la thèse implicite de la scène de la remise de la médaille, en citant au moins un procédé (ironie, satire, point de vue).
\item Pourquoi peut-on dire que Kelara, la femme de Meka, incarne la lucidité africaine face à l'aliénation ?
\end{enumerate}
\end{exercice}

\begin{corrige}[Exercice « À vous de jouer »]
\begin{enumerate}
\item Exemples attendus : « vieux nègre » (mépris racial), « ivresse » (ridicule), « coup de pied », « cellule » (violence physique). Ces mots montrent que derrière la fête se cache la brutalité réelle de la colonisation.
\item Le paragraphe doit suivre le schéma I.A.E. : \textbf{Idée} (la cérémonie est une mascarade), \textbf{Argument} (procédé : ironie, Meka est humilié au moment où il est honoré), \textbf{Exemple} (le vin, les rires des Blancs, l'arrestation nocturne), \textbf{Mini-conclusion} (Oyono dénonce l'hypocrisie coloniale).
\item Kelara se méfie des Blancs et annonce que la médaille n'apportera rien ; les événements lui donnent raison. Elle représente la conscience claire que la colonisation est une tromperie, contrairement à Meka, longtemps aveuglé.
\end{enumerate}
\end{corrige}

\begin{aretenir}[Bilan de la section]
\begin{itemize}
\item \emph{Le Vieux nègre et la médaille} (Ferdinand Oyono, 1956) dénonce l'aliénation coloniale à travers le destin de Meka.
\item La lecture méthodique dégage la situation d'énonciation, le thème, la thèse (implicite) et les procédés.
\item La thèse se formule en une phrase : « l'auteur montre/dénonce que... »
\item Les arguments du texte s'appuient sur des exemples précis du récit ; les procédés majeurs sont l'ironie, la satire et le point de vue naïf.
\item Les champs lexicaux (médaille / honneur / humiliation) confirment la thèse.
\item \textbf{Lien avec la dissertation} : la structure Argument / Exemple repérée dans le texte est le modèle exact du paragraphe argumentatif (Idée directrice -- Argument -- Exemple).
\end{itemize}
\end{aretenir}

\coursec{Le texte argumentatif : stratégies argumentatives}

Dans la section précédente, nous avons lu méthodiquement la première partie du \emph{Vieux nègre et la médaille} de Ferdinand Oyono. Nous y avons découvert Meka, ce vieillard que l'administration coloniale veut décorer, et nous avons remarqué que le roman entier repose sur un débat : la médaille est-elle une récompense sincère ou un instrument de domination ? Autrement dit, Oyono construit un \cle{texte argumentatif}. Pour bien lire ce roman — et pour réussir la dissertation au Probatoire et au Baccalauréat technique — il faut maintenant comprendre \emph{comment} un auteur défend une idée : c'est l'objet de cette section.

\courssub{Qu'est-ce qu'une stratégie argumentative ?}

Imaginons une scène de la vie courante : un élève veut convaincre son père de lui acheter un téléphone. Il ne dit pas seulement « donne-moi un téléphone ». Il avance des raisons (« tous mes camarades en ont un », « je pourrai réviser avec les applications »), il donne des exemples (« Marlyse a eu le sien et ses notes ont monté »), il joue parfois sur l'émotion (« tu veux que je sois le seul exclu du groupe ? »). Sans le savoir, cet élève met en œuvre une \cle{stratégie argumentative}.

\begin{definition}[Stratégie argumentative]
La \cle{stratégie argumentative} est l'ensemble des moyens — \cle{arguments}, \cle{exemples} et \cle{procédés rhétoriques} — mis en œuvre par un locuteur ou un auteur pour faire adhérer le lecteur ou l'auditeur à une \cle{thèse} (une opinion, un point de vue). Elle répond à la question : \emph{comment l'auteur s'y prend-il pour me faire penser comme lui ?}
\end{definition}

\begin{aretenir}[Structure du texte argumentatif — rappel]
Tout texte argumentatif s'organise autour de quatre éléments :
\begin{enumerate}
\item la \cle{thèse} : l'idée défendue (ex. « la médaille est un piège ») ;
\item les \cle{arguments} : les raisons générales qui soutiennent la thèse ;
\item les \cle{exemples} : les cas précis, concrets, qui illustrent les arguments ;
\item la \cle{conclusion} : le point d'arrivée qui réaffirme la thèse.
\end{enumerate}
\end{aretenir}

\begin{popfigure}
\begin{center}
\begin{tikzpicture}[scale=1]
% Pyramide argumentative
\fill[popblue!70] (0,0) -- (1.2,1.6) -- (-1.2,1.6) -- cycle;
\fill[poporange!60] (-1.2,1.6) -- (1.2,1.6) -- (2.4,3.2) -- (-2.4,3.2) -- cycle;
\fill[popgreen!55] (-2.4,3.2) -- (2.4,3.2) -- (3.6,4.8) -- (-3.6,4.8) -- cycle;
\draw[popdark,thick] (0,0) -- (3.6,4.8) -- (-3.6,4.8) -- cycle;
\draw[popdark] (-1.2,1.6) -- (1.2,1.6);
\draw[popdark] (-2.4,3.2) -- (2.4,3.2);
\node[white,font=\bfseries] at (0,0.9) {THÈSE};
\node[popdark,font=\bfseries] at (0,2.4) {ARGUMENTS};
\node[popdark,font=\bfseries] at (0,4.0) {EXEMPLES};
% Annotations
\node[right,align=left,font=\small] at (4.2,0.9) {l'idée défendue\\ (une seule, claire)};
\node[right,align=left,font=\small] at (4.2,2.4) {les raisons générales\\ (2 ou 3 par partie)};
\node[right,align=left,font=\small] at (4.2,4.0) {les cas concrets qui prouvent\\ (faits, textes, statistiques)};
\draw[->,popmuted] (4.1,0.9) -- (1.1,0.9);
\draw[->,popmuted] (4.1,2.4) -- (2.1,2.4);
\draw[->,popmuted] (4.1,4.0) -- (3.1,4.0);
\end{tikzpicture}
\end{center}
\legende{La pyramide argumentative : la thèse domine, elle est soutenue par les arguments, eux-mêmes appuyés sur les exemples. Sans exemples, les arguments restent abstraits ; sans arguments, la thèse n'est qu'une affirmation gratuite.}
\end{popfigure}

\courssub{Les types d'arguments}

Tous les arguments ne se valent pas. Pour analyser un texte (ou construire ta propre dissertation), il faut savoir les reconnaître et les classer.

\textbf{1. L'argument logique (ou rationnel).} Il s'adresse à la \cle{raison} du lecteur : faits vérifiables, statistiques, démonstrations, constats d'expérience. Exemple : « Le taux de réussite au Baccalauréat technique progresse lorsque les élèves s'entraînent régulièrement à la dissertation » — affirmation vérifiable par les chiffres officiels.

\textbf{2. L'argument d'autorité.} Il s'appuie sur la parole d'une personne reconnue : expert, chef traditionnel, écrivain célèbre, texte sacré ou loi. Sa force vient du \emph{prestige} de celui qu'on cite, non d'une démonstration.

\begin{exemplebox}[L'argument d'autorité dans le contexte camerounais]
Dans un débat sur la coutume, un élève écrit : « Comme l'a déclaré le sultan des Bamouns, la tradition n'empêche pas le progrès, elle l'éclaire. » Citer un chef traditionnel respecté (ou un recteur, un médecin-chef, un prix Nobel) donne du poids à la thèse. Attention : l'argument d'autorité ne dispense pas de prouver — un chef peut se tromper.
\end{exemplebox}

\textbf{3. L'argument affectif (ou pathétique).} Il touche les \cle{émotions} : pitié, peur, fierté, colère, honte. Exemple : « Imaginez ce vieillard, après cinquante ans de loyaux services, humilié devant tout le village le jour de sa médaille. » Cet argument est puissant mais peut devenir manipulateur s'il remplace totalement la raison.

\begin{methode}[Identifier et classer les arguments d'un texte]
\begin{enumerate}
\item \textbf{Repère la thèse} : demande-toi « que veut me faire croire l'auteur ? ».
\item \textbf{Surligne les raisons avancées} : cherche les connecteurs (\emph{car, parce que, en effet, d'abord, ensuite, enfin}).
\item \textbf{Classe chaque argument} : est-il logique (fait, chiffre), d'autorité (citation d'une personne reconnue) ou affectif (appel aux émotions) ?
\item \textbf{Vérifie l'exemple associé} : chaque argument est-il illustré par un cas précis ?
\item \textbf{Juge l'efficacité} : l'argument est-il solide, pertinent, ou fallacieux (sophisme) ?
\end{enumerate}
\end{methode}

\begin{attention}[Erreur fréquente : confondre argument et exemple]
L'\cle{argument} est une \textbf{raison générale} (« les récompenses coloniales divisent les communautés ») ; l'\cle{exemple} est un \textbf{cas précis} qui l'illustre (« Meka est rejeté par les siens après avoir reçu la médaille »). Au Bac, écrire « l'argument est Meka » est une faute d'analyse : Meka est un exemple, pas un argument. Formule correcte : « Oyono illustre par l'exemple de Meka l'argument selon lequel la médaille isole l'individu de sa communauté. »
\end{attention}

\courssub{Convaincre, persuader, délibérer : les trois visées de l'argumentation}

\begin{definition}[Convaincre et persuader]
\cle{Convaincre} consiste à obtenir l'adhésion par la \textbf{raison} : preuves, logique, démonstration (registre du discours scientifique, du débat d'idées). \cle{Persuader} consiste à obtenir l'adhésion par les \textbf{émotions} et les sentiments : images, rythme, ironie, pathos (registre du discours publicitaire, politique, littéraire). \cle{Délibérer} consiste à peser le pour et le contre avant de trancher : c'est la visée de la dissertation dialectique (thèse / antithèse / synthèse).
\end{definition}

\begin{popfigure}
\begin{center}
\begin{tikzpicture}
\node[draw=popblue,fill=popblueL,rounded corners,minimum width=5.6cm,minimum height=1cm,font=\bfseries] (conv) at (-3.2,3.4) {CONVAINCRE};
\node[draw=poppink,fill=poppinkL,rounded corners,minimum width=5.6cm,minimum height=1cm,font=\bfseries] (pers) at (3.2,3.4) {PERSUADER};
\node[align=left,font=\small,anchor=north west] at (-6.2,2.6) {\textbf{Arme :} la raison\\ \textbf{Moyens :} faits, chiffres,\\ démonstrations, logique\\ \textbf{Exemple :} « Les statistiques\\ montrent que la scolarisation\\ des filles réduit la pauvreté. »};
\node[align=left,font=\small,anchor=north west] at (0.6,2.6) {\textbf{Arme :} l'émotion\\ \textbf{Moyens :} images, ironie,\\ gradation, pathos\\ \textbf{Exemple :} « Fermerez-vous\\ les yeux pendant que ces\\ enfants grandissent sans école ? »};
\node[draw=popgreen,fill=popgreenL,rounded corners,align=center,font=\small\bfseries] (delib) at (0,-2.2) {DÉLIBÉRER : examiner le pour et le contre,\\ puis trancher en connaissance de cause (dissertation dialectique)};
\draw[->,popmuted,thick] (conv.south) -- (delib.north west);
\draw[->,popmuted,thick] (pers.south) -- (delib.north east);
\end{tikzpicture}
\end{center}
\legende{Tableau comparatif schématisé : convaincre s'adresse à la raison, persuader aux émotions ; la délibération combine les deux en examinant les deux faces d'un problème avant de conclure.}
\end{popfigure}

Dans la pratique, les grands orateurs et les grands écrivains \textbf{combinent} les deux registres : Oyono, dans \emph{Le Vieux nègre et la médaille}, persuade par l'ironie et le comique de situation, mais il convainc aussi en accumulant les faits qui démontrent l'absurdité du système colonial.

\courssub{Les procédés rhétoriques}

Les procédés rhétoriques sont les « outils de style » de la persuasion. En voici les six principaux à connaître pour le Bac technique.

\begin{definition}[Principaux procédés rhétoriques]
\begin{itemize}
\item La \cle{question rhétorique} : question qui n'attend pas de réponse, car celle-ci est évidente. « N'avons-nous pas assez souffert ? » Elle fait adhérer le lecteur en l'obligeant à répondre mentalement « si ».
\item L'\cle{ironie} : dire le contraire de ce qu'on pense, pour ridiculiser. « Quelle belle récompense que cette médaille, pour cinquante ans de misère ! » C'est l'arme favorite d'Oyono.
\item L'\cle{antithèse} : opposition de deux termes ou idées contraires. « Ils l'appellent honneur ; nous appelons cela chaîne. » Elle structure la pensée en mettant en lumière un contraste.
\item La \cle{gradation} : énumération de termes de force croissante (ou décroissante). « Il a donné sa jeunesse, ses forces, sa vie. » Elle crée un effet d'amplification.
\item La \cle{répétition} (anaphore quand elle est en tête de phrase) : « Jamais il n'a été payé, jamais il n'a été remercié, jamais il n'a été respecté. » Elle martèle l'idée.
\item La \cle{concession} : admettre un point de vue adverse pour mieux le réfuter. « Certes, la médaille flatte l'orgueil de Meka ; mais cet orgueil le coupe des siens. » Elle donne une impression d'honnêteté et d'équilibre.
\end{itemize}
\end{definition}

\begin{exoresolu}[Repérer les stratégies dans un extrait]
\textbf{Énoncé.} Dans \emph{Le Vieux nègre et la médaille}, Oyono écrit à peu près : « Le commandant en second déclara que Meka était un homme bon, un homme fidèle, un homme dévoué, et qu'il méritait la médaille que la France, dans sa grande générosité, accordait à ses meilleurs serviteurs. » Relève deux procédés rhétoriques et explique leur effet.
\tcblower
\textbf{Solution détaillée.}
\begin{enumerate}
\item \textbf{Gradation} : « un homme bon, un homme fidèle, un homme dévoué » — les qualificatifs s'enchaînent en s'amplifiant pour dresser un portrait élogieux. Effet : le discours officiel gonfle les louanges pour justifier la médaille.
\item \textbf{Ironie} : « la France, dans sa grande générosité » — Oyono fait semblant d'admirer cette générosité, alors que le roman montre que la médaille est un instrument de propagande. Effet : le lecteur comprend le contraire de ce qui est dit ; la colonisation est ridiculisée.
\end{enumerate}
\textbf{Conclusion} : Oyono persuade par l'ironie et la gradation ; sous l'éloge apparent, il dénonce. C'est une stratégie de dénonciation indirecte, typique du roman satirique.
\end{exoresolu}

\textbf{Application à la presse camerounaise.} Prenons un éditorial de presse sur l'état des routes : « Combien de vies faut-il encore perdre sur l'axe Yaoundé--Douala ? Des promesses, toujours des promesses, rien que des promesses. Certes, le gouvernement a lancé des travaux ; mais qui ne voit que les nids-de-poule renaissent après chaque pluie ? » On y relève : une \emph{question rhétorique} (indignation), une \emph{répétition} (« promesses » martelé), une \emph{concession} (« certes... mais »). L'éditorialiste persuade par l'émotion tout en convainquant par le constat factuel.

\courssub{Complément Bac : argument direct et argument implicite}

\begin{definition}[Argument direct et argument implicite]
Un \cle{argument direct} est explicitement formulé dans le texte (« la médaille divise le village, car les anciens amis de Meka se détournent de lui »). Un \cle{argument implicite} (ou sous-entendu) n'est pas écrit : le lecteur doit le reconstituer à partir d'indices (ironie, situation, silences du texte). Quand Oyono décrit la joie naïve de Meka sans commentaire, le sous-entendu est : « cet homme ne voit pas le piège » — c'est au lecteur de le comprendre.
\end{definition}

\begin{attention}[Au Bac]
Quand tu analyses un texte, ne te contente pas des arguments explicites : les correcteurs valorisent le candidat qui sait dégager un \cle{sous-entendu} et le justifier par un indice précis du texte (un mot ironique, une répétition, un contraste).
\end{attention}

\begin{savaistu}[Le point Bac]
Au Baccalauréat technique, la question « Quelles stratégies l'auteur utilise-t-il pour convaincre (ou persuader) ? » rapporte souvent \textbf{2 à 4 points} dans les questions de lecture méthodique. Pour les obtenir, il suffit de nommer le procédé, de citer le passage et d'expliquer son effet : trois réflexes qui font la différence.
\end{savaistu}

\begin{exercice}[Entraînement]
Dans un discours, un candidat aux élections locales déclare : « Mes chers frères, faut-il que nos enfants continuent à étudier à la lampe tempête ? Le ministre de l'Énergie lui-même reconnaît que 40 \% de nos villages n'ont pas l'électricité. Votez pour moi, et l'école, le centre de santé et le marché seront électrifiés ! »
\begin{enumerate}
\item Identifie la thèse du candidat.
\item Relève un argument d'autorité, un argument logique et un argument affectif.
\item Repère deux procédés rhétoriques et explique leur effet.
\item Ce discours cherche-t-il plutôt à convaincre ou à persuader ? Justifie.
\end{enumerate}
\end{exercice}

\begin{corrige}[Exercice]
\begin{enumerate}
\item Thèse : « il faut voter pour moi car je résoudrai le problème de l'électrification rurale ».
\item Argument d'autorité : « le ministre de l'Énergie lui-même reconnaît... » ; argument logique : la statistique « 40 \% de nos villages n'ont pas l'électricité » ; argument affectif : « nos enfants... à la lampe tempête » (pitié, tendresse envers les enfants).
\item Question rhétorique : « faut-il que nos enfants continuent... ? » — elle force l'auditeur à répondre « non ». Gradation : « l'école, le centre de santé et le marché » — énumération croissante qui promet de plus en plus.
\item Le discours cherche surtout à \textbf{persuader} (appel aux émotions, promesses), mais il s'appuie aussi sur un argument logique pour convaincre : les deux registres sont combinés, comme dans tout bon discours politique.
\end{enumerate}
\end{corrige}

\begin{aretenir}[L'essentiel de la section]
\begin{itemize}
\item La stratégie argumentative = thèse + arguments + exemples + procédés rhétoriques.
\item Trois types d'arguments : logiques, d'autorité, affectifs.
\item Convaincre vise la raison ; persuader vise l'émotion ; délibérer pèse le pour et le contre.
\item Six procédés à maîtriser : question rhétorique, ironie, antithèse, gradation, répétition, concession.
\item Distingue toujours l'argument (raison générale) de l'exemple (cas précis), et l'argument direct du sous-entendu.
\end{itemize}
\end{aretenir}

Ces outils d'analyse nous serviront immédiatement dans la suite de la leçon : avant de rédiger notre propre dissertation, nous devrons maîtriser les \emph{outils de liaison dans la phrase}, qui permettent d'enchaîner arguments et exemples avec clarté.

\coursec{Les outils de liaison dans la phrase}

Dans la section précédente, nous avons étudié les \cle{stratégies argumentatives} : nous savons désormais choisir nos arguments, les hiérarchiser et utiliser les procédés rhétoriques (question rhétorique, concession, ironie) pour convaincre. Mais des arguments, même excellents, ne suffisent pas : il faut encore les \emph{relier} entre eux. Une dissertation est un tissu de phrases ; les outils de liaison sont les fils qui cousent ce tissu. Sans eux, le devoir ressemble à une liste de courses ; avec eux, il devient un raisonnement fluide, ce que le correcteur du Baccalauréat technique attend et récompense.

\courssub{Le connecteur logique : définition et intuition}

Imaginez deux phrases : \og Le téléphone portable distrait les élèves. Le téléphone portable facilite la recherche documentaire. \fg{} Juxtaposées, elles laissent le lecteur perplexe : l'auteur est-il pour ou contre ? Ajoutez un seul mot --- \emph{mais}, \emph{car}, \emph{donc} --- et tout s'éclaire : \og Le téléphone portable distrait les élèves, \textbf{mais} il facilite la recherche documentaire. \fg{} Ce petit mot a établi une \emph{relation logique} (ici, l'opposition) entre les deux énoncés.

\begin{definition}[Le connecteur logique]
Un \cle{connecteur logique} (ou outil de liaison) est un mot ou un groupe de mots qui établit une \textbf{relation logique} entre deux énoncés : addition, opposition, cause, conséquence, temps, explication, comparaison. Il peut être une conjonction de coordination, une conjonction de subordination (ou locution conjonctive), un adverbe de liaison, une locution adverbiale ou un pronom relatif.
\end{definition}

\begin{attention}[Pourquoi les correcteurs y sont sensibles]
Deux erreurs fréquentes pénalisent les copies : \textbf{enchaîner les phrases sans aucun connecteur} (style télégraphique, raisonnement invisible) et \textbf{employer un connecteur familier ou impropre} dans un texte formel. Ainsi, \og par contre \fg{} est à éviter dans une dissertation : on lui préfère \og en revanche \fg{} ou \og cependant \fg{}. De même, \og comme \fg{} au sens de \og parce que \fg{} en début de phrase est toléré à l'oral mais déconseillé à l'écrit formel.
\end{attention}

\courssub{Les conjonctions de coordination}

\begin{definition}[Conjonction de coordination]
Une \cle{conjonction de coordination} relie deux mots, deux groupes de mots ou deux propositions de \textbf{même fonction} et de \textbf{même rang}. Il en existe sept : \textbf{mais, ou, et, donc, or, ni, car} (moyen mnémotechnique : \og Mais où est donc Ornicar ? \fg{}).
\end{definition}

Chacune a une valeur logique précise. Attention : \textbf{car} ne relie que des propositions (pas des mots ou groupes de mots).

\begin{center}
\begin{tabularx}{\linewidth}{|l|l|X|}
\hline
\rowcolor{popblueL} \textbf{Conjonction} & \textbf{Valeur logique} & \textbf{Exemple} \\
\hline
\cle{et} & addition & Meka est vieux \textbf{et} il est pauvre. \\
\hline
\cle{mais} & opposition & Meka est décoré, \textbf{mais} il reste méprisé. \\
\hline
\cle{ou} & alternative & Obéira-t-il \textbf{ou} résistera-t-il ? \\
\hline
\cle{donc} & conséquence & Il a compris la supercherie, \textbf{donc} il s'enfuit. \\
\hline
\cle{or} & passage à un fait nouveau, renforcement & Il croyait être honoré ; \textbf{or}, il était ridiculisé. \\
\hline
\cle{ni} & addition négative & Il n'a ni force \textbf{ni} soutien. \\
\hline
\cle{car} & cause, justification (propositions seulement) & Il s'enfuit, \textbf{car} il a tout compris. \\
\hline
\end{tabularx}
\end{center}

\begin{attention}[Pièges de la coordination]
\begin{itemize}
\item \textbf{car} n'ouvre jamais une phrase : on écrit \og Il part, car il a compris \fg{}, jamais \og Car il a compris, il part \fg{}.
\item \textbf{donc} exprime la conséquence logique ; ne le confondez pas avec \textbf{car} (la cause). La cause précède, la conséquence suit.
\item \textbf{or} ne signifie pas \og mais \fg{} : il introduit un fait nouveau, souvent inattendu, qui fait avancer le raisonnement.
\item \textbf{ni} s'emploie toujours après une négation (ne... ni... ni...) ou redouble une négation.
\end{itemize}
\end{attention}

\courssub{Les conjonctions de subordination (et locutions conjonctives)}

\begin{definition}[Conjonction de subordination]
Une \cle{conjonction de subordination} (ou locution conjonctive) introduit une \textbf{proposition subordonnée}, c'est-à-dire une proposition qui dépend d'une proposition principale et ne peut pas se suffire à elle-même. Contrairement à la coordination qui met deux énoncés sur le même plan, la subordination crée une \textbf{hiérarchie}.
\end{definition}

\begin{center}
\begin{tabularx}{\linewidth}{|l|l|X|}
\hline
\rowcolor{popblueL} \textbf{Connecteur} & \textbf{Valeur logique} & \textbf{Exemple} \\
\hline
\cle{que} & complétive & Le commandant pense \textbf{que} Meka mérite une médaille. \\
\hline
\cle{parce que} & cause & Meka est décoré \textbf{parce qu'}il est le plus vieux. \\
\hline
\cle{puisque} & cause connue, évidente & \textbf{Puisque} la médaille est une farce, il la rejette. \\
\hline
\cle{bien que} & concession (+ subjonctif) & \textbf{Bien qu'}il soit décoré, Meka reste humilié. \\
\hline
\cle{si} & condition, hypothèse & \textbf{Si} l'école interdit le portable, les élèves protesteront. \\
\hline
\cle{lorsque} & temps & \textbf{Lorsque} la cérémonie commence, Meka est fier. \\
\hline
\cle{afin que} & but (+ subjonctif) & L'auteur écrit \textbf{afin que} l'on comprenne l'injustice. \\
\hline
\end{tabularx}
\end{center}

\begin{attention}[Mode après \og bien que \fg{} et \og afin que \fg{}]
\og Bien que \fg{} et \og afin que \fg{} exigent le \textbf{subjonctif} : \og Bien qu'il \emph{soit} riche... \fg{} ; \og Il lutte \emph{afin qu'}il \emph{obtienne} justice \fg{}. C'est une faute classique qui coûte des points au Bac. À l'inverse, \og après que \fg{} se construit avec l'indicatif (même si l'usage hésite, la norme scolaire impose l'indicatif).
\end{attention}

\begin{popfigure}
\begin{center}
\begin{tikzpicture}[
  prop/.style={rectangle, rounded corners, draw=popblue, fill=popblueL, align=center, minimum height=1cm, inner sep=6pt},
  conj/.style={rectangle, draw=poporange, fill=poporangeL, align=center, inner sep=4pt, font=\bfseries},
  fleche/.style={-{Stealth[length=3mm]}, thick, popdark}
]
\node[prop, align=center] (pp) at (0,0){Proposition principale\\ \emph{Meka rejette la médaille}};
\node[conj] (c) at (5.6,0) {parce que};
\node[prop, draw=poppurple, fill=poppurpleL, align=center] (ps) at (10.6,0){Proposition subordonnée\\ \emph{il a compris la supercherie}};
\draw[fleche] (pp) -- (c);
\draw[fleche] (c) -- (ps);
\draw[fleche, poppurple, dashed, bend right=35] (ps.south) to node[below, font=\itshape, text=poppurple]{la subordonnée dépend de la principale} (pp.south);
\end{tikzpicture}
\legende{Schéma de la subordination : le connecteur (orange) relie la proposition principale à la proposition subordonnée ; la flèche violette en pointillés montre la dépendance de la subordonnée, qui ne peut exister seule.}
\end{center}
\end{popfigure}

\courssub{Les adverbes de liaison et les locutions adverbiales}

Dans une dissertation, les adverbes et locutions de liaison jouent un rôle capital : ils articulent les \emph{paragraphes} et les \emph{étapes du raisonnement}, là où les conjonctions relient surtout des propositions à l'intérieur de la phrase.

\begin{definition}[Adverbe et locution adverbiale de liaison]
Un \cle{adverbe de liaison} (\emph{d'abord, ensuite, enfin, cependant, néanmoins, d'ailleurs}) ou une \cle{locution adverbiale de liaison} (\emph{en effet, en outre, par conséquent, en revanche, c'est-à-dire}) établit une relation logique entre deux phrases ou deux paragraphes. Il se place généralement en tête de phrase, suivi d'une virgule.
\end{definition}

\begin{center}
\begin{tabularx}{\linewidth}{|l|X|}
\hline
\rowcolor{popblueL} \textbf{Relation logique} & \textbf{Connecteurs utiles (registre formel)} \\
\hline
Ordre, plan & d'abord, ensuite, puis, enfin, en premier lieu, pour conclure \\
\hline
Addition & d'ailleurs, en outre, de plus, également \\
\hline
Opposition & cependant, néanmoins, en revanche, pourtant, toutefois \\
\hline
Cause, explication & en effet, car, effectivement \\
\hline
Conséquence & donc, par conséquent, c'est pourquoi, ainsi, dès lors \\
\hline
Reformulation & c'est-à-dire, autrement dit, en d'autres termes \\
\hline
\end{tabularx}
\end{center}

\begin{exemplebox}[Opposition dans \emph{Le Vieux Nègre et la médaille}]
\og Meka reçoit une médaille ; \textbf{cependant}, cette distinction le conduit à l'humiliation. \fg{} Le connecteur \emph{cependant} signale nettement l'opposition entre l'honneur apparent et la réalité vécue : le lecteur comprend immédiatement que la médaille est un leurre. Sans ce mot, le rapport entre les deux idées resterait flou.
\end{exemplebox}

\begin{popfigure}
\begin{center}
\begin{tikzpicture}[
  cat/.style={rectangle, rounded corners, draw=popdark, fill=popcream, align=center, minimum width=3.4cm, minimum height=0.85cm, font=\bfseries},
  mot/.style={align=center, font=\small},
  lien/.style={-{Stealth[length=2.5mm]}, thick, popmuted}
]
\node[cat, fill=popblueL, draw=popblue] (add) at (0,3.2) {Addition};
\node[mot, align=center] at (0,2.2){et, ni\\ d'ailleurs, en outre, de plus};
\node[cat, fill=poporangeL, draw=poporange] (opp) at (5.2,3.2) {Opposition};
\node[mot, align=center] at (5.2,2.2){mais, bien que\\ cependant, néanmoins, en revanche};
\node[cat, fill=popgreenL, draw=popgreen] (cau) at (0,0.6) {Cause};
\node[mot, align=center] at (0,-0.4){car, parce que, puisque\\ en effet};
\node[cat, fill=poppurpleL, draw=poppurple] (cons) at (5.2,0.6) {Conséquence};
\node[mot, align=center] at (5.2,-0.4){donc, si bien que\\ par conséquent, c'est pourquoi};
\node[cat, fill=popgoldL, draw=popgold] (temps) at (10.4,3.2) {Temps / Ordre};
\node[mot, align=center] at (10.4,2.2){lorsque, quand\\ d'abord, ensuite, enfin};
\node[cat, fill=poppinkL, draw=poppink] (ref) at (10.4,0.6) {Reformulation};
\node[mot, align=center] at (10.4,-0.4){c'est-à-dire\\ autrement dit};
\end{tikzpicture}
\legende{Classification des connecteurs logiques par fonction : addition, opposition, cause, conséquence, temps/ordre et reformulation. Chaque famille regroupe conjonctions, adverbes et locutions de même valeur.}
\end{center}
\end{popfigure}

\courssub{Les pronoms relatifs : reprendre et enchaîner les idées}

\begin{definition}[Pronom relatif]
Le \cle{pronom relatif} (\textbf{qui, que, où, dont, lequel} et ses formes composées : \emph{lequel, laquelle, lesquels, lesquelles, auquel, auxquels, auquel, auxquelles, duquel, desquels, desquelles}) remplace un nom déjà mentionné (l'\textbf{antécédent}) et introduit une \textbf{proposition subordonnée relative}. Il évite les répétitions et resserre la phrase.
\end{definition}

Comparez : \og Meka reçoit une médaille. Cette médaille est un symbole trompeur. \fg{} et \og Meka reçoit une médaille \textbf{qui} est un symbole trompeur. \fg{} La seconde version, plus dense et plus élégante, est celle que l'on attend en dissertation.

\begin{center}
\begin{tabularx}{\linewidth}{|l|l|X|}
\hline
\rowcolor{popblueL} \textbf{Pronom} & \textbf{Fonction} & \textbf{Exemple} \\
\hline
\cle{qui} & sujet & Le roman \textbf{qui} dénonce le colonialisme... \\
\hline
\cle{que} & COD & L'argument \textbf{que} l'auteur développe... \\
\hline
\cle{où} & complément circonstanciel de lieu ou de temps & Le village \textbf{où} vit Meka... ; L'année \textbf{où} il est décoré... \\
\hline
\cle{dont} & complément de l'antécédent introduit par \og de \fg{} (objet prépositionnel) & L'humiliation \textbf{dont} Meka est victime... (on dit : victime \emph{de} quelque chose). \\
\hline
\cle{lequel} & complément introduit par une préposition \emph{autre que de} (\emph{à, pour, sur, avec, sans...}) & Le prétexte \textbf{sur lequel} repose la médaille... \\
\hline
\end{tabularx}
\end{center}

\begin{attention}[Accord et choix du relatif]
\begin{itemize}
\item Avec \textbf{que} (COD), le participe passé s'accorde avec l'antécédent placé \textbf{avant} : \og les médailles \emph{que} le commandant a \emph{décernées} \fg{}.
\item Ne confondez pas \textbf{où} (lieu/temps) et \textbf{ou} (alternative), ni \textbf{dont} et \textbf{que} : \og l'homme \emph{dont} je parle \fg{} car on dit \og parler \emph{de} quelqu'un \fg{}.
\item \textbf{Lequel} et ses formes s'accordent en genre et en nombre avec l'antécédent.
\end{itemize}
\end{attention}

\courssub{Méthode : choisir le bon connecteur}

\begin{methode}[Choisir le connecteur adapté]
\begin{enumerate}
\item \textbf{Identifier la relation logique} entre les deux idées : s'agit-il d'une opposition, d'une cause, d'une conséquence, d'une addition, d'un ordre chronologique ?
\item \textbf{Choisir la famille de connecteurs} correspondante à l'aide du tableau de classification.
\item \textbf{Vérifier le niveau de langue} : pour un texte formel (dissertation), préférer \og en revanche \fg{} à \og par contre \fg{}, \og par conséquent \fg{} à \og alors \fg{}, \og toutefois \fg{} à \og par contre \fg{}.
\item \textbf{Vérifier la construction grammaticale} : subjonctif après \og bien que / afin que \fg{}, place de \og car \fg{} (jamais en tête de phrase), virgule après un adverbe de liaison en tête de phrase.
\item \textbf{Relire} : supprimer le connecteur ; si le sens devient ambigu, il est utile ; sinon, il alourdit.
\end{enumerate}
\end{methode}

\begin{exoresolu}[Relier un paragraphe argumentatif]
Énoncé : reliez les phrases suivantes par les connecteurs appropriés pour obtenir un paragraphe argumentatif cohérent sur le sujet : \og Faut-il interdire le téléphone portable à l'école ? \fg{}

\og Le téléphone portable favorise l'accès à l'information. Il constitue une source de distraction permanente. Les élèves consultent les réseaux sociaux en classe. Leurs résultats scolaires baissent. Une régulation stricte s'impose. \fg{}
\tcblower
Solution détaillée : \og Le téléphone portable favorise l'accès à l'information ; \textbf{cependant}, il constitue une source de distraction permanente. \textbf{En effet}, les élèves consultent les réseaux sociaux en classe, \textbf{si bien que} leurs résultats scolaires baissent. \textbf{Par conséquent}, une régulation stricte s'impose. \fg{}

Analyse des choix : \emph{cependant} marque l'opposition entre l'avantage et l'inconvénient ; \emph{en effet} introduit l'explication de la distraction ; \emph{si bien que} exprime la conséquence ; \emph{par conséquent} conclut le raisonnement. Le paragraphe suit désormais une progression logique visible : thèse, nuance, preuve, conséquence, conclusion.
\end{exoresolu}

\begin{savaistu}[La langue compte dans la note]
Au Baccalauréat technique, la correction grammaticale, l'orthographe et la cohérence textuelle --- dont l'usage pertinent des connecteurs --- représentent une part notable de la note de dissertation. Deux copies au contenu comparable peuvent être séparées de plusieurs points uniquement par la qualité de la langue et des liaisons. Maîtriser les connecteurs, c'est donc gagner des points sans apprendre un seul argument supplémentaire.
\end{savaistu}

\begin{exercice}[À vous de jouer]
Complétez le paragraphe suivant avec les connecteurs de votre choix, en justifiant chaque choix par la relation logique visée : \og Les réseaux sociaux permettent aux jeunes Camerounais de s'informer rapidement. \dots{}, ils exposent aussi à la désinformation. \dots{}, de nombreux élèves croient des rumeurs sans les vérifier. \dots{}, leur esprit critique s'affaiblit. \dots{}, il appartient à l'école d'enseigner l'éducation aux médias. \fg{}
\end{exercice}

\begin{corrige}[Exercice]
Proposition de corrigé : \og Les réseaux sociaux permettent aux jeunes Camerounais de s'informer rapidement. \textbf{Cependant}, ils exposent aussi à la désinformation. \textbf{En effet}, de nombreux élèves croient des rumeurs sans les vérifier. \textbf{Par conséquent}, leur esprit critique s'affaiblit. \textbf{C'est pourquoi} il appartient à l'école d'enseigner l'éducation aux médias. \fg{} Justifications : \emph{cependant} = opposition ; \emph{en effet} = explication ; \emph{par conséquent} et \emph{c'est pourquoi} = conséquence. D'autres choix sont acceptables (\emph{en revanche}, \emph{ainsi}) s'ils respectent la relation logique.
\end{corrige}

\begin{aretenir}[L'essentiel]
\begin{itemize}
\item Un \cle{connecteur logique} établit une relation logique (addition, opposition, cause, conséquence, temps) entre deux énoncés.
\item Les \textbf{conjonctions de coordination} (mais, ou, et, donc, or, ni, car) relient des éléments de même rang ; \textbf{car} ne relie que des propositions.
\item Les \textbf{conjonctions de subordination} (que, parce que, bien que, si, lorsque, puisque, afin que...) créent une hiérarchie entre propositions ; \textbf{bien que} et \textbf{afin que} exigent le subjonctif.
\item Les \textbf{adverbes et locutions de liaison} (cependant, en effet, par conséquent, en revanche) articulent phrases et paragraphes de la dissertation.
\item Les \textbf{pronoms relatifs} (qui, que, où, dont, lequel) évitent les répétitions et resserrent le texte ; accord du participe passé avec \textbf{que} (COD antéposé).
\item Le bon connecteur se choisit en identifiant d'abord la relation logique visée, puis en respectant le niveau de langue formel et la construction grammaticale.
\end{itemize}
\end{aretenir}

\coursec{Rédaction de l'introduction et de la conclusion}

Dans la section précédente, nous avons appris à \cle{relier} nos phrases et nos idées grâce aux conjonctions, aux adverbes de liaison et aux pronoms relatifs. Ces outils trouvent leur emploi le plus noble dans la dissertation : c'est en effet grâce à eux que l'introduction enchaîne harmonieusement ses étapes et que la conclusion boucle la réflexion. Nous allons maintenant apprendre à construire ces deux moments capitaux de la copie : l'\cle{introduction}, qui ouvre la réflexion, et la \cle{conclusion}, qui la referme.

\courssub{Fonction de l'introduction}

Imaginez un receveur de copies qui corrige deux cents dissertations en quelques jours. Les dix premières lignes de votre copie sont votre poignée de main : elles donnent immédiatement une impression de maîtrise ou de maladresse. L'introduction n'est donc pas un préambule décoratif : elle remplit trois fonctions précises.

\begin{aretenir}[Les trois fonctions de l'introduction]
\begin{enumerate}
\item \textbf{Capter l'attention} du lecteur par une entrée en matière intéressante (l'amorce).
\item \textbf{Poser le problème} : présenter le sujet et formuler la question à laquelle toute la dissertation va répondre (la problématique).
\item \textbf{Annoncer la démarche} : indiquer les grandes étapes du raisonnement (l'annonce du plan).
\end{enumerate}
\end{aretenir}

\courssub{Les quatre étapes de l'introduction}

Une introduction de dissertation se construit comme un \cle{entonnoir} : on part du général et l'on resserre progressivement jusqu'au particulier, c'est-à-dire jusqu'au plan précis de votre devoir. Elle comporte quatre étapes obligatoires.

\begin{enumerate}
\item \textbf{L'amorce (ou accroche)} : une phrase générale qui met le lecteur dans le bain du thème.
\item \textbf{La présentation du sujet} : on situe le sujet, on en définit éventuellement les termes importants, on montre son intérêt.
\item \textbf{La problématique} : la question centrale, formulée clairement, souvent sous forme interrogative.
\item \textbf{L'annonce du plan} : on présente les deux ou trois parties du développement.
\end{enumerate}

\begin{popfigure}
\begin{center}
\begin{tikzpicture}[scale=1]
\fill[popblue!15] (0,4.4) -- (8,4.4) -- (7.2,3.3) -- (0.8,3.3) -- cycle;
\fill[popblue!30] (0.8,3.3) -- (7.2,3.3) -- (6.4,2.2) -- (1.6,2.2) -- cycle;
\fill[popblue!50] (1.6,2.2) -- (6.4,2.2) -- (5.6,1.1) -- (2.4,1.1) -- cycle;
\fill[popblue!75] (2.4,1.1) -- (5.6,1.1) -- (4.8,0) -- (3.2,0) -- cycle;
\node[text=popdark,font=\bfseries] at (4,3.85) {1. Amorce (général)};
\node[text=popdark,font=\bfseries] at (4,2.75) {2. Présentation du sujet};
\node[text=white,font=\bfseries] at (4,1.65) {3. Problématique};
\node[text=white,font=\bfseries] at (4,0.55) {4. Annonce du plan};
\draw[->,thick,popdark] (8.6,4.2) -- (8.6,0.2) node[midway,right,align=left,font=\small] {du général\\vers le\\particulier};
\end{tikzpicture}
\end{center}
\legende{L'entonnoir de l'introduction : chaque étape resserre le propos, de l'idée la plus large (l'amorce) jusqu'à la plus précise (le plan).}
\end{popfigure}

\courssub{Techniques d'amorce}

L'amorce est la première phrase de la copie : elle doit être juste, élégante et en rapport avec le thème. Quatre techniques sont à votre disposition.

\begin{center}
\begin{tabularx}{\linewidth}{|l|X|}
\hline
\rowcolor{popblueL} \textbf{Technique} & \textbf{Exemple d'amorce} \\
\hline
Fait d'actualité & « À chaque cérémonie de fin d'année, les tableaux d'honneur distribuent médailles et prix aux élèves méritants. » \\
\hline
Citation & « "La récompense flatte celui qui la reçoit autant qu'elle interroge celui qui la donne", écrit un moraliste. » \\
\hline
Constat général & « De tout temps, les sociétés humaines ont décerné des distinctions à ceux qu'elles voulaient honorer. » \\
\hline
Anecdote & « Un vieux soldat, décoré après des années d'oubli, aurait répondu : "Cette médaille arrive trop tard pour sécher mes larmes." » \\
\hline
\end{tabularx}
\end{center}

\begin{attention}[Pièges de l'amorce]
Évitez les amorces trop vagues et sans rapport avec le sujet (« Depuis la nuit des temps, l'homme... »), les définitions de dictionnaire (« Le Petit Robert définit la médaille comme... ») et les affirmations fausses ou invérifiables. L'amorce doit mener \emph{naturellement} au sujet.
\end{attention}

\courssub{La problématique}

Après l'amorce et la présentation du sujet vient le cœur de l'introduction : la problématique.

\begin{definition}[La problématique]
La \cle{problématique} est la \textbf{question centrale} qui oriente toute la réflexion et à laquelle la dissertation tout entière se propose de répondre. Elle naît d'une tension, d'un doute ou d'une contradiction contenus dans le sujet. Elle peut s'accompagner de \textbf{questions secondaires} qui préparent les parties du plan.
\end{definition}

Pour bien formuler une problématique, il faut d'abord repérer le \emph{problème} que cache le sujet. Prenons le sujet : « La médaille peut-elle réparer une injustice ? » Le problème est le suivant : d'un côté, la médaille est un honneur qui reconnaît un mérite ; de l'autre, une injustice passée (souffrance, humiliation) semble impossible à effacer par un simple objet symbolique. La problématique s'écrit alors naturellement sous forme interrogative : « Une distinction honorifique, si tardive soit-elle, suffit-elle à effacer le tort subi, ou n'est-elle qu'un geste symbolique impuissant à réparer le passé ? »

\begin{methode}[Rédiger une introduction en quatre étapes]
\begin{enumerate}
\item \textbf{Relisez le sujet} et repérez ses mots-clés ; choisissez une technique d'amorce adaptée au thème.
\item \textbf{Rédigez l'amorce} (une ou deux phrases), puis enchaînez avec la présentation du sujet : situez-le, définissez les termes essentiels, montrez son actualité ou son intérêt. Utilisez un connecteur (« En effet », « Ainsi », « Or »).
\item \textbf{Formulez la problématique} : faites apparaître la tension du sujet (« Mais... », « Cependant... ») puis posez la question principale, éventuellement suivie d'une question secondaire.
\item \textbf{Annoncez le plan} en une ou deux phrases élégantes et impersonnelles, qui laissent deviner les deux ou trois parties du développement.
\end{enumerate}
\end{methode}

\begin{exemplebox}[Introduction rédigée et commentée]
\textbf{Sujet : « L'école garantit-elle encore la réussite sociale au Cameroun ? »}

\emph{Texte de l'introduction :}

« De tout temps, l'école a été présentée comme l'ascenseur social par excellence, celui qui permet à l'enfant de paysan de devenir médecin, magistrat ou ingénieur. Au Cameroun, nombre de familles consentent d'immenses sacrifices pour scolariser leurs enfants, convaincues que le diplôme ouvre toutes les portes. Or, face au chômage des jeunes diplômés et à la réussite éclatante de certains autodidactes, cette belle certitude semble aujourd'hui ébranlée. L'école garantit-elle donc encore la réussite sociale au Cameroun ? N'est-elle pas devenue, pour beaucoup, une promesse incertaine ? Pour répondre à cette question, nous verrons d'abord que l'école demeure un instrument irremplaçable de promotion sociale, avant de montrer qu'elle ne suffit plus, à elle seule, à garantir la réussite. »

\emph{Commentaire :} la première phrase est une amorce par constat général ; les deux phrases suivantes présentent le sujet et son actualité ; le connecteur « Or » introduit la tension ; la problématique est posée clairement, avec une question secondaire ; la dernière phrase annonce un plan en deux parties de manière impersonnelle.
\end{exemplebox}

\begin{attention}[Erreur fréquente : l'annonce de plan maladroite]
N'écrivez jamais : « \emph{Je vais d'abord parler de..., ensuite je parlerai de..., enfin je conclurai.} » Cette annonce à la première personne est scolaire et disgracieuse. Préférez une formulation élégante et impersonnelle : « \emph{Nous verrons dans un premier temps que..., puis nous nous demanderons si...} » ou « \emph{Il conviendra d'examiner..., avant d'envisager...} ».
\end{attention}

\courssub{Fonction et construction de la conclusion}

Si l'introduction est un entonnoir, la conclusion est un \cle{entonnoir inversé} : on part du particulier (ce que le devoir a démontré) pour revenir vers le général (une ouverture sur le monde). Elle remplit trois fonctions.

\begin{aretenir}[Les trois fonctions de la conclusion]
\begin{enumerate}
\item \textbf{Faire le bilan} : résumer de façon synthétique le chemin parcouru, \emph{sans répéter mot à mot} les phrases du développement.
\item \textbf{Répondre à la problématique} : donner clairement votre réponse à la question posée dans l'introduction.
\item \textbf{Ouvrir la réflexion} : terminer par une question ou une perspective nouvelle, en lien avec le sujet, qui prolonge la pensée sans la contredire.
\end{enumerate}
\end{aretenir}

\begin{popfigure}
\begin{center}
\begin{tikzpicture}[scale=1]
\fill[poporange!70] (2.4,3.3) -- (5.6,3.3) -- (6.4,2.2) -- (1.6,2.2) -- cycle;
\fill[poporange!45] (1.6,2.2) -- (6.4,2.2) -- (7.2,1.1) -- (0.8,1.1) -- cycle;
\fill[poporange!20] (0.8,1.1) -- (7.2,1.1) -- (8,0) -- (0,0) -- cycle;
\node[text=white,font=\bfseries] at (4,2.75) {1. Bilan synthétique};
\node[text=popdark,font=\bfseries] at (4,1.65) {2. Réponse à la problématique};
\node[text=popdark,font=\bfseries] at (4,0.55) {3. Ouverture};
\draw[->,thick,popdark] (8.6,3.1) -- (8.6,0.2) node[midway,right,align=left,font=\small] {du particulier\\vers le\\général};
\end{tikzpicture}
\end{center}
\legende{L'entonnoir inversé de la conclusion : on s'élargit du bilan précis du devoir vers une perspective générale.}
\end{popfigure}

\begin{attention}[Erreur fréquente : l'argument nouveau]
La conclusion ne doit \textbf{jamais} introduire un argument ou un exemple nouveau : tout argument appartient au développement. Si une idée vous revient au moment de conclure, soit elle est déjà traitée (alors contentez-vous de la rappeler brièvement), soit elle est absente (alors il est trop tard : ne l'ajoutez pas). La conclusion \emph{rassemble}, elle n'\emph{ajoute} pas.
\end{attention}

\courssub{Travaux pratiques guidés : « La médaille peut-elle réparer une injustice ? »}

Appliquons la méthode au sujet tiré de notre lecture du \emph{Vieux nègre et la médaille} de Ferdinand Oyono. Dans ce récit, Meka, vieillard africain, reçoit une médaille des mains de l'administration coloniale en échange de terres cédées et de fils morts à la guerre : cette décoration prétend « récompenser » des sacrifices immenses et irréparables.

\begin{exoresolu}[Rédaction guidée d'une introduction]
Rédiger l'introduction de la dissertation sur le sujet : « La médaille peut-elle réparer une injustice ? »
\tcblower
\textbf{Étape 1 — Amorce (constat général) :} « De tout temps, les hommes ont cherché à réparer les torts subis en offrant des compensations : honneurs, récompenses, distinctions. »

\textbf{Étape 2 — Présentation du sujet :} « C'est précisément cette logique qui préside, dans \emph{Le Vieux nègre et la médaille} de Ferdinand Oyono, à la remise d'une médaille à Meka, vieillard qui a donné ses terres à la mission et ses deux fils à la guerre. Cette décoration se veut le signe de la reconnaissance coloniale. »

\textbf{Étape 3 — Problématique :} « Or peut-on échanger des vies humaines et des souffrances contre un morceau de métal doré ? La médaille peut-elle réparer une injustice, ou n'est-elle qu'un leurre destiné à faire oublier l'irréparable ? »

\textbf{Étape 4 — Annonce du plan :} « Nous montrerons d'abord que la médaille possède une véritable valeur symbolique de reconnaissance, avant de démontrer qu'elle demeure impuissante à effacer l'injustice subie. »
\end{exoresolu}

\begin{exoresolu}[Rédaction guidée d'une conclusion]
Rédiger la conclusion de la même dissertation.
\tcblower
\textbf{Étape 1 — Bilan synthétique :} « Au terme de notre réflexion, il apparaît que la médaille, en tant que signe officiel de gratitude, honore celui qui la reçoit et reconnaît publiquement son sacrifice. »

\textbf{Étape 2 — Réponse à la problématique :} « Cependant, comme le montre le destin tragique de Meka, humilié le soir même de sa décoration, aucune médaille ne saurait ressusciter les morts, restituer les terres ni laver l'humiliation : elle ne répare pas l'injustice, elle ne fait que la maquiller. »

\textbf{Étape 3 — Ouverture :} « Reste alors à se demander si la véritable réparation ne passerait pas plutôt par la justice, la restitution et la mémoire, plutôt que par les symboles. »
\end{exoresolu}

\begin{exercice}[À vous de jouer]
Rédigez une introduction complète (quatre étapes) et une conclusion complète (trois étapes) pour le sujet suivant : « Les récompenses scolaires (prix, tableaux d'honneur) sont-elles utiles aux élèves ? » Vous choisirez une technique d'amorce différente de celle du corrigé, et vous veillerez à annoncer le plan de manière impersonnelle.
\end{exercice}

\begin{corrige}[Exercice — éléments de réponse]
\emph{Introduction attendue :} amorce possible par fait d'actualité (« À chaque fin d'année scolaire, les cérémonies de remise de prix rassemblent élèves, parents et enseignants... ») ; présentation du sujet (les récompenses sont une tradition scolaire destinée à encourager le mérite) ; problématique (« Mais ces distinctions motivent-elles réellement les élèves, ou créent-elles jalousies et découragements ? ») ; annonce impersonnelle du plan (« Nous verrons d'abord que les récompenses stimulent l'effort, puis qu'elles peuvent aussi produire des effets néfastes. »).

\emph{Conclusion attendue :} bilan (les récompenses encouragent le travail mais peuvent frustrer les non-récompensés) ; réponse nuancée (elles sont utiles si elles valorisent le progrès de tous et non seulement les meilleurs) ; ouverture (« Ne faudrait-il pas repenser les formes de reconnaissance scolaire pour qu'aucun élève ne se sente exclu ? »).
\end{corrige}

\begin{savaistu}[Le pouvoir des premières lignes]
Les études sur la correction des copies le confirment : une introduction claire, bien construite et sans faute conditionne \emph{favorablement} le correcteur dès les premières lignes. C'est ce qu'on appelle l'« effet de halo » : un lecteur séduit par le début lit la suite avec bienveillance, tandis qu'une introduction brouillonne l'incite à la sévérité. Au Baccalauréat technique, soigner l'introduction et la conclusion, c'est donc déjà gagner des points avant même que le fond ne soit examiné !
\end{savaistu}

En résumé, l'introduction descend en entonnoir de l'amorce vers le plan, et la conclusion remonte en entonnoir inversé du bilan vers l'ouverture. Ces deux paragraphes encadrent le développement comme deux piliers encadrent un pont. Il nous reste maintenant à apprendre à construire les arches de ce pont : les paragraphes du développement, avec leur idée directrice, leurs arguments et leurs exemples — ce sera l'objet de la section suivante.

\coursec{Rédaction du paragraphe argumentatif}

Dans la section précédente, nous avons appris à construire l'introduction, qui présente le sujet et annonce le plan, et la conclusion, qui fait le bilan de la démonstration. Entre ces deux moments stratégiques se trouve le cœur de la dissertation : le développement. Or le développement n'est pas un bloc continu : il est constitué de \cle{paragraphes}, chacun portant une étape précise du raisonnement. Savoir rédiger un bon paragraphe argumentatif, c'est donc maîtriser la brique élémentaire de toute la dissertation. C'est l'objet de cette section.

\courssub{Qu'est-ce qu'un paragraphe argumentatif ?}

\begin{definition}[Le paragraphe argumentatif]
Le \cle{paragraphe argumentatif} est une unité de texte, matérialisée par un alinéa, qui développe \textbf{une seule idée} au service de la démonstration. Il obéit à une structure en quatre temps : l'\textbf{idée directrice} (phrase d'ouverture), l'\textbf{argument} (explication et justification), l'\textbf{exemple} (illustration concrète) et la \textbf{mini-conclusion} ou \textbf{transition} (bilan et lien avec la suite).
\end{definition}

Pour bien comprendre, faisons une comparaison simple. Une dissertation ressemble à une maison : l'introduction est le porche, la conclusion est le toit, et les paragraphes sont les murs. Un mur mal construit fragilise toute la maison ; de même, un paragraphe confus affaiblit toute la démonstration. Le correcteur du Probatoire ou du Baccalauréat technique lit votre copie paragraphe par paragraphe : chacun doit être clair, complet et solidaire des autres.

\begin{definition}[L'idée directrice]
L'\cle{idée directrice} est la phrase qui énonce l'idée principale du paragraphe et qui s'inscrit dans la démonstration. Elle se place en tête du paragraphe et annonce ce que l'on va prouver. Elle ne doit jamais être un simple mot ni une question : c'est une phrase complète, affirmative, clairement reliée à la problématique.
\end{definition}

\begin{exemplebox}[Idée directrice réussie et idée directrice ratée]
Sujet : \emph{La littérature peut-elle dénoncer les injustices sociales ?}
\begin{itemize}
\item \textbf{Ratée} : « La dénonciation. » (ce n'est pas une phrase ; on ne sait pas ce que le paragraphe va prouver).
\item \textbf{Réussie} : « La littérature constitue d'abord une arme efficace de dénonciation, car l'écrivain, par le récit, révèle au public des injustices que le pouvoir préférerait cacher. »
\end{itemize}
La seconde formulation annonce clairement l'idée du paragraphe et prépare l'argument qui va suivre.
\end{exemplebox}

\courssub{La structure en quatre temps : la méthode I-A-E-T}

\begin{methode}[Construire un paragraphe en quatre temps : I-A-E-T]
Pour rédiger un paragraphe argumentatif complet, suivez toujours les quatre étapes de la méthode \cle{I-A-E-T} :
\begin{enumerate}
\item \textbf{I — Idée directrice} : formulez en une phrase l'idée du paragraphe, en la reliant à la problématique du sujet.
\item \textbf{A — Argument} : expliquez et justifiez cette idée. Répondez aux questions : \emph{pourquoi ? comment ? en quoi ?} Montrez le lien logique entre l'idée et la thèse défendue.
\item \textbf{E — Exemple} : illustrez l'argument par un cas précis et pertinent : un passage de l'œuvre étudiée (\emph{Le Vieux Nègre et la Médaille} de Ferdinand Oyono), un fait de l'actualité camerounaise, une référence de culture générale.
\item \textbf{T — Transition} : concluez brièvement le paragraphe et annoncez, grâce à un connecteur logique, l'idée du paragraphe suivant.
\end{enumerate}
\end{methode}

\begin{popfigure}
\begin{center}
\begin{tikzpicture}[
  couche/.style={rectangle, draw=popblue, thick, fill=popblueL, rounded corners=3pt,
                 minimum width=9.5cm, minimum height=0.95cm, align=center, font=\small},
  fleche/.style={-{Stealth[length=3mm]}, thick, poporange}
]
\node[couche] (i) at (0,0) {\textbf{I — Idée directrice} : la phrase d'ouverture qui annonce l'idée};
\node[couche, fill=poporangeL, draw=poporange] (a) at (0,-1.5) {\textbf{A — Argument} : explication, justification, lien avec la problématique};
\node[couche, fill=popgreenL, draw=popgreen] (e) at (0,-3) {\textbf{E — Exemple} : illustration précise (œuvre, actualité, culture)};
\node[couche, fill=poppurpleL, draw=poppurple] (t) at (0,-4.5) {\textbf{T — Transition} : mini-conclusion et lien vers le paragraphe suivant};
\draw[fleche] (i.south) -- (a.north);
\draw[fleche] (a.south) -- (e.north);
\draw[fleche] (e.south) -- (t.north);
\node[font=\small\itshape, popmuted] at (0,-5.5) {Un paragraphe = une seule idée, développée en quatre temps};
\end{tikzpicture}
\end{center}
\legende{Schéma de la structure en couches du paragraphe argumentatif (méthode I-A-E-T) : chaque couche repose sur la précédente, comme les assises d'un mur.}
\end{popfigure}

\courssub{Développer l'argument : expliquer, justifier, relier}

L'argument est le cœur du paragraphe : c'est lui qui \emph{prouve}. Beaucoup d'élèves se contentent d'affirmer ; or affirmer n'est pas démontrer. Pour développer un argument, il faut procéder en trois mouvements.

\begin{enumerate}
\item \textbf{Expliquer} : reformuler l'idée directrice en la décomposant, en définissant ses termes si nécessaire. « Autrement dit… », « En effet… », « C'est-à-dire… ».
\item \textbf{Justifier} : donner la raison, la cause, le mécanisme. « Car… », « puisque… », « dans la mesure où… ». C'est ici que les conjonctions de subordination étudiées dans cette unité prennent toute leur utilité.
\item \textbf{Relier à la problématique} : montrer explicitement que cet argument fait avancer la réponse à la question posée par le sujet. « Ainsi… », « C'est pourquoi… », « On comprend dès lors que… ».
\end{enumerate}

\begin{exemplebox}[Un argument développé]
Idée directrice : « Le rire, chez Ferdinand Oyono, est une arme de dénonciation. »
Développement : « \textbf{En effet}, l'ironie permet à l'écrivain de critiquer sans s'exposer directement à la censure. \textbf{Car} en ridiculisant les colonisateurs, il révèle leur injustice tout en restant dans le domaine de la fiction. \textbf{Ainsi}, le comique n'est pas gratuit : il sert la contestation, ce qui montre bien que la littérature peut dénoncer les injustices sociales. »
\end{exemplebox}

\courssub{Choisir et exploiter l'exemple}

L'exemple rend l'argument concret et crédible. Mais un exemple ne vaut que s'il respecte trois conditions :

\begin{itemize}
\item \textbf{Précis} : citez un passage, un personnage, une scène identifiable, et non une allusion vague. Préférez « Meka, dans \emph{Le Vieux Nègre et la Médaille}, reçoit une médaille pour avoir donné ses terres et ses fils à la France » à « dans le roman, il y a une injustice ».
\item \textbf{Pertinent} : l'exemple doit illustrer exactement l'argument, pas un autre. Avant de l'écrire, demandez-vous : « cet exemple prouve-t-il bien ce que je viens de dire ? »
\item \textbf{Commenté} : après l'exemple, ajoutez une ou deux phrases qui montrent \emph{en quoi} il confirme l'argument. Un exemple abandonné à lui-même ne prouve rien.
\end{itemize}

Les sources d'exemples autorisées sont variées : l'œuvre au programme (\emph{Le Vieux Nègre et la Médaille}), d'autres œuvres africaines ou françaises étudiées, l'actualité camerounaise (débats publics, faits de société), l'histoire et la culture générale. Varier les sources témoigne d'une culture riche, très appréciée des correcteurs.

\courssub{La transition : assurer la continuité de la démonstration}

La dernière phrase du paragraphe a une double fonction : elle \textbf{boucle} l'idée (mini-conclusion) et elle \textbf{ouvre} sur la suite (transition). C'est ici que les outils de liaison étudiés dans cette unité — adverbes de liaison, locutions adverbiales, pronoms relatifs — deviennent indispensables.

\begin{center}
\begin{tabularx}{\linewidth}{|l|X|}\hline
\rowcolor{popblueL} \textbf{Fonction} & \textbf{Connecteurs utiles} \\ \hline
Addition & de plus, en outre, également, par ailleurs \\ \hline
Opposition & cependant, néanmoins, en revanche, pourtant \\ \hline
Cause & car, en effet, puisque, dans la mesure où \\ \hline
Conséquence & ainsi, donc, c'est pourquoi, par conséquent \\ \hline
Progression & d'abord, ensuite, enfin, non seulement… mais encore \\ \hline
\end{tabularx}
\end{center}

\begin{exemplebox}[Une transition réussie]
« Ainsi, la médaille de Meka apparaît comme le symbole dérisoire d'une reconnaissance trompeuse. \textbf{Cependant}, la dénonciation de l'injustice coloniale ne passe pas seulement par le symbole : elle s'exprime aussi à travers le portrait des personnages. » — La première phrase conclut le paragraphe sur le symbole ; « Cependant » annonce le paragraphe suivant, consacré aux personnages.
\end{exemplebox}

\courssub{Un paragraphe complet annoté}

\begin{exoresolu}[Rédiger un paragraphe à partir d'un plan]
\textbf{Énoncé.} Sujet : \emph{La littérature peut-elle dénoncer les injustices sociales ?} À partir du plan de paragraphe suivant, rédigez un paragraphe complet (8 à 12 lignes) : Idée : Oyono dénonce l'hypocrisie coloniale. Argument : la médaille récompense la soumission, pas le mérite. Exemple : Meka perd ses terres et ses fils et ne reçoit qu'une médaille.
\tcblower
\textbf{Solution détaillée.} « La littérature permet d'abord de dénoncer l'hypocrisie des discours officiels, comme le montre Ferdinand Oyono dans \emph{Le Vieux Nègre et la Médaille}. En effet, la médaille remise au vieux Meka ne récompense pas un mérite réel : elle couronne au contraire sa soumission et ses sacrifices. Car le vieil homme a donné ses terres à la mission et ses deux fils à la guerre, sans recevoir en échange la moindre contrepartie véritable. Ainsi, lors de la cérémonie de remise de la médaille, le lecteur mesure l'écart scandaleux entre les grands discours des autorités coloniales et la misère du vieux nègre, écart qui révèle toute l'hypocrisie du système. Cette scène prouve que le roman peut démasquer les injustices que le pouvoir dissimule sous les honneurs. Cependant, Oyono ne dénonce pas seulement par les symboles : il utilise aussi le rire comme arme de contestation. »
\par\smallskip
Ce paragraphe respecte la méthode I-A-E-T : idée directrice en tête, argument expliqué (« En effet… Car… »), exemple précis commenté, transition ouvrant sur le paragraphe suivant (« Cependant… »).
\end{exoresolu}

\begin{popfigure}
\begin{center}
\begin{tikzpicture}[
  zone/.style={rectangle, draw=popline, rounded corners=2pt, align=justify,
               text width=8.6cm, inner sep=5pt, font=\footnotesize},
  tag/.style={rectangle, draw=popdark, fill=popgoldL, rounded corners=2pt,
              font=\footnotesize\bfseries, inner sep=3pt},
  lien/.style={-{Stealth[length=2.5mm]}, thick, popdark}
]
\node[zone, fill=popblueL] (p1) at (0,0) {La littérature permet d'abord de dénoncer l'hypocrisie des discours officiels, comme le montre Oyono dans \emph{Le Vieux Nègre et la Médaille}.};
\node[zone, fill=poporangeL] (p2) at (0,-2.1) {En effet, la médaille remise au vieux Meka ne récompense pas un mérite réel : elle couronne sa soumission et ses sacrifices, car il a donné ses terres et ses deux fils sans contrepartie véritable.};
\node[zone, fill=popgreenL] (p3) at (0,-4.4) {Ainsi, lors de la cérémonie, le lecteur mesure l'écart scandaleux entre les grands discours des autorités et la misère du vieux nègre, écart qui révèle l'hypocrisie du système.};
\node[zone, fill=poppurpleL] (p4) at (0,-6.5) {Cependant, Oyono ne dénonce pas seulement par les symboles : il utilise aussi le rire comme arme de contestation.};
\node[tag] (t1) at (7.6,0) {Idée directrice};
\node[tag] (t2) at (7.6,-2.1) {Argument};
\node[tag] (t3) at (7.6,-4.4) {Exemple commenté};
\node[tag] (t4) at (7.6,-6.5) {Transition};
\draw[lien] (t1.west) -- (p1.east);
\draw[lien] (t2.west) -- (p2.east);
\draw[lien] (t3.west) -- (p3.east);
\draw[lien] (t4.west) -- (p4.east);
\end{tikzpicture}
\end{center}
\legende{Paragraphe annoté : chaque composante de la méthode I-A-E-T est identifiée par une flèche et une étiquette.}
\end{popfigure}

\courssub{Longueur, équilibre et erreurs à éviter}

\begin{exemplebox}[Une question de mesure]
Un paragraphe de dissertation efficace compte en moyenne \textbf{8 à 12 lignes}. En dessous de 5 lignes, l'idée n'est pas développée (l'argument ou l'exemple manque) ; au-delà de 15 lignes, le paragraphe contient probablement deux idées qu'il faut séparer en deux paragraphes distincts.
\end{exemplebox}

\begin{attention}[Erreurs fréquentes à éviter]
\begin{itemize}
\item \textbf{Accumuler les exemples sans argument} : une liste d'exemples n'est pas une démonstration. Un seul exemple bien commenté vaut mieux que trois exemples jetés en vrac.
\item \textbf{Écrire des paragraphes d'une seule phrase} : un paragraphe réduit à une ligne signifie que l'idée n'est pas développée ; c'est une simple affirmation.
\item \textbf{Mettre plusieurs idées dans un même paragraphe} : une idée = un paragraphe. Si vous changez d'idée, faites un alinéa.
\item \textbf{Oublier la transition} : sans connecteur final, les paragraphes se juxtaposent au lieu de s'enchaîner, et la démonstration perd sa fluidité.
\item \textbf{Perdre de vue la problématique} : chaque paragraphe doit visiblement répondre à la question du sujet.
\end{itemize}
\end{attention}

\begin{savaistu}[Et au Baccalauréat technique ?]
Au Probatoire et au Baccalauréat technique, le développement de la dissertation comporte généralement \textbf{deux ou trois grandes parties}, chacune subdivisée en paragraphes construits selon la méthode I-A-E-T. Les correcteurs disposent d'une grille qui valorise explicitement la « cohérence du paragraphe » et l'« enchaînement des idées » : un candidat qui maîtrise le paragraphe argumentatif sécurise donc une part importante de sa note, avant même que soit jugée la richesse de sa culture.
\end{savaistu}

\courssub{TP guidé : du plan au paragraphe rédigé}

\begin{experience}[Transformer un plan de paragraphe en paragraphe rédigé]
\textbf{Matériel} : le roman \emph{Le Vieux Nègre et la Médaille}, un brouillon, la fiche méthode I-A-E-T.
\par\smallskip
\textbf{Protocole} :
\begin{enumerate}
\item Lisez le plan de paragraphe suivant : \emph{Idée : le rire est une arme de dénonciation chez Oyono. Argument : l'ironie ridiculise les colonisateurs et révèle leur injustice. Exemple : la scène de la cérémonie de remise de la médaille, où les discours pompeux contrastent avec la réalité vécue par Meka.}
\item Rédigez d'abord l'idée directrice en une phrase complète, reliée à la problématique.
\item Développez l'argument en trois phrases minimum, en utilisant au moins deux connecteurs de cause ou d'explication (\emph{en effet, car, puisque}).
\item Intégrez l'exemple et commentez-le en une ou deux phrases (\emph{ainsi, on le voit}).
\item Terminez par une transition qui annonce le paragraphe suivant (par exemple sur les conséquences de cette dénonciation pour le lecteur).
\item Relisez en vérifiant : une seule idée ? 8 à 12 lignes ? connecteurs présents ?
\end{enumerate}
\textbf{Observations attendues} : le paragraphe obtenu doit pouvoir être découpé en quatre zones correspondant aux quatre temps I-A-E-T ; si une zone manque, le protocole a été mal suivi et le paragraphe doit être complété.
\par\smallskip
\textbf{Interprétation} : cet exercice montre que la rédaction n'est pas une inspiration mystérieuse mais l'exécution ordonnée d'un plan : qui sait construire un plan de paragraphe sait rédiger un paragraphe.
\end{experience}

\begin{exercice}[À vous de rédiger]
Sujet : \emph{Faut-il que l'écrivain s'engage dans les débats de son époque ?} Rédigez un paragraphe argumentatif complet (8 à 12 lignes) à partir de ce plan : Idée : l'écrivain engagé éveille les consciences. Argument : par son œuvre, il rend visibles des problèmes que l'opinion ignore. Exemple : un fait de société camerounais ou un passage d'\emph{Le Vieux Nègre et la Médaille}. N'oubliez ni l'idée directrice, ni la transition finale.
\end{exercice}

\begin{corrige}[Exercice : rédiger un paragraphe]
Éléments de correction attendus :
\begin{itemize}
\item \textbf{Idée directrice} : une phrase complète annonçant que l'écrivain engagé éveille les consciences, reliée au sujet.
\item \textbf{Argument} : au moins trois phrases d'explication et de justification, avec des connecteurs de cause (\emph{en effet, car, puisque}).
\item \textbf{Exemple} : précis (scène du roman ou fait camerounais identifiable) et suivi d'un commentaire.
\item \textbf{Transition} : une phrase conclusive avec un connecteur (\emph{cependant, ainsi, en outre}) ouvrant sur le paragraphe suivant.
\item \textbf{Forme} : un seul alinéa, 8 à 12 lignes, une seule idée développée.
\end{itemize}
\end{corrige}

\begin{aretenir}[L'essentiel du paragraphe argumentatif]
\begin{itemize}
\item Un paragraphe = \textbf{une seule idée}, matérialisée par un alinéa.
\item Structure en quatre temps (méthode \cle{I-A-E-T}) : \textbf{I}dée directrice, \textbf{A}rgument développé, \textbf{E}xemple précis et commenté, \textbf{T}ransition.
\item Longueur moyenne : \textbf{8 à 12 lignes}.
\item Les connecteurs logiques (conjonctions, adverbes de liaison) assurent la cohérence interne et l'enchaînement des paragraphes.
\item Pièges à éviter : exemples sans argument, paragraphes d'une phrase, plusieurs idées dans un même paragraphe.
\end{itemize}
\end{aretenir}

\coursec{Analyse du sujet, plan détaillé et production complète}

Dans la section précédente, nous avons appris à construire le \cle{paragraphe argumentatif} : une idée directrice, un argument, un exemple, une phrase de liaison. Nous savons donc désormais fabriquer les « briques » de la dissertation. Mais des briques, même excellentes, ne font pas une maison solide si l'on n'a pas d'abord dressé les plans du bâtiment. Cette dernière section est le moment de la synthèse : nous allons apprendre à \cle{analyser un sujet}, à \cle{élaborer un plan détaillé}, puis à \cle{rédiger une dissertation complète} dans le temps de l'épreuve, avant de découvrir comment \cle{évaluer et améliorer} nos productions. C'est l'aboutissement de toute la séquence, et c'est exactement ce qui sera exigé de vous au \cle{Probatoire} et au \cle{Baccalauréat technique}.

\courssub{Analyser le sujet : la première victoire se gagne au brouillon}

\textbf{L'intuition d'abord.} Imaginez un menuisier de Douala à qui un client commande « un meuble ». S'il se met aussitôt à scier des planches sans demander : un meuble pour quoi faire ? pour quelle pièce ? de quelles dimensions ?, il fabriquera peut-être une superbe armoire... alors que le client voulait un lit. Tout son travail sera perdu. En dissertation, le \cle{hors-sujet} est exactement cela : une copie parfois bien écrite, mais qui répond à une question que l'examinateur n'a pas posée. L'analyse du sujet est donc l'étape qui conditionne toutes les autres.

\begin{definition}[L'analyse du sujet]
L'\cle{analyse du sujet} est la première étape du travail de dissertation. Elle consiste à lire attentivement l'énoncé pour : repérer les \cle{mots-clés}, définir chaque terme important, délimiter le sujet dans le \cle{temps} et dans l'\cle{espace}, et dégager le \cle{problème} (la \cle{problématique}) que le sujet pose. Elle se fait entièrement au brouillon et dure environ quinze minutes.
\end{definition}

\begin{methode}[Analyser un sujet en cinq questions]
Face à n'importe quel sujet, posez-vous dans l'ordre les cinq questions suivantes et notez vos réponses au brouillon :
\begin{enumerate}
\item \textbf{De quoi parle-t-on ?} Identifiez le \cle{thème général} du sujet (l'école, la tradition, la justice, le travail...).
\item \textbf{Quels sont les mots-clés ?} Soulignez les termes porteurs de sens et \textbf{définissez chacun} avec vos propres mots. Méfiez-vous des mots qui semblent évidents : ce sont souvent eux qui cachent le problème.
\item \textbf{Quelles sont les limites du sujet ?} Le sujet est-il restreint à une époque (aujourd'hui ? autrefois ?), à un lieu (le Cameroun ? l'Afrique ? le monde ?), à un domaine (l'école ? la famille ?) ? Tout ce qui sort de ces limites est du hors-sujet.
\item \textbf{Quel problème le sujet pose-t-il ?} Transformez le sujet en une véritable question de débat : c'est la \cle{problématique}. Un bon sujet appelle toujours au moins deux réponses opposées.
\item \textbf{Quel type de plan convient ?} Selon la formulation du sujet (question ouverte, affirmation à discuter, citation à commenter), choisissez le plan dialectique, thématique ou analytique.
\end{enumerate}
\end{methode}

\textbf{La recherche d'idées.} Une fois la problématique dégagée, il faut trouver de la matière. Deux techniques complémentaires :
\begin{itemize}
\item le \cle{brainstorming} (remue-méninges) : pendant cinq minutes, notez au brouillon, sans trier et sans juger, tout ce qui vous vient à l'esprit sur le thème : idées, faits de société, proverbes, exemples vécus ;
\item la \cle{mobilisation des lectures} : rappelez-vous les œuvres étudiées en classe. Par exemple, \emph{Le Vieux nègre et la médaille} de Ferdinand Oyono fournit des exemples précieux sur la colonisation, l'illusion, l'hypocrisie ou la dignité : le personnage de Meka, décoré pour avoir « donné » ses terres et ses fils à la France, illustre parfaitement le décalage entre les apparences honorées et la réalité vécue.
\end{itemize}
Ensuite seulement, triez : barrez les idées hors limites, regroupez celles qui se ressemblent, et faites apparaître deux ou trois grandes orientations.

\begin{attention}[Erreur fréquente : le hors-sujet]
Le \cle{hors-sujet} est la faute la plus grave et la plus fréquente en dissertation. Il résulte presque toujours d'une analyse trop rapide des mots-clés : l'élève lit « langues locales » et rédige aussitôt tout ce qu'il sait sur « les langues » en général, ou confond « préserver » avec « enseigner ». Règle d'or : ne commencez \textbf{jamais} à rédiger avant d'avoir défini chaque mot-clé et reformulé le sujet sous forme de question. Dix minutes d'analyse sauvent trois heures de rédaction.
\end{attention}

\courssub{Choisir le type de plan}

Trois types de plans sont admis à l'examen. Le choix dépend de la formulation du sujet.

\begin{definition}[Le plan détaillé]
Le \cle{plan détaillé} est l'organisation hiérarchisée, rédigée au brouillon avant toute rédaction, des \textbf{parties} et \textbf{sous-parties} du devoir, avec pour chacune l'\cle{idée directrice}, les \cle{arguments} et les \cle{exemples} qui seront développés. Il est le squelette complet de la dissertation : celui qui le lit doit pouvoir deviner tout le contenu de la copie.
\end{definition}

\begin{itemize}
\item Le \cle{plan dialectique} (thèse / antithèse / synthèse) convient aux sujets polémiques, aux questions qui appellent un débat : « Faut-il... ? », « Peut-on... ? ». On expose d'abord une position (thèse), puis la position opposée (antithèse), puis on dépasse l'opposition par une réponse nuancée et personnelle (synthèse).
\item Le \cle{plan thématique} convient aux sujets qui invitent à décrire ou à classer les aspects d'une réalité : on examine successivement plusieurs facettes du thème (causes, manifestations, conséquences, par exemple).
\item Le \cle{plan analytique} convient aux sujets qui demandent d'expliquer un phénomène : constat, puis causes, puis conséquences ou solutions.
\end{itemize}

\begin{popfigure}
\begin{center}
\begin{tikzpicture}[font=\small,
boite/.style={draw=popblue, thick, fill=popblueL, rounded corners=2pt, text width=2.6cm, align=center, minimum height=0.8cm},
titre/.style={font=\bfseries\small, text=popdark}]
% Plan dialectique
\node[titre] at (0,3.4) {Plan dialectique};
\node[boite] (t) at (0,2.4) {Thèse : une position};
\node[boite] (a) at (0,1.3) {Antithèse : position opposée};
\node[boite] (s) at (0,0.2) {Synthèse : dépassement};
\draw[-{Stealth}, thick, popblue] (t) -- (a);
\draw[-{Stealth}, thick, popblue] (a) -- (s);
% Plan thématique
\node[titre] at (5.2,3.4) {Plan thématique};
\node[boite, draw=popgreen, fill=popgreenL] (th) at (5.2,2.4) {Thème};
\node[boite, draw=popgreen, fill=popgreenL] (a1) at (3.4,0.9) {Aspect 1};
\node[boite, draw=popgreen, fill=popgreenL] (a2) at (5.2,0.9) {Aspect 2};
\node[boite, draw=popgreen, fill=popgreenL] (a3) at (7.0,0.9) {Aspect 3};
\draw[-{Stealth}, thick, popgreen] (th) -- (a1);
\draw[-{Stealth}, thick, popgreen] (th) -- (a2);
\draw[-{Stealth}, thick, popgreen] (th) -- (a3);
% Plan analytique
\node[titre] at (10.4,3.4) {Plan analytique};
\node[boite, draw=poporange, fill=poporangeL] (c) at (10.4,2.4) {Constat};
\node[boite, draw=poporange, fill=poporangeL] (ca) at (10.4,1.3) {Causes};
\node[boite, draw=poporange, fill=poporangeL] (co) at (10.4,0.2) {Conséquences / solutions};
\draw[-{Stealth}, thick, poporange] (c) -- (ca);
\draw[-{Stealth}, thick, poporange] (ca) -- (co);
\end{tikzpicture}
\end{center}
\legende{Les trois types de plans admis à l'examen : le plan dialectique oppose deux positions avant de les dépasser ; le plan thématique déploie les aspects d'un thème ; le plan analytique va du constat aux causes puis aux conséquences.}
\end{popfigure}

\begin{exemplebox}[Plan dialectique détaillé : « Faut-il préserver les langues locales à l'école camerounaise ? »]
\textbf{Analyse.} Mots-clés : \emph{préserver} (garder de la disparition), \emph{langues locales} (langues camerounaises : ewondo, duala, fulfuldé, etc.), \emph{école} (limitation : le cadre scolaire, au Cameroun, aujourd'hui). Problématique : l'école camerounaise doit-elle faire une place aux langues locales, ou doit-elle se consacrer aux langues officielles ?\par
\textbf{I. Thèse : il faut préserver les langues locales à l'école.}
\begin{itemize}
\item 1. Elles sont un patrimoine culturel menacé de disparition (ex. : des langues camerounaises ne sont plus transmises aux enfants des villes).
\item 2. L'enfant apprend mieux dans une langue qu'il maîtrise (ex. : écolier du village qui comprend les explications données dans sa langue).
\end{itemize}
\textbf{II. Antithèse : l'école doit privilégier le français et l'anglais.}
\begin{itemize}
\item 1. Les langues officielles unifient un pays de plus de 250 langues (ex. : le français comme langue commune à l'administration).
\item 2. Elles ouvrent sur le monde, les études supérieures et l'emploi (ex. : concours, manuels scientifiques rédigés en français ou en anglais).
\end{itemize}
\textbf{III. Synthèse : concilier les deux exigences.}
\begin{itemize}
\item 1. Les langues locales et officielles ne sont pas ennemies : le bilinguisme est une richesse.
\item 2. Solution : introduire les langues locales au primaire comme aide à l'apprentissage, tout en consolidant les langues officielles (ex. : expériences d'écoles bilingues au Cameroun).
\end{itemize}
\end{exemplebox}

\courssub{Rédiger la dissertation complète en temps limité}

Le jour de l'épreuve, le temps est votre adversaire le plus redoutable. Une dissertation se gère comme un match : chaque mi-temps a sa fonction.

\begin{methode}[Gérer le temps de l'épreuve]
Pour une épreuve de deux heures, répartissez le temps ainsi :
\begin{enumerate}
\item \textbf{Analyse du sujet et recherche d'idées : 15 minutes} (mots-clés, limites, problématique, brainstorming) ;
\item \textbf{Élaboration du plan détaillé : 20 minutes} (parties, sous-parties, arguments et exemples au brouillon) ;
\item \textbf{Rédaction : environ 1 h 15} (introduction, développement, conclusion, en rédigeant directement au propre avec soin) ;
\item \textbf{Relecture : 10 minutes} (orthographe, accords, connecteurs, présentation).
\end{enumerate}
Ne sacrifiez jamais les deux premières étapes : une copie bien planifiée se rédige vite et bien.
\end{methode}

\begin{popfigure}
\begin{center}
\begin{tikzpicture}[font=\small]
% Camembert : 15 + 20 + 75 + 10 = 120 min
% angles : 15->45°, 20->60°, 75->225°, 10->30° (départ 90°)
\fill[popblueL, draw=popblue, thick] (0,0) -- (90:2.2) arc (90:45:2.2) -- cycle;
\fill[popgreenL, draw=popgreen, thick] (0,0) -- (45:2.2) arc (45:-15:2.2) -- cycle;
\fill[poporangeL, draw=poporange, thick] (0,0) -- (-15:2.2) arc (-15:-240:2.2) -- cycle;
\fill[poppinkL, draw=poppink, thick] (0,0) -- (-240:2.2) arc (-240:-270:2.2) -- cycle;
\node[font=\footnotesize\bfseries, text=popdark] at (67:1.4) {15 min};
\node[font=\footnotesize\bfseries, text=popdark] at (15:1.4) {20 min};
\node[font=\footnotesize\bfseries, text=popdark] at (-128:1.4) {75 min};
\node[font=\footnotesize\bfseries, text=popdark] at (-255:1.5) {10 min};
\node[right, font=\footnotesize] at (2.7,1.2) {\textcolor{popblue}{\rule{0.35cm}{0.18cm}} Analyse du sujet};
\node[right, font=\footnotesize] at (2.7,0.6) {\textcolor{popgreen}{\rule{0.35cm}{0.18cm}} Plan détaillé};
\node[right, font=\footnotesize] at (2.7,0.0) {\textcolor{poporange}{\rule{0.35cm}{0.18cm}} Rédaction};
\node[right, font=\footnotesize] at (2.7,-0.6) {\textcolor{poppink}{\rule{0.35cm}{0.18cm}} Relecture};
\end{tikzpicture}
\end{center}
\legende{Répartition conseillée du temps pour une épreuve de dissertation de deux heures (120 minutes) : l'analyse et le plan représentent près d'un tiers du temps, la relecture est obligatoire.}
\end{popfigure}

\textbf{De la rédaction à la copie.} Avec le plan détaillé sous les yeux, la rédaction consiste à transformer chaque sous-partie en un ou deux paragraphes construits selon la méthode de la section précédente (idée directrice, argument, exemple, transition), à rédiger l'introduction (accroche, présentation du sujet, problématique, annonce du plan) et la conclusion (bilan, réponse à la problématique, ouverture). Les \cle{connecteurs logiques} étudiés dans cette leçon — \emph{d'abord, ensuite, en effet, cependant, par conséquent, enfin} — assurent la liaison entre les phrases et entre les parties.

\begin{attention}[Erreur fréquente : négliger la relecture]
Beaucoup de copies perdent plusieurs points non pas à cause d'idées faibles, mais à cause de fautes que dix minutes de relecture auraient effacées : accords sujet-verbe oubliés, homophones confondus (\emph{a/à}, \emph{ou/où}, \emph{ce/se}, \emph{son/sont}), connecteurs répétitifs ou absents, mots oubliés. Prévoyez toujours ces dix minutes finales : relisez phrase par phrase, en vérifiant d'abord les accords, puis les homophones, puis la ponctuation.
\end{attention}

\courssub{Confrontation et amélioration des productions}

Une fois la copie rédigée, le travail n'est pas fini : en classe, la \cle{confrontation des productions} permet de progresser en s'évaluant soi-même (auto-évaluation) puis en évaluant la copie d'un camarade (évaluation croisée), à l'aide d'une grille commune.

\begin{center}
\begin{tabularx}{\linewidth}{|l|X|c|}
\hline
\rowcolor{popblueL} \textbf{Critère} & \textbf{Questions à se poser} & \textbf{Barème} \\
\hline
Respect du sujet & Les mots-clés sont-ils définis ? La problématique répond-elle au sujet posé ? Aucune partie ne sort-elle des limites ? & /4 \\
\hline
Cohérence du plan & Le type de plan est-il adapté ? Les parties s'enchaînent-elles logiquement ? Les transitions existent-elles ? & /4 \\
\hline
Qualité des arguments & Chaque paragraphe contient-il idée directrice, argument et exemple ? Les exemples sont-ils précis ? & /4 \\
\hline
Langue et connecteurs & Orthographe, accords, homophones, ponctuation ; variété des connecteurs logiques. & /4 \\
\hline
Présentation & Écriture lisible, paragraphes visibles, introduction et conclusion nettement séparées. & /4 \\
\hline
\end{tabularx}
\end{center}

\begin{exoresolu}[Correction collective d'une copie-type]
\textbf{Énoncé.} Voici l'introduction d'un élève au sujet « Faut-il préserver les langues locales à l'école camerounaise ? » : \emph{« Les langues sont très importantes dans la vie. Au Cameroun il y a beaucoup de langues. Je vais vous parler des langues. »} Appliquez la grille d'évaluation et proposez une version améliorée.
\tcblower
\textbf{Solution détaillée.} Analyse selon les critères :
\begin{itemize}
\item \textbf{Respect du sujet} : faible. Aucun mot-clé n'est défini ; « les langues » en général élargit le sujet au lieu de le délimiter (langues \emph{locales}, à \emph{l'école}, au \emph{Cameroun}).
\item \textbf{Problématique} : absente. Aucune question n'est posée, donc aucun débat n'est possible.
\item \textbf{Annonce du plan} : absente. « Je vais vous parler des langues » n'annonce ni parties ni démarche, et le « je vous parle » relève de l'oral, pas de la dissertation.
\item \textbf{Langue} : correcte mais pauvre et répétitive (« langues » trois fois).
\end{itemize}
\textbf{Version améliorée} : \emph{« Le Cameroun compte plus de deux cent cinquante langues locales, véritables trésors de son patrimoine culturel. Pourtant, à l'école, seuls le français et l'anglais occupent une place officielle. Faut-il alors préserver les langues locales en les introduisant dans l'enseignement, ou l'école doit-elle se consacrer aux seules langues officielles ? Nous verrons d'abord que les langues locales méritent d'être préservées à l'école, puis que les langues officielles présentent des avantages décisifs, avant de montrer qu'il est possible de concilier les deux exigences. »} Cette version définit le cadre, pose la problématique sous forme de question et annonce clairement un plan dialectique en trois parties.
\end{exoresolu}

\textbf{Compte rendu et remédiation.} Après l'évaluation croisée, le professeur fait le \cle{compte rendu} collectif : il corrige au tableau une copie-type, recense les erreurs fréquentes de la classe (hors-sujet, paragraphes sans exemple, homophones, absence de transitions) et propose des exercices de \cle{remédiation} ciblés : refaire une introduction, transformer une liste d'idées en plan détaillé, réécrire un paragraphe en insérant des connecteurs. La remédiation n'est pas une punition : c'est un deuxième entraînement sur le point précis qui a fait défaut.

\begin{savaistu}[Comment la dissertation est-elle notée au Baccalauréat technique ?]
Au \cle{Baccalauréat technique}, la dissertation est évaluée selon quatre grandes exigences : la \textbf{pertinence du contenu} (respect du sujet, richesse des idées et des exemples), la \textbf{cohérence du plan} (organisation logique, transitions), la \textbf{qualité de la langue} (orthographe, syntaxe, vocabulaire, connecteurs) et la \textbf{présentation} de la copie (lisibilité, aération des paragraphes). Autrement dit, un candidat aux idées moyennes mais rigoureux dans la méthode obtient souvent une meilleure note qu'un candidat brillant mais désorganisé : la méthode que vous venez d'apprendre est votre meilleure alliée.
\end{savaistu}

\begin{exercice}[Entraînement complet]
Sujet : \emph{« Le travail scolaire est-il la seule voie de la réussite ? »}
\begin{enumerate}
\item Analysez le sujet en répondant aux cinq questions de la méthode (thème, mots-clés définis, limites, problématique, type de plan).
\item Rédigez au brouillon un plan détaillé en trois parties, avec arguments et exemples pour chaque sous-partie.
\item Rédigez l'introduction complète, puis chronométrez-vous : vous devez avoir terminé ces trois étapes en 45 minutes.
\end{enumerate}
\end{exercice}

\begin{corrige}[Exercice : éléments de correction]
\textbf{1. Analyse.} Thème : la réussite et ses conditions. Mots-clés : \emph{travail scolaire} (effort fourni dans les études), \emph{seule voie} (caractère exclusif — c'est le mot qui porte le débat), \emph{réussite} (succès dans la vie, pas seulement aux examens). Limites : aucune imposée, on peut raisonner dans le contexte camerounais actuel. Problématique : le travail scolaire suffit-il à garantir la réussite, ou d'autres voies existent-elles ? Plan adapté : dialectique.\par
\textbf{2. Plan détaillé (exemple).} I. Le travail scolaire est une voie essentielle de la réussite : 1. il donne des diplômes exigés par les emplois (ex. : concours de la fonction publique) ; 2. il forme l'esprit et la persévérance. II. Mais il n'est pas la seule voie : 1. la formation technique et professionnelle mène à des métiers recherchés (ex. : techniciens, artisans prospères) ; 2. l'initiative personnelle et l'entrepreneuriat réussissent parfois sans longs parcours scolaires (ex. : jeunes entrepreneurs camerounais). III. Synthèse : la réussite naît de l'alliance du travail, scolaire ou non, avec la persévérance et le sens de l'effort ; l'essentiel est de travailler, quel que soit le chemin choisi.\par
\textbf{3. Introduction attendue} : accroche sur la réussite comme aspiration universelle, présentation du sujet, problématique sous forme de question, annonce des trois parties du plan dialectique.
\end{corrige}

\begin{aretenir}[L'essentiel de la section]
\begin{itemize}
\item L'\cle{analyse du sujet} (mots-clés définis, limites, problématique) est l'étape qui protège du hors-sujet : ne jamais rédiger avant de l'avoir achevée.
\item Le \cle{plan détaillé} organise au brouillon parties, sous-parties, arguments et exemples ; on choisit le plan \cle{dialectique}, \cle{thématique} ou \cle{analytique} selon la formulation du sujet.
\item Le temps de l'épreuve se gère : 15 min d'analyse, 20 min de plan, environ 1 h 15 de rédaction, 10 min de relecture obligatoire.
\item La \cle{grille d'évaluation} (respect du sujet, cohérence, arguments, langue, présentation) permet l'auto-correction, l'évaluation croisée et la remédiation.
\end{itemize}
\end{aretenir}

\newpage

\coursec{Activité d'intégration}

\courssub{Objectifs de l'activité}

À l'issue de cet atelier d'écriture, tu seras capable de :

\begin{itemize}
\item analyser un sujet de dissertation (dégager le thème, la thèse implicite, la problématique) ;
\item élaborer un plan détaillé cohérent (thèse, antithèse, synthèse) au brouillon ;
\item rédiger une introduction complète (amorce, présentation du sujet, problématique, annonce du plan) et une conclusion (bilan, ouverture) ;
\item construire un paragraphe argumentatif respectant la triade \cle{idée directrice} -- \cle{argument} -- \cle{exemple} ;
\item employer des outils de liaison variés (conjonctions de coordination et de subordination, adverbes de liaison, locutions adverbiales, pronoms relatifs) ;
\item évaluer la production d'un pair à l'aide d'une grille, puis améliorer ta propre production (remédiation).
\end{itemize}

\courssub{Situation}

Le lycée technique de Nkongsamba organise, comme chaque année, sa semaine du livre et de la francophonie. Cette année, le thème retenu est : \emph{« Honneur et dignité : que vaut une médaille ? »}. Le club littéraire, dont tu es membre, a été chargé de produire un texte de réflexion qui sera lu lors de la cérémonie de clôture, en présence du proviseur, des parents d'élèves et d'un représentant de la délégation départementale des enseignements secondaires.

Le débat est d'actualité : au Cameroun, les distinctions honorifiques (médailles du travail, ordres nationaux, titres traditionnels) font régulièrement l'objet de discussions dans la presse et sur les réseaux sociaux. Certains y voient la juste récompense du mérite ; d'autres dénoncent des honneurs distribués sans égard à la valeur réelle des récipiendaires. La lecture du roman \emph{Le Vieux Nègre et la médaille} de Ferdinand Oyono, étudiée en classe, nourrit cette réflexion : Meka, le vieux nègre, reçoit une médaille pour avoir « donné » ses terres à la mission et ses fils à la guerre, mais cet honneur ne lui apporte ni respect véritable ni bien-être.

Le sujet proposé par le jury du concours d'écriture est le suivant :

\begin{center}
\textbf{« L'honneur reçu des autres suffit-il à donner sa dignité à un homme ? »}
\end{center}

Pour t'aider, le jury met à ta disposition un dossier documentaire de quatre documents :

\begin{itemize}
\item \textbf{Document 1} (extrait du roman) : Meka, après la remise de la médaille, constate que sa vie n'a pas changé ; il reste le même vieil homme pauvre, moqué par certains, oublié par l'administration dès la fête terminée.
\item \textbf{Document 2} (extrait du roman) : avant la cérémonie, les villageois envient Meka ; après, ils rient de lui, car la médaille n'a rempli ni sa case ni sa marmite.
\item \textbf{Document 3} (article de presse camerounaise) : un éditorialiste s'interroge sur la multiplication des distinctions honorifiques et rappelle que « la dignité se construit par les actes, non par les rubans ».
\item \textbf{Document 4} (citation attribuée à Ferdinand Oyono) : « On ne donne pas la dignité à un homme comme on épingle une médaille sur sa poitrine. »
\end{itemize}

Le jury impose des contraintes précises : la dissertation comportera une introduction d'environ 10 lignes, trois grandes parties et une conclusion ; chaque paragraphe devra utiliser au moins \textbf{cinq connecteurs variés} ; le travail se fait en binôme et sera évalué sur une grille de 20 points répartis ainsi : pertinence des arguments (6 points), cohérence du plan (4 points), qualité des liaisons logiques (5 points), correction de la langue (5 points).

\courssub{Consignes}

\begin{enumerate}
\item \textbf{Analyse du sujet.} Souligne les mots-clés du sujet (« honneur », « reçu des autres », « suffit », « dignité »). Définis chacun d'eux en une phrase, puis reformule le sujet sous forme d'une problématique claire.
\item \textbf{Mobilisation du dossier.} Pour chacun des quatre documents, indique en deux lignes s'il soutient plutôt la thèse (l'honneur suffit), l'antithèse (l'honneur ne suffit pas) ou la synthèse (l'honneur a une valeur relative).
\item \textbf{Plan détaillé.} Au brouillon, élabore un plan dialectique en trois parties. Pour chaque partie, note l'idée directrice, deux arguments et un exemple tiré du roman, du dossier ou de la vie sociale camerounaise.
\item \textbf{Rédaction de l'introduction.} Rédige l'introduction complète : amorce, présentation du sujet, problématique, annonce du plan.
\item \textbf{Rédaction d'un paragraphe type.} Rédige le paragraphe de développement de la deuxième partie (antithèse), en veillant à la chaîne idée directrice -- argument -- exemple -- analyse, et en utilisant au moins cinq connecteurs variés (par exemple : \emph{en effet, cependant, parce que, qui, c'est pourquoi, néanmoins}).
\item \textbf{Évaluation croisée.} Échange ta production avec celle d'un autre binôme. Évalue-la à l'aide de la grille du jury et note deux points forts et deux points à améliorer.
\item \textbf{Remédiation.} Réécris ton paragraphe en tenant compte des remarques reçues, puis participe au compte rendu collectif au tableau.
\end{enumerate}

\courssub{Ressources à mobiliser}

\begin{aretenir}[Outils de la dissertation]
\begin{itemize}
\item \textbf{Structure de l'introduction} : amorce (phrase d'accroche) ; présentation du sujet ; problématique ; annonce du plan.
\item \textbf{Structure du paragraphe} : idée directrice ; argument ; exemple précis ; analyse de l'exemple ; phrase de transition.
\item \textbf{Structure de la conclusion} : bilan de la réflexion (réponse à la problématique) ; ouverture.
\item \textbf{Outils de liaison} : conjonctions de coordination (\emph{mais, ou, et, donc, or, ni, car}) ; conjonctions de subordination (\emph{parce que, puisque, bien que, si, lorsque}) ; adverbes et locutions de liaison (\emph{d'abord, ensuite, en effet, cependant, néanmoins, en outre, enfin, c'est pourquoi}) ; pronoms relatifs (\emph{qui, que, dont, où}).
\item \textbf{Stratégies argumentatives} : argument d'autorité (citation d'Oyono), argument d'expérience (cas de Meka), argument de valeur (dignité, justice), argument logique (cause-conséquence).
\end{itemize}
\end{aretenir}

\begin{attention}[Pièges fréquents]
Évite le hors-sujet (parler du roman sans répondre à la question), le résumé de l'œuvre à la place de l'argumentation, les connecteurs répétitifs (aligner cinq fois « ensuite »), et l'annonce du plan absente de l'introduction. Une citation doit toujours être introduite et commentée.
\end{attention}

\courssub{Démarche et corrigé commenté}

\begin{exoresolu}[Atelier d'écriture : de l'analyse du sujet à la remédiation]
Énoncé. À partir du sujet « L'honneur reçu des autres suffit-il à donner sa dignité à un homme ? » et du dossier documentaire, produire en binôme l'analyse du sujet, un plan détaillé, une introduction et un paragraphe type, puis procéder à l'évaluation croisée et à la réécriture.

\tcblower

\textbf{Étape 1 : analyse du sujet.} Les mots-clés sont : \emph{honneur reçu des autres} (reconnaissance extérieure : médaille, titre, distinction), \emph{suffit} (est-ce la condition unique et pleine ?), \emph{dignité} (valeur intrinsèque, respect de soi et d'autrui). Le sujet oppose donc une valeur extérieure (l'honneur) à une valeur intérieure (la dignité). Problématique possible : \emph{la reconnaissance accordée par autrui peut-elle, à elle seule, fonder la dignité d'un homme, ou celle-ci naît-elle d'abord de ses actes et de sa conscience ?}

\textbf{Étape 2 : exploitation du dossier.} Documents 1 et 2 : antithèse (la médaille de Meka ne change ni sa condition ni le regard réel des autres). Document 3 : antithèse et synthèse (la dignité se construit par les actes). Document 4 : antithèse (argument d'autorité). La thèse de départ peut s'appuyer sur l'expérience sociale : les distinctions motivent, récompensent le mérite, donnent une place dans la communauté.

\textbf{Étape 3 : plan détaillé (dialectique).}
\begin{itemize}
\item \textbf{I. L'honneur reçu semble conférer la dignité.} Argument 1 : la distinction reconnaît publiquement le mérite (exemple : médailles du travail décernées le 1\textsuperscript{er} mai au Cameroun). Argument 2 : elle donne un statut et une estime sociale (exemple : le respect dont jouissent les notables titrés).
\item \textbf{II. Mais l'honneur extérieur ne suffit pas.} Argument 1 : il peut être immérité ou hypocrite (exemple : Meka, médaille « récompensant » sa spoliation). Argument 2 : il ne change pas la condition réelle (exemple : Meka reste pauvre et moqué, documents 1 et 2 ; citation d'Oyono, document 4).
\item \textbf{III. La dignité naît des actes et de la conscience de soi.} Argument 1 : elle se construit par le travail, l'honnêteté, le courage (exemple : l'éditorialiste du document 3). Argument 2 : l'honneur n'a de valeur que s'il couronne une dignité déjà acquise (exemple : le sportif ou l'enseignant méritant).
\end{itemize}

\textbf{Étape 4 : introduction rédigée (modèle).} « Dans nos sociétés, les cérémonies de décoration sont des moments forts où la communauté honore ceux qu'elle juge méritants. Au Cameroun comme ailleurs, médailles et titres sont censés couronner une vie de service. C'est toute l'ambiguïté de cette reconnaissance que Ferdinand Oyono met en scène dans \emph{Le Vieux Nègre et la médaille}, où Meka reçoit une médaille qui ne lui apporte ni respect ni bonheur. On peut alors se demander : l'honneur reçu des autres suffit-il à donner sa dignité à un homme ? Nous verrons d'abord que l'honneur semble conférer la dignité, puis qu'il ne suffit pas à la fonder, enfin que la dignité véritable naît des actes et de la conscience de soi. »

\textbf{Étape 5 : paragraphe type rédigé (partie II, modèle).} « Cependant, l'honneur extérieur ne suffit pas à donner sa dignité à un homme. En effet, une distinction peut être imméritée ou servir des intérêts qui humilient celui qu'elle prétend honorer. Meka, \emph{qui} a perdu ses terres et ses fils, reçoit une médaille \emph{parce que} l'administration coloniale veut se donner bonne conscience ; \emph{or} cette médaille ne remplit ni sa case ni sa marmite, \emph{et} les villageois, qui l'enviaient la veille, rient de lui le lendemain. Oyono le dit clairement : on ne donne pas la dignité « comme on épingle une médaille sur sa poitrine ». \emph{C'est pourquoi} l'honneur reçu, lorsqu'il ne s'accompagne d'aucune justice réelle, demeure une décoration vide. » On compte bien plus de cinq connecteurs variés : \emph{cependant, en effet, qui, parce que, or, et, c'est pourquoi}.

\textbf{Étape 6 : évaluation croisée.} Le binôme évaluateur applique la grille : pertinence (/6), cohérence (/4), liaisons (/5), langue (/5). Exemple de retour : « Arguments pertinents et exemple bien analysé (5/6) ; transition vers la troisième partie absente (3/4) ; connecteurs variés (5/5) ; deux accords à revoir (4/5). »

\textbf{Étape 7 : remédiation.} La réécriture ajoute une phrase de transition (« Ainsi, si l'honneur ne suffit pas, c'est que la dignité se cherche ailleurs : dans les actes. ») et corrige les accords signalés. Le compte rendu au tableau permet de mutualiser les problématiques et les meilleurs connecteurs repérés dans les copies.
\end{exoresolu}

\courssub{Bilan de l'activité}

Cet atelier a permis de mobiliser et de vérifier l'ensemble des savoir-faire de la leçon. L'analyse du sujet a montré qu'une problématique naît de la tension entre deux mots-clés (ici, honneur extérieur et dignité intérieure). L'élaboration du plan détaillé a confirmé qu'un bon plan dialectique réserve la synthèse au dépassement des deux premières positions. La rédaction de l'introduction et du paragraphe a mis en pratique la triade idée directrice -- argument -- exemple et l'usage varié des outils de liaison, sans lesquels le texte perd sa cohérence. Enfin, l'évaluation croisée et la remédiation ont fait apparaître que réécrire est une étape indispensable de l'écriture : c'est en corrigeant les transitions, les accords et la précision des exemples qu'une copie gagne en pertinence et en clarté, compétence directement exigible à l'épreuve de dissertation du Baccalauréat technique.

\newpage

\coursec{Méthodes et exercices résolus}

\coursec{Leçon 08 : Dissertation — Rédiger une dissertation}

\courssub{Méthode 1 : Analyser un sujet de dissertation}

\begin{methode}[Analyser le sujet en 5 étapes]
\begin{enumerate}
\item \textbf{Lire le sujet plusieurs fois} et souligner les \cle{mots-clés} (notions principales, verbe de consigne : « discutez », « commentez », « pensez-vous que... »).
\item \textbf{Définir chaque mot-clé} : donner un sens précis aux termes importants pour éviter le hors-sujet.
\item \textbf{Repérer la tension} du sujet : la plupart des sujets opposent deux idées (par exemple : la médaille honore-t-elle ou humilie-t-elle le vieux Meka ?).
\item \textbf{Transformer le sujet en problématique} : formuler une question précise à laquelle toute la dissertation répondra.
\item \textbf{Mobiliser ses connaissances} : noter au brouillon les arguments, les exemples (œuvres lues, faits de société, proverbes) et choisir le type de plan adapté (dialectique : thèse/antithèse/synthèse ; thématique ; analytique).
\end{enumerate}
\textbf{Réflexe} : si le sujet contient une citation ou une affirmation à discuter, le plan dialectique est presque toujours le plus sûr.
\end{methode}

\begin{exoresolu}[Analyse d'un sujet type Bac]
\textbf{Énoncé.} Sujet : « La distinction sociale est-elle toujours source de bonheur ? Vous discuterez cette affirmation en vous appuyant sur \emph{Le Vieux nègre et la médaille} de Ferdinand Oyono et sur votre expérience personnelle. » Dégagez les mots-clés, la problématique et le type de plan adapté.
\tcblower
\textbf{Solution détaillée.}
\begin{enumerate}
\item \textbf{Mots-clés} : « distinction sociale » (honneur, médaille, titre, prestige conféré par la société) ; « toujours » (terme absolu, à discuter) ; « bonheur » (bien-être, satisfaction profonde).
\item \textbf{Définitions} : la distinction sociale est une marque de reconnaissance accordée par un groupe ou une autorité ; le bonheur est un état durable de satisfaction.
\item \textbf{Tension du sujet} : d'un côté, la distinction flatte et valorise (Meka est d'abord fier de sa médaille) ; de l'autre, elle peut humilier, isoler ou tromper (la médaille de Meka est un instrument de propagande coloniale, et sa nuit de honte au poste révèle la vanité de cet honneur).
\item \textbf{Problématique} : la reconnaissance sociale apporte-t-elle un bonheur authentique, ou n'est-elle qu'un leurre qui masque une domination ?
\item \textbf{Plan choisi} : plan \cle{dialectique}. Thèse : la distinction sociale procure fierté et considération (ex. : la joie de Meka, l'envie des villageois). Antithèse : elle peut être un instrument de manipulation et une source d'humiliation (ex. : la médaille comme récompense de la soumission, l'arrestation de Meka). Synthèse : seule une reconnaissance fondée sur la dignité réelle, et non sur l'apparence, rend heureux.
\end{enumerate}
\textbf{Conclusion} : le sujet appelle une discussion nuancée ; le mot « toujours » invite à dépasser la réponse simple par un plan dialectique.
\end{exoresolu}

\courssub{Méthode 2 : Rédiger l'introduction}

\begin{methode}[Les 4 temps de l'introduction]
\begin{enumerate}
\item \textbf{L'amorce (accroche)} : une phrase générale liée au thème (réalité sociale, constat, proverbe) — jamais une définition de dictionnaire banale.
\item \textbf{L'amenée} : on resserre progressivement vers le sujet et on cite le sujet ou l'œuvre de référence.
\item \textbf{La problématique} : la question centrale, clairement formulée (une phrase interrogative).
\item \textbf{L'annonce du plan} : « Nous verrons d'abord que..., ensuite que..., enfin que... » (deux ou trois parties, sans détailler les arguments).
\end{enumerate}
\textbf{Longueur} : 8 à 12 lignes. \textbf{Interdit} : donner son avis personnel dans l'introduction.
\end{methode}

\begin{exoresolu}[Rédaction d'une introduction complète]
\textbf{Énoncé.} Rédigez l'introduction de la dissertation sur le sujet : « La distinction sociale est-elle toujours source de bonheur ? » (appui sur \emph{Le Vieux nègre et la médaille}).
\tcblower
\textbf{Solution détaillée} (modèle rédigé et commenté).

\emph{Amorce} : « Dans toutes les sociétés, les hommes recherchent la reconnaissance de leurs semblables : titres, décorations et honneurs semblent couronner une vie réussie. »

\emph{Amenée} : « C'est précisément cette quête de considération que Ferdinand Oyono met en scène dans \emph{Le Vieux nègre et la médaille} : le vieux Meka reçoit une médaille des mains de l'administration coloniale pour avoir donné ses terres et ses fils à la France. Mais cet honneur se solde par une nuit d'humiliation au poste de police. »

\emph{Problématique} : « Dès lors, une question s'impose : la distinction sociale est-elle toujours source de bonheur, ou peut-elle devenir un instrument de domination et de désillusion ? »

\emph{Annonce du plan} : « Nous montrerons d'abord que la distinction sociale procure fierté et prestige ; nous verrons ensuite qu'elle peut masquer une manipulation et engendrer la honte ; nous soutiendrons enfin que seule une reconnaissance fondée sur la dignité procure un bonheur authentique. »

\textbf{Vérification} : les 4 temps sont présents, dans l'ordre ; la problématique reprend les mots-clés du sujet ; le plan annonce trois parties sans révéler les arguments.
\end{exoresolu}

\courssub{Méthode 3 : Construire un paragraphe argumentatif}

\begin{methode}[La règle I-A-E-E]
Un paragraphe de dissertation suit quatre étapes obligatoires :
\begin{enumerate}
\item \textbf{Idée directrice} : une phrase qui annonce l'argument du paragraphe (jamais d'exemple en premier).
\item \textbf{Argument} : on développe et on justifie l'idée par un raisonnement logique.
\item \textbf{Exemple} : une référence précise (scène du roman, fait historique, fait de société camerounais ou africain).
\item \textbf{Explication / transition} : on montre que l'exemple prouve bien l'idée, puis on annonce le paragraphe suivant.
\end{enumerate}
\textbf{Réflexes} : un paragraphe = une seule idée ; utiliser des connecteurs logiques (\cle{en effet}, \cle{par conséquent}, \cle{cependant}) ; 8 à 15 lignes par paragraphe.
\end{methode}

\begin{exoresolu}[Rédaction d'un paragraphe argumentatif]
\textbf{Énoncé.} Rédigez un paragraphe complet soutenant l'idée suivante : « La médaille de Meka est un instrument de manipulation coloniale. »
\tcblower
\textbf{Solution détaillée.}

\emph{Idée directrice} : « La médaille décernée à Meka, loin d'être une récompense sincère, apparaît comme un instrument de manipulation au service de la propagande coloniale. »

\emph{Argument} : « En effet, l'administration coloniale n'honore pas Meka pour sa valeur personnelle : elle célèbre sa soumission. Récompenser un vieillard qui a donné ses terres à la mission et perdu ses deux fils à la guerre, c'est envoyer à tous les villageois le message que l'obéissance paie. La distinction devient ainsi un moyen de fidéliser les populations et de masquer la spoliation. »

\emph{Exemple} : « La scène de la cérémonie du 14 juillet l'illustre parfaitement : le commandant accroche la médaille sur la poitrine de Meka comme on décore un objet, tandis que les notables blancs se congratulent ; le vieil homme, ému, ne comprend même pas la portée politique de la fête organisée autour de lui. »

\emph{Explication et transition} : « Cet épisode prouve que la distinction n'honore pas l'homme mais instrumentalise sa naïveté. Cependant, la médaille ne trompe pas seulement Meka : elle finit par se retourner contre lui, comme le montre la suite du récit. »

\textbf{Vérification} : les quatre temps I-A-E-E sont présents ; l'exemple est précis et rattaché à l'argument par « cet épisode prouve que... ».
\end{exoresolu}

\courssub{Méthode 4 : Employer les stratégies argumentatives}

\begin{methode}[Choisir sa stratégie argumentative]
\begin{enumerate}
\item \textbf{Identifier la stratégie de l'énonciateur} : argumentation directe (essai, discours : l'auteur assume son point de vue avec « je ») ou indirecte (conte, roman, satire : le message passe par une fiction, comme chez Oyono).
\item \textbf{Repérer les procédés} : pour convaincre (raisonnement logique, preuves, exemples) ; pour persuader (émotions, images, ironie, hyperbole, registre pathétique ou polémique).
\item \textbf{Adapter sa propre stratégie au sujet} : dans une dissertation, privilégier l'argumentation directe et logique ; réserver l'émotion à quelques exemples bien choisis.
\item \textbf{Hiérarchiser les arguments} : du plus faible au plus fort dans chaque partie, pour terminer en force.
\item \textbf{Anticiper les objections} : « On objectera que..., mais... » renforce la démonstration.
\end{enumerate}
\end{methode}

\begin{exoresolu}[Analyse d'une stratégie argumentative]
\textbf{Énoncé.} Dans \emph{Le Vieux nègre et la médaille}, Oyono dénonce la colonisation sans jamais écrire « la colonisation est injuste ». Identifiez la stratégie employée et montrez son efficacité à travers deux procédés.
\tcblower
\textbf{Solution détaillée.}
\begin{enumerate}
\item \textbf{Stratégie} : Oyono choisit l'argumentation \cle{indirecte} : c'est le récit romanesque, et non un discours, qui porte la dénonciation. Le lecteur conclut lui-même à l'injustice coloniale, ce qui rend la critique plus percutante.
\item \textbf{Procédé 1 — l'ironie satirique} : la médaille, symbole supposé de gloire, récompense en réalité la perte des terres et des fils de Meka. Ce contraste entre l'apparence (honneur) et la réalité (spoliation) crée une ironie amère qui ridiculise le discours colonial.
\item \textbf{Procédé 2 — le registre pathétique} : la nuit passée par Meka au poste, humilié, mouillé, traité comme un délinquant alors qu'il venait d'être décoré, suscite la pitié et l'indignation du lecteur.
\end{enumerate}
\textbf{Conclusion} : par l'ironie et le pathétique, l'argumentation indirecte d'Oyono fait ressentir l'hypocrisie coloniale au lieu de la décrire : la fiction devient ainsi une arme de dénonciation plus efficace qu'un pamphlet.
\end{exoresolu}

\courssub{Méthode 5 : Assurer la liaison dans la phrase et entre les phrases}

\begin{methode}[Bien utiliser les outils de liaison]
\begin{enumerate}
\item \textbf{Conjonctions de coordination} (mais, ou, et, donc, or, ni, car) : relient deux éléments de même nature. « \emph{Car} » introduit une cause, « \emph{donc} » une conséquence.
\item \textbf{Conjonctions de subordination} : cause (parce que, puisque), conséquence (si bien que, de sorte que), concession (bien que, quoique + subjonctif), but (pour que, afin que), condition (si), temps (quand, lorsque).
\item \textbf{Adverbes et locutions de liaison} entre phrases et paragraphes : d'abord, ensuite, enfin ; en effet, en réalité ; cependant, néanmoins ; par conséquent, c'est pourquoi.
\item \textbf{Pronoms relatifs} (qui, que, où, dont, lequel) : évitent les répétitions et créent des phrases complexes claires.
\item \textbf{Vérifier la logique} : chaque connecteur doit correspondre au rapport logique réel (ne pas écrire « donc » pour une opposition).
\end{enumerate}
\end{methode}

\begin{exoresolu}[Exercice de liaison]
\textbf{Énoncé.} Complétez avec le connecteur logique approprié et justifiez : (1) « Meka est fier de sa médaille, \dots{} elle symbolise à ses yeux la reconnaissance des Blancs. » (2) « \dots{} Meka ait été décoré, il passe la nuit au poste de police. » (3) « La médaille est un honneur, \dots{} elle ne protège pas Meka de l'humiliation. » (4) « Les villageois admirent Meka ; \dots{}, certains envient sa distinction. »
\tcblower
\textbf{Solution détaillée.}
\begin{enumerate}
\item « Meka est fier de sa médaille, \textbf{car} elle symbolise à ses yeux la reconnaissance des Blancs. » Justification : la seconde proposition donne la \emph{cause} de la fierté ; « car » (coordination) convient. « Parce que » serait possible mais exige une subordonnée.
\item « \textbf{Bien que} Meka ait été décoré, il passe la nuit au poste de police. » Justification : rapport de \emph{concession} (opposition admise) ; « bien que » exige le subjonctif (« ait été »), ce qui est respecté.
\item « La médaille est un honneur, \textbf{mais} elle ne protège pas Meka de l'humiliation. » Justification : \emph{opposition} directe entre deux propositions coordonnées ; « mais » est le connecteur attendu. « Cependant » conviendrait aussi en tête de seconde phrase.
\item « Les villageois admirent Meka ; \textbf{cependant}, certains envient sa distinction. » Justification : l'envie \emph{nuance} l'admiration générale ; « cependant » (adverbe de liaison) marque cette restriction. « En effet » serait faux, car il n'y a pas explication mais opposition.
\end{enumerate}
\textbf{Conclusion} : le choix du connecteur dépend toujours du rapport logique (cause, opposition, concession, conséquence) entre les idées, et non de l'habitude.
\end{exoresolu}

\courssub{Méthode 6 : Rédiger la conclusion}

\begin{methode}[Les 3 temps de la conclusion]
\begin{enumerate}
\item \textbf{Le bilan} : rappeler brièvement le chemin parcouru (« Nous avons montré que..., puis que... ») sans recopier l'introduction.
\item \textbf{La réponse à la problématique} : trancher clairement la question posée, avec nuance si le plan est dialectique.
\item \textbf{L'ouverture} : élargir vers une question nouvelle, une réalité plus vaste ou une autre œuvre — sans commencer un nouveau développement.
\end{enumerate}
\textbf{Interdits} : introduire un argument nouveau, écrire « à mon avis » sans justification, terminer brutalement.
\textbf{Longueur} : 5 à 8 lignes.
\end{methode}

\begin{exoresolu}[Rédaction d'une conclusion complète]
\textbf{Énoncé.} Rédigez la conclusion de la dissertation : « La distinction sociale est-elle toujours source de bonheur ? »
\tcblower
\textbf{Solution détaillée} (modèle rédigé et commenté).

\emph{Bilan} : « Au terme de cette réflexion, nous avons vu que la distinction sociale flatte l'amour-propre et confère un prestige réel, mais que la médaille de Meka révèle aussi sa face sombre : instrument de propagande, elle n'empêche ni l'humiliation ni la désillusion. »

\emph{Réponse à la problématique} : « La distinction sociale n'est donc pas toujours source de bonheur : elle ne rend heureux que lorsqu'elle reconnaît la dignité véritable de l'homme, et non lorsqu'elle récompense sa soumission. »

\emph{Ouverture} : « Cette leçon dépasse le cadre colonial : dans nos sociétés actuelles, où titres et honneurs fascinent encore, ne convient-il pas de s'interroger sur la valeur réelle des distinctions que l'on recherche ? »

\textbf{Vérification} : les trois temps sont présents ; la réponse reprend le mot-clé « toujours » du sujet ; l'ouverture élargit sans ouvrir un nouveau développement.
\end{exoresolu}

\courssub{Méthode 7 : Confronter et améliorer sa production}

\begin{methode}[Relecture et amélioration en 5 points]
\begin{enumerate}
\item \textbf{Respect du sujet} : vérifier que chaque partie répond à la problématique (chasser le hors-sujet).
\item \textbf{Structure} : contrôler la présence de l'introduction en 4 temps, des paragraphes I-A-E-E, de la conclusion en 3 temps.
\item \textbf{Liaison} : souligner les connecteurs ; s'il y en a moins d'un par phrase, en ajouter ; vérifier leur exactitude logique.
\item \textbf{Langue} : corriger orthographe, accords (sujet-verbe, participe passé), conjugaison (présent de vérité générale), ponctuation.
\item \textbf{Échange} : lire la copie d'un camarade avec une grille (sujet, plan, arguments, exemples, langue) et formuler deux conseils précis ; accepter la critique sur sa propre copie et réécrire les passages faibles.
\end{enumerate}
\end{methode}

\begin{exoresolu}[Correction d'un passage défectueux]
\textbf{Énoncé.} Voici le début d'un devoir d'élève : « La médaille c'est quelque chose de bien. Meka il est content. Moi je pense que la médaille rend heureux. Dans le roman Meka reçoit une médaille. » Relevez quatre défauts et proposez une version améliorée.
\tcblower
\textbf{Solution détaillée.}
\begin{enumerate}
\item \textbf{Défaut 1 — absence d'amorce et de problématique} : le texte entre brutalement dans le sujet.
\item \textbf{Défaut 2 — avis personnel prématuré} : « Moi je pense » n'a pas sa place dans l'introduction.
\item \textbf{Défaut 3 — phrases juxtaposées sans connecteurs} : aucun outil de liaison.
\item \textbf{Défaut 4 — langue relâchée et répétitions} : « c'est quelque chose de bien », « Meka il est content », reprise de « médaille ».
\end{enumerate}
\textbf{Version améliorée} : « Les honneurs occupent une place importante dans la vie des hommes, car chacun aspire à être reconnu par la société. C'est ce que montre Ferdinand Oyono dans \emph{Le Vieux nègre et la médaille}, où le vieux Meka, décoré par l'administration coloniale, éprouve d'abord une fierté immense. Cependant, cet honneur débouche sur l'humiliation. On peut donc se demander si la distinction sociale est toujours source de bonheur. »

\textbf{Conclusion} : la version corrigée respecte les 4 temps de l'introduction, utilise des connecteurs logiques (car, cependant, donc) et bannit l'avis personnel.
\end{exoresolu}

\begin{attention}[Les 5 erreurs qui coûtent des points au Bac]
\begin{enumerate}
\item \textbf{Le hors-sujet} : réciter l'œuvre ou parler du thème en général sans répondre à la question précise posée. Chaque paragraphe doit pouvoir être rattaché à la problématique.
\item \textbf{Raconter l'histoire au lieu d'argumenter} : résumer les péripéties du \emph{Vieux nègre et la médaille} ne vaut rien si la scène citée ne sert pas un argument clairement formulé (règle I-A-E-E).
\item \textbf{Donner son avis dans l'introduction} ou oublier d'annoncer le plan : l'introduction doit présenter le sujet et le plan, jamais la conclusion personnelle.
\item \textbf{Les connecteurs logiques faux ou absents} : écrire « donc » pour une opposition, juxtaposer les phrases sans liaison, ou oublier le subjonctif après « bien que » : la logique et la grammaire des connecteurs sont sanctionnées.
\item \textbf{Négliger la relecture} : fautes d'accord, oubli de la conclusion ou de l'ouverture, absence de transition entre les parties. Réserver toujours 10 à 15 minutes à la relecture complète de la copie.
\end{enumerate}
\end{attention}

\newpage

\coursec{Exercices}

\courssub{Je vérifie mes connaissances}

\begin{exercice}[QCM : les étapes de la dissertation]
Pour chaque question, choisis la (ou les) bonne(s) réponse(s).
\begin{enumerate}
\item L'introduction d'une dissertation comprend, dans l'ordre :
\begin{enumerate}
\item[a)] problématique -- amorce -- annonce du plan -- présentation du sujet ;
\item[b)] amorce -- présentation du sujet -- problématique -- annonce du plan ;
\item[c)] annonce du plan -- amorce -- problématique -- présentation du sujet.
\end{enumerate}
\item Un paragraphe argumentatif bien construit contient :
\begin{enumerate}
\item[a)] une idée directrice, un argument, un exemple, une phrase de transition ;
\item[b)] uniquement des exemples ;
\item[c)] une citation et une conclusion générale.
\end{enumerate}
\item La conclusion d'une dissertation :
\begin{enumerate}
\item[a)] introduit un nouvel argument ;
\item[b)] répond à la problématique et peut s'ouvrir sur une question nouvelle ;
\item[c)] répète mot pour mot l'introduction.
\end{enumerate}
\item Dans \emph{Le Vieux nègre et la médaille} de Ferdinand Oyono, la médaille remise à Meka symbolise :
\begin{enumerate}
\item[a)] la réussite sociale du vieillard ;
\item[b)] l'hypocrisie et l'illusion de la récompense coloniale ;
\item[c)] la richesse matérielle du village.
\end{enumerate}
\end{enumerate}
\end{exercice}

\begin{exercice}[Vrai ou faux, en justifiant]
Dis si chaque affirmation est vraie ou fausse, puis justifie ta réponse en une ou deux phrases.
\begin{enumerate}
\item L'annonce du plan doit citer explicitement les mots « premièrement, deuxièmement ».
\item Une conjonction de subordination relie une proposition principale à une proposition subordonnée.
\item La problématique est une question (ou un ensemble de questions) qui naît de la confrontation des idées du sujet.
\item Dans une dissertation, les exemples peuvent être tirés uniquement de l'œuvre imposée au programme.
\end{enumerate}
\end{exercice}

\begin{exercice}[Définitions]
Définis en deux ou trois phrases chacun des termes suivants, puis illustre chacun par un exemple de ton invention :
\begin{enumerate}
\item la \cle{thèse} ;
\item l'\cle{argument} ;
\item la \cle{problématique} ;
\item le \cle{connecteur logique} (ou outil de liaison).
\end{enumerate}
\end{exercice}

\begin{exercice}[Classement des outils de liaison]
Classe les mots et expressions suivants dans le tableau : \emph{parce que ; cependant ; qui ; d'abord ; mais ; en effet ; bien que ; ensuite ; où ; par conséquent}.
\begin{center}
\begin{tabularx}{\linewidth}{|l|X|X|X|}\hline
\rowcolor{popblueL} & \textbf{Conjonctions} & \textbf{Adverbes / locutions} & \textbf{Pronoms relatifs} \\ \hline
Coordination & & & \\ \hline
Subordination & & & \\ \hline
\end{tabularx}
\end{center}
\end{exercice}

\courssub{Je m'entraîne}

\begin{exercice}[Analyse d'un sujet]
Soit le sujet : « Les réseaux sociaux détruisent-ils la communication entre les jeunes Camerounais ? »
\begin{enumerate}
\item Souligne les mots-clés du sujet et définis chacun d'eux.
\item Dégage les deux thèses en présence.
\item Formule la problématique en une question claire.
\item Propose deux arguments (avec un exemple chacun) pour chacune des deux thèses.
\end{enumerate}
\end{exercice}

\begin{exercice}[Réécriture avec connecteurs]
Réécris le paragraphe suivant en y insérant au moins cinq outils de liaison appropriés (au moins une conjonction de subordination, un adverbe de liaison, une locution adverbiale et un pronom relatif). Tu peux modifier légèrement les phrases.

\emph{Meka a donné ses terres à la mission. Meka a donné ses fils à la guerre. Le commandant lui remet une médaille. Meka croit être devenu quelqu'un d'important. La médaille ne change rien à sa condition. Meka comprend qu'on s'est moqué de lui.}
\end{exercice}

\begin{exercice}[Lecture méthodique]
Lis attentivement l'extrait suivant du \emph{Vieux nègre et la médaille} (Ferdinand Oyono) : après la cérémonie, Meka, épuisé et humilié par la pluie, les moqueries et l'arrestation de la nuit, mesure le vide de la récompense reçue.

« Il avait donné ses terres, il avait donné ses fils... et voilà ce qu'on lui avait donné en échange : une médaille ! » (paraphrase de l'extrait étudié en classe).
\begin{enumerate}
\item Dégage la thèse que défend l'auteur dans cet épisode.
\item Relève deux stratégies argumentatives utilisées (ironie, accumulation, contraste, réquisitoire...) et justifie chacune par une courte citation.
\item Explique en cinq lignes comment cet épisode sert la critique de la colonisation.
\end{enumerate}
\end{exercice}

\begin{exercice}[Amélioration d'une copie]
Voici l'introduction rédigée par un élève pour le sujet : « Le roman ne fait que divertir le lecteur. »

\emph{Le roman est un genre littéraire. Beaucoup de gens lisent des romans. Je vais d'abord montrer que le roman divertit, ensuite qu'il instruit. Donc le roman est très intéressant pour tout le monde.}
\begin{enumerate}
\item Relève quatre faiblesses de cette introduction (amorce absente ou banale, problématique manquante, annonce du plan maladroite, conclusion prématurée...).
\item Réécris entièrement cette introduction en respectant les quatre étapes : amorce, présentation du sujet, problématique, annonce du plan.
\end{enumerate}
\end{exercice}

\courssub{Je me prépare au Bac}

\begin{exercice}[Sujet de dissertation type Bac]
\textbf{Sujet :} « Le roman ne fait que divertir le lecteur. » Discutez cette opinion en vous appuyant sur \emph{Le Vieux nègre et la médaille} de Ferdinand Oyono et sur votre culture personnelle.

\begin{enumerate}
\item Analyse le sujet : mots-clés, définitions, thèses en présence, problématique.
\item Élabore un plan détaillé (deux ou trois parties), avec pour chaque partie l'idée directrice, deux arguments et un exemple précis.
\item Rédige l'introduction complète (quinze à vingt lignes).
\item Rédige le premier paragraphe du développement (idée directrice, argument, exemple commenté, transition).
\item Rédige la conclusion (réponse à la problématique et ouverture).
\end{enumerate}
\emph{Barème indicatif : analyse du sujet et plan 4 pts ; introduction 6 pts ; paragraphe 6 pts ; conclusion 4 pts.}
\end{exercice}

\begin{exercice}[Dissertation complète contextualisée]
\textbf{Sujet :} Dans le monde d'aujourd'hui, les jeunes affirment souvent que « les récompenses et les honneurs ne sont que des apparences trompeuses ». Partagez-vous ce point de vue ? Vous développerez votre réponse en vous appuyant sur \emph{Le Vieux nègre et la médaille} et sur des exemples tirés de la vie sociale camerounaise.

\begin{enumerate}
\item Dégage la problématique et annonce un plan dialectique (thèse / antithèse / synthèse).
\item Rédige la dissertation complète (introduction, développement en trois paragraphes, conclusion), en veillant à utiliser au moins huit outils de liaison variés, que tu souligneras.
\item À la fin de ta copie, dresse la liste des connecteurs utilisés en précisant leur fonction (opposition, cause, conséquence, addition...).
\end{enumerate}
\emph{Durée conseillée : 2 h 30. Longueur attendue : deux à trois pages.}
\end{exercice}

\begin{exercice}[Confrontation et amélioration de productions]
Voici le paragraphe rédigé par un candidat au Baccalauréat technique :

\emph{La colonisation est mauvaise. Meka reçoit une médaille. Cette médaille ne sert à rien. Oyono critique les Blancs. Donc la colonisation est une injustice pour les Africains.}

\begin{enumerate}
\item Relève les défauts de ce paragraphe (enchaînement, absence d'idée directrice claire, exemple non commenté, argumentation pauvre, connecteurs absents).
\item Réécris ce paragraphe en suivant le schéma : idée directrice -- argument -- exemple cité et commenté -- phrase de conclusion-transition, avec au moins quatre outils de liaison.
\item Compare ta version avec celle d'un camarade : note sur 5 la clarté de l'idée directrice, la pertinence de l'exemple et la qualité des liaisons, puis propose une amélioration supplémentaire.
\end{enumerate}
\end{exercice}

\newpage

\coursec{Corrigés des exercices}

\coursec{Lecture méthodique du \emph{Vieux nègre et la médaille} (1re partie)}

\begin{corrige}[Exercice 7 — Lecture méthodique : thèse, stratégies et critique]
\begin{enumerate}
\item \textbf{Thèse de l'auteur :} Ferdinand Oyono défend l'idée que la colonisation repose sur le mensonge et l'exploitation : les récompenses coloniales (la médaille) ne sont qu'un leurre destiné à faire accepter aux Africains leurs sacrifices, sans jamais les dédommager réellement.
\item \textbf{Deux stratégies argumentatives :}
\begin{itemize}
\item \emph{L'accumulation} : « il avait donné ses terres, il avait donné ses fils » — la répétition de la structure « il avait donné » souligne l'immensité des sacrifices de Meka et prépare le contraste avec la récompense dérisoire.
\item \emph{L'ironie et le contraste} : « et voilà ce qu'on lui avait donné en échange : une médaille ! » — le point d'exclamation et le mot « échange » dénoncent l'inégalité scandaleuse entre le don immense (terres, fils) et le don ridicule (un bout de métal).
\end{itemize}
\item \textbf{Critique de la colonisation (environ cinq lignes) :} Cet épisode condense la dénonciation de tout le système colonial. Meka a tout perdu — sa terre, source de vie, et ses fils, morts pour la France — et l'administration croit l'honorer d'une médaille sans valeur. L'humiliation de la nuit (pluie, moqueries, arrestation) révèle la vérité : l'Africain « évolué » reste un sujet méprisé. Oyono montre ainsi que les honneurs coloniaux masquent la spoliation et l'injustice, et invite le lecteur à rejeter cette illusion.
\end{enumerate}
\end{corrige}

\coursec{Le texte argumentatif : stratégies argumentatives}

\begin{corrige}[Exercice 1 — QCM : les étapes de la dissertation et structure du paragraphe]
\begin{enumerate}
\item \textbf{Réponse b)} : amorce -- présentation du sujet -- problématique -- annonce du plan. C'est l'ordre canonique exigé au Probatoire / Baccalauréat technique : on part du général (amorce), on resserre sur le sujet, on pose la question directrice, puis on annonce le plan.
\item \textbf{Réponse a)} : un paragraphe argumentatif suit le schéma \cle{idée directrice} -- \cle{argument} -- \cle{exemple} commenté -- phrase de transition. Les réponses b) et c) sont fausses : un paragraphe ne peut être une simple accumulation d'exemples, et la conclusion générale n'appartient pas au paragraphe.
\item \textbf{Réponse b)} : la conclusion \textbf{répond à la problématique} (bilan) et peut s'ouvrir sur une question nouvelle. Elle ne doit jamais introduire un nouvel argument (a) ni répéter l'introduction (c).
\item \textbf{Réponse b)} : la médaille symbolise l'\cle{hypocrisie} et l'illusion de la récompense coloniale : Meka a donné ses terres et ses fils, et ne reçoit en échange qu'un bout de métal qui ne change rien à sa condition.
\end{enumerate}
\textbf{Points de vigilance :} au Bac, l'ordre des étapes de l'introduction est sanctionné ; toute inversion (problématique avant la présentation du sujet, par exemple) est pénalisée.
\end{corrige}

\begin{corrige}[Exercice 2 — Vrai ou faux, en justifiant]
\begin{enumerate}
\item \textbf{Faux.} L'annonce du plan doit être \emph{élégante et implicite} : on annonce les deux (ou trois) axes dans une phrase fluide (« Nous verrons d'abord que..., avant de montrer que... »), sans énumération scolaire du type « premièrement, deuxièmement ».
\item \textbf{Vrai.} C'est la définition même de la conjonction de subordination (\emph{parce que, bien que, puisque, si...}) : elle introduit une proposition subordonnée qui dépend d'une proposition principale.
\item \textbf{Vrai.} La problématique naît de la \emph{tension} entre les idées du sujet (par exemple : divertir / instruire) ; elle s'exprime par une question directrice qui oriente tout le devoir.
\item \textbf{Faux.} Les exemples peuvent et doivent être variés : œuvre au programme, mais aussi culture personnelle, autres œuvres littéraires, faits historiques, vie sociale camerounaise. Se limiter à une seule œuvre appauvrit l'argumentation.
\end{enumerate}
\end{corrige}

\begin{corrige}[Exercice 3 — Définitions des notions clés de l'argumentation]
\begin{enumerate}
\item \textbf{La thèse} : c'est le point de vue, l'opinion que l'on défend ou que l'on examine dans un texte argumentatif. Elle peut être soutenue par l'auteur ou attribuée à un adversaire. \emph{Exemple :} « Le travail manuel est aussi noble que le travail intellectuel » est une thèse.
\item \textbf{L'argument} : c'est une raison avancée pour justifier ou réfuter une thèse. Il doit être clair, pertinent et développé, puis illustré par un exemple. \emph{Exemple :} pour défendre que la lecture instruit, on peut arguer qu'elle transmet des connaissances sur d'autres cultures.
\item \textbf{La problématique} : c'est la question (ou l'ensemble de questions) qui naît de la confrontation des idées du sujet et qui guide l'ensemble du devoir. \emph{Exemple :} « Le roman a-t-il pour seule fonction de divertir, ou peut-il aussi instruire et dénoncer ? »
\item \textbf{Le connecteur logique} : c'est un mot ou une expression (conjonction, adverbe, locution, pronom relatif) qui relie les phrases et les idées en marquant une relation logique (cause, opposition, conséquence, addition...). \emph{Exemple :} « Meka a tout sacrifié ; \textbf{cependant}, il n'a reçu qu'une médaille dérisoire » (\emph{cependant} marque l'opposition).
\end{enumerate}
\end{corrige}

\coursec{Les outils de liaison dans la phrase}

\begin{corrige}[Exercice 4 — Classement des outils de liaison]
\begin{center}
\begin{tabularx}{\linewidth}{|l|X|X|X|}\hline
\rowcolor{popblueL} & \textbf{Conjonctions} & \textbf{Adverbes / locutions} & \textbf{Pronoms relatifs} \\ \hline
Coordination & mais & d'abord, ensuite, en effet, par conséquent, cependant & --- \\ \hline
Subordination & parce que, bien que & --- & qui, où \\ \hline
\end{tabularx}
\end{center}
\textbf{Justifications :}
\begin{itemize}
\item \emph{mais} : conjonction de coordination (opposition) ; elle relie deux propositions de même niveau.
\item \emph{parce que} (cause) et \emph{bien que} (concession) : conjonctions de subordination, elles introduisent une subordonnée.
\item \emph{d'abord, ensuite} : adverbes de liaison (ordre, énumération) ; \emph{en effet} : locution adverbiale d'explication ; \emph{par conséquent} : locution de conséquence ; \emph{cependant} : adverbe d'opposition.
\item \emph{qui} et \emph{où} : pronoms relatifs, ils introduisent une subordonnée relative et remplacent un nom antécédent.
\end{itemize}
\textbf{Point de vigilance :} \emph{cependant} n'est pas une conjonction de coordination : c'est un adverbe de liaison, souvent précédé d'un point-virgule.
\end{corrige}

\begin{corrige}[Exercice 6 — Réécriture d'un texte avec connecteurs variés]
\textbf{Version possible} (les outils de liaison sont en gras) :

\emph{Meka a donné ses terres à la mission \textbf{et} il a même offert ses fils à la guerre. \textbf{C'est pourquoi} le commandant, \textbf{lors de la grande cérémonie}, lui remet une médaille. Meka, \textbf{qui} a tant sacrifié, croit \textbf{d'abord} être devenu quelqu'un d'important. \textbf{Cependant}, la médaille ne change rien à sa condition, \textbf{bien qu'}elle brille sur sa poitrine. \textbf{En effet}, ce bout de métal ne lui rend ni ses terres ni ses fils. \textbf{Finalement}, Meka comprend qu'on s'est moqué de lui.}

\textbf{Connecteurs utilisés (au moins cinq exigés) :} \emph{et} (coordination, addition) ; \emph{c'est pourquoi} (conséquence) ; \emph{lors de la grande cérémonie} (locution de temps) ; \emph{qui} (pronom relatif) ; \emph{d'abord} (adverbe d'ordre) ; \emph{cependant} (adverbe d'opposition) ; \emph{bien que} (conjonction de subordination, concession) ; \emph{en effet} (locution d'explication) ; \emph{finalement} (adverbe de conclusion).

\textbf{Point de vigilance :} vérifier que chaque connecteur exprime bien la relation logique voulue ; un connecteur mal choisi (par exemple \emph{par conséquent} pour une opposition) fausse le sens.
\end{corrige}

\coursec{Rédaction de l'introduction et de la conclusion}

\begin{corrige}[Exercice 8 — Amélioration d'une copie : l'introduction]
\begin{enumerate}
\item \textbf{Quatre faiblesses de l'introduction proposée :}
\begin{itemize}
\item amorce absente et banale : « Le roman est un genre littéraire » est une définition plate, sans ouverture sur le monde ou la vie des hommes ;
\item problématique manquante : aucune question ne naît de la confrontation des idées (divertir / instruire) ;
\item annonce du plan maladroite et scolaire : « Je vais d'abord montrer... ensuite... » avec la première personne, à bannir ;
\item conclusion prématurée : « Donc le roman est très intéressant » juge avant d'avoir argumenté, ce qui revient à plaquer la conclusion dans l'introduction.
\end{itemize}
\item \textbf{Introduction réécrite (modèle) :}

\emph{Depuis toujours, les hommes racontent des histoires : autour du feu au village, dans les livres, à l'écran. Le roman, genre né de ce besoin de récits, occupe une place centrale dans la vie culturelle, en Europe comme au Cameroun. Certains lecteurs estiment pourtant qu'il ne sert qu'à divertir, à faire passer le temps. Cette opinion est-elle fondée ? Le roman n'a-t-il pour fonction que le plaisir du lecteur, ou peut-il aussi instruire, dénoncer et transformer la société ? Nous verrons d'abord que le roman est bien une source de divertissement, puis nous montrerons qu'il est également un instrument de connaissance et de critique sociale.}
\end{enumerate}
\textbf{Points de vigilance :} bannir la première personne (« je vais montrer ») ; ne jamais conclure dans l'introduction ; la problématique doit être une vraie question.
\end{corrige}

\coursec{Rédaction du paragraphe argumentatif}

\begin{corrige}[Exercice 9 — Sujet de dissertation type Bac : analyse, plan, introduction, paragraphe, conclusion]
\textbf{Barème indicatif :} analyse du sujet et plan 4 pts ; introduction 6 pts ; paragraphe 6 pts ; conclusion 4 pts.

\begin{enumerate}
\item \textbf{Analyse du sujet (1,5 pt) :} Mots-clés : \emph{roman} (genre narratif en prose), \emph{ne... que} (restriction : c'est le cœur du débat), \emph{divertir} (distraire, faire passer un moment agréable), \emph{lecteur}. Thèse implicite du sujet : le roman ne sert qu'à divertir. Thèse opposée : le roman a d'autres fonctions (instruire, dénoncer, éduquer). \textbf{Problématique :} le roman se limite-t-il au divertissement, ou remplit-il aussi des fonctions de connaissance et de critique sociale ?

\item \textbf{Plan détaillé dialectique (2,5 pts) :}

\textbf{I. Le roman divertit le lecteur.}
Idée directrice : le plaisir du récit est la première fonction du roman.
Argument 1 : le roman fait vivre des aventures par procuration ; \emph{exemple :} le lecteur suit les péripéties de Meka le jour de la cérémonie avec suspense et amusement.
Argument 2 : l'humour et la satire procurent le rire ; \emph{exemple :} les scènes comiques du \emph{Vieux nègre et la médaille} (le commandant, le père Vandermayer) font sourire.

\textbf{II. Mais le roman fait bien plus : il instruit et dénonce.}
Idée directrice : le roman est un miroir et une arme critique.
Argument 1 : il fait connaître une réalité historique et sociale ; \emph{exemple :} Oyono fait découvrir la vie des villages camerounais sous la colonisation.
Argument 2 : il dénonce l'injustice et éveille les consciences ; \emph{exemple :} la médaille de Meka dénonce l'hypocrisie coloniale ; de même, la littérature engagée africaine a nourri les luttes d'indépendance.

\item \textbf{Introduction rédigée (6 pts) :}

\emph{L'homme est un être d'histoires : des veillées du village africain aux bibliothèques des grandes villes, le récit accompagne la vie des hommes. Le roman, héritier moderne de cette tradition, est devenu le genre littéraire le plus lu au monde, y compris chez les jeunes Camerounais. Pourtant, certains esprits pressés affirment qu'il ne fait que divertir le lecteur, comme un simple passe-temps sans portée. Cette opinion réductrice mérite d'être discutée : le roman se limite-t-il au plaisir immédiat, ou possède-t-il aussi le pouvoir d'instruire, de dénoncer et de transformer celui qui le lit ? Nous montrerons d'abord que le roman est bien une source de divertissement, avant d'établir qu'il est également un instrument de connaissance et de critique sociale.}

\item \textbf{Premier paragraphe du développement (6 pts) :}

\emph{Il est indéniable, d'abord, que le roman divertit le lecteur. En effet, par le jeu du récit, il arrache celui-ci à sa vie quotidienne et le transporte dans un monde imaginaire peuplé de personnages attachants. Ainsi, dans \emph{Le Vieux nègre et la médaille}, Ferdinand Oyono entraîne son lecteur à la suite de Meka, le vieux villageois convoqué à Douala pour recevoir une médaille : les préparatifs de la cérémonie, l'agitation du village, les réactions comiques des notables blancs tiennent le lecteur en haleine et le font sourire. Le plaisir de la lecture naît donc de cette aventure vécue par procuration. Cependant, réduire le roman à ce seul divertissement serait méconnaître sa dimension profonde, comme nous allons le voir.}

\item \textbf{Conclusion (4 pts) :}

\emph{En définitive, le roman divertit certes le lecteur, mais il ne fait pas que cela. Comme le montre \emph{Le Vieux nègre et la médaille}, il instruit sur l'histoire et les sociétés, dénonce les injustices et invite à la réflexion. L'opinion examinée est donc trop étroite : le roman est à la fois un plaisir et une leçon. On peut dès lors se demander si, à l'heure des écrans et des réseaux sociaux, le roman conserve encore ce pouvoir auprès des jeunes générations.}
\end{enumerate}
\textbf{Points de vigilance :} la problématique doit figurer dans l'introduction ; chaque exemple doit être \emph{commenté} (un exemple cité sans explication ne vaut rien) ; la conclusion répond explicitement à la problématique, sans nouvel argument.
\end{corrige}

\begin{corrige}[Exercice 11 — Confrontation et amélioration de productions : le paragraphe]
\begin{enumerate}
\item \textbf{Défauts du paragraphe d'élève :}
\begin{itemize}
\item absence d'idée directrice claire : « La colonisation est mauvaise » est un jugement vague, non une idée directrice développable ;
\item exemple non commenté : « Meka reçoit une médaille » est cité brut, sans explication de sa signification ;
\item argumentation pauvre et répétitive : « Cette médaille ne sert à rien » répète l'idée sans l'approfondir ;
\item connecteurs absents : un seul « Donc » mal utilisé (la conclusion ne découle pas logiquement de ce qui précède) ;
\item enchaînement saccadé : phrases courtes juxtaposées, sans liant.
\end{itemize}
\item \textbf{Paragraphe réécrit (modèle) :}

\emph{La colonisation, loin d'être une œuvre civilisatrice, repose sur l'injustice et le mépris des Africains. \textbf{En effet}, elle spolie les hommes de leurs biens tout en prétendant les honorer. \textbf{Ainsi}, dans \emph{Le Vieux nègre et la médaille}, Meka, \textbf{qui} a sacrifié ses terres et ses fils pour la France, reçoit en récompense une simple médaille : ce bout de métal, \textbf{qui} ne lui rend rien de ce qu'il a perdu, révèle l'hypocrisie du système colonial. \textbf{Cependant}, l'humiliation de la nuit qui suit la cérémonie achève de lui ouvrir les yeux. \textbf{Par conséquent}, à travers le destin de ce vieillard, Oyono dénonce la colonisation comme une injustice profonde, ce que confirmera l'analyse des autres épisodes du roman.}

Connecteurs utilisés : \emph{en effet} (explication), \emph{ainsi} (illustration/conséquence), \emph{qui} (pronoms relatifs), \emph{cependant} (opposition), \emph{par conséquent} (conclusion).
\item \textbf{Grille de comparaison entre pairs (sur 5) :}
\begin{center}
\begin{tabularx}{\linewidth}{|X|c|}\hline
\rowcolor{popblueL} \textbf{Critère} & \textbf{Note /5} \\ \hline
Clarté de l'idée directrice (phrase d'ouverture qui annonce l'argument) & \\ \hline
Pertinence de l'exemple (précis, tiré de l'œuvre, cité et commenté) & \\ \hline
Qualité des liaisons (au moins quatre connecteurs variés et corrects) & \\ \hline
\end{tabularx}
\end{center}
\emph{Amélioration supplémentaire possible :} ajouter une courte citation exacte du roman pour renforcer l'exemple, et varier la longueur des phrases pour le rythme.
\end{enumerate}
\textbf{Points de vigilance :} la critique d'une copie doit être argumentée (nommer précisément le défaut) ; la réécriture doit respecter le schéma imposé ; l'évaluation par les pairs développe l'esprit critique exigé au Probatoire / Baccalauréat technique.
\end{corrige}

\coursec{Analyse du sujet, plan détaillé et production complète}

\begin{corrige}[Exercice 5 — Analyse d'un sujet d'actualité : réseaux sociaux et communication]
\begin{enumerate}
\item \textbf{Mots-clés et définitions :}
\begin{itemize}
\item \emph{réseaux sociaux} : plateformes numériques (WhatsApp, Facebook, TikTok...) permettant d'échanger messages, images et vidéos ;
\item \emph{détruisent} : verbe fort, qui suppose une dégradation, une disparition ;
\item \emph{communication} : échange d'informations, de sentiments, de paroles entre personnes ;
\item \emph{jeunes Camerounais} : catégorie sociale précise, ancrée dans un contexte national (familles, quartiers, écoles).
\end{itemize}
\item \textbf{Les deux thèses en présence :}
\begin{itemize}
\item Thèse 1 : les réseaux sociaux détruisent la communication (isolement, dialogues superficiels, disparition des veillées et des échanges directs) ;
\item Thèse 2 : les réseaux sociaux ne détruisent pas la communication, ils la transforment et l'élargissent (lien avec les proches éloignés, information rapide).
\end{itemize}
\item \textbf{Problématique :} « Les réseaux sociaux détruisent-ils réellement la communication entre les jeunes Camerounais, ou ne font-ils que la transformer ? »
\item \textbf{Arguments et exemples :}
\begin{itemize}
\item \emph{Thèse 1, argument 1 :} les échanges virtuels remplacent les conversations directes ; \emph{exemple :} des amis assis ensemble dans un quartier de Yaoundé, chacun absorbé par son téléphone.
\item \emph{Thèse 1, argument 2 :} la communication devient superficielle ; \emph{exemple :} les « likes » et messages abrégés remplacent les longues discussions familiales.
\item \emph{Thèse 2, argument 1 :} les réseaux maintiennent le lien à distance ; \emph{exemple :} un étudiant de Bafoussam échange chaque soir par WhatsApp avec sa famille restée au village.
\item \emph{Thèse 2, argument 2 :} ils diffusent l'information utile ; \emph{exemple :} les groupes scolaires où les élèves partagent les leçons et les sujets d'examen.
\end{itemize}
\end{enumerate}
\end{corrige}

\begin{corrige}[Exercice 10 — Dissertation complète contextualisée : honneurs et apparences]
\begin{enumerate}
\item \textbf{Problématique :} les récompenses et les honneurs ne sont-ils que des apparences trompeuses, ou conservent-ils une valeur réelle ? \textbf{Plan dialectique :} I. Les honneurs sont souvent des leurres (thèse des jeunes) ; II. Pourtant, certaines récompenses ont une valeur authentique ; III. Synthèse : il faut distinguer les honneurs sincères des honneurs manipulateurs.

\item \textbf{Dissertation rédigée} (connecteurs soulignés dans la copie ; ici en gras) :

\emph{\textbf{Introduction.} Dans nos sociétés, on couvre de médailles, de titres et de décorations ceux que l'on veut honorer : chefs, fonctionnaires, anciens combattants. \textbf{Cependant}, beaucoup de jeunes affirment aujourd'hui que ces récompenses ne sont que des apparences trompeuses. Ce point de vue est-il fondé ? Les honneurs sont-ils toujours des leurres, ou peuvent-ils avoir une valeur réelle ? Nous verrons \textbf{d'abord} que les récompenses peuvent tromper, \textbf{ensuite} qu'elles gardent parfois un sens authentique, \textbf{enfin} qu'il appartient à chacun de distinguer le vrai honneur de l'illusion.}

\emph{\textbf{I.} \textbf{En effet}, l'expérience montre que les honneurs servent souvent à masquer l'injustice. \textbf{Par exemple}, dans \emph{Le Vieux nègre et la médaille}, Meka, \textbf{qui} a donné ses terres et ses fils à la France, ne reçoit \textbf{en échange} qu'une médaille dérisoire : l'honneur colonial cache la spoliation. \textbf{De même}, dans la vie sociale camerounaise, certaines distinctions récompensent la complaisance plutôt que le mérite. Les jeunes ont donc des raisons de se méfier.}

\emph{\textbf{II.} \textbf{Néanmoins}, toutes les récompenses ne sont pas trompeuses. \textbf{Parce qu'}elles reconnaissent un effort réel, elles peuvent encourager et élever celui \textbf{qui} les reçoit : le diplôme du bachelier, le trophée du sportif, la distinction du travailleur méritant ont une valeur concrète. \textbf{Ainsi}, l'honneur sincère reste un moteur du progrès individuel et collectif.}

\emph{\textbf{III.} \textbf{Par conséquent}, le vrai problème n'est pas la récompense elle-même, \textbf{mais} l'intention de celui \textbf{qui} la donne. \textbf{Bien que} la médaille de Meka soit un leurre, d'autres honneurs, décernés avec justice, méritent le respect. \textbf{Finalement}, c'est au regard lucide de chacun — comme celui que Meka finit par ouvrir — de démasquer l'apparence et de reconnaître la valeur véritable.}

\emph{\textbf{Conclusion.} Les jeunes ont donc partiellement raison : les honneurs peuvent tromper, comme le prouve l'expérience amère de Meka. \textbf{Toutefois}, ils ne sont pas tous des apparences : décernés avec justice, ils reconnaissent le mérite. Reste à savoir si nos sociétés modernes sauront rendre aux récompenses leur sincérité perdue.}

\item \textbf{Liste des connecteurs utilisés et leur fonction :}
\begin{center}
\begin{tabularx}{\linewidth}{|l|X|}\hline
\rowcolor{popblueL} \textbf{Connecteur} & \textbf{Fonction logique} \\ \hline
cependant / néanmoins / toutefois & opposition \\ \hline
d'abord, ensuite, enfin, finalement & ordre, énumération \\ \hline
en effet & explication, justification \\ \hline
parce que & cause \\ \hline
par conséquent / ainsi & conséquence \\ \hline
mais & opposition (coordination) \\ \hline
bien que & concession (subordination) \\ \hline
qui & pronom relatif (reprise de l'antécédent) \\ \hline
par exemple / de même & illustration, comparaison \\ \hline
\end{tabularx}
\end{center}
\end{enumerate}
\textbf{Points de vigilance :} au moins huit connecteurs \emph{variés} et soulignés ; respecter le plan dialectique annoncé ; chaque paragraphe suit le schéma idée directrice -- argument -- exemple commenté ; la conclusion répond à la problématique sans nouvel argument.
\end{corrige}

\newpage

\coursec{Fiche bilan}

\begin{popfigure}
\centering
\begin{tikzpicture}[mindmap, grow cyclic,
  every node/.style={concept, text=white, font=\small},
  root concept/.append style={concept color=popdark, font=\bfseries},
  level 1/.append style={level distance=3.4cm, sibling angle=60},
  level 1 concept/.append style={font=\footnotesize}]
\node[root concept]{La dissertation}
  child[concept color=popblue]{ node[align=center]{Lecture méthodique\\ \emph{Le Vieux nègre}\\ \emph{et la médaille}} }
  child[concept color=poporange]{ node[align=center]{Texte argumentatif\\ stratégies : convaincre,\\ persuader, délibérer} }
  child[concept color=popgreen]{ node[align=center]{Outils de liaison\\ conjonctions, adverbes,\\ pronoms relatifs} }
  child[concept color=poppurple]{ node[align=center]{Introduction\\ accroche, sujet,\\ problématique, plan} }
  child[concept color=poppink]{ node[align=center]{Paragraphe\\ idée + argument\\ + exemple} }
  child[concept color=popgold]{ node[align=center]{Conclusion\\ bilan + ouverture} }
  child[concept color=popmuted]{ node[align=center]{Plan détaillé\\ analyse du sujet,\\ production, relecture} };
\end{tikzpicture}
\legende{Carte mentale de la leçon : les sept savoirs à mobiliser pour rédiger une dissertation au Baccalauréat technique.}
\end{popfigure}

\begin{bilan}
\textbf{1. Définitions clés}
\begin{itemize}
  \item \cle{Dissertation} : exercice écrit qui expose de manière ordonnée et argumentée une réflexion personnelle sur un sujet donné.
  \item \cle{Problématique} : question centrale, issue de l'analyse du sujet, à laquelle toute la dissertation répond.
  \item \cle{Thèse} : point de vue défendu ; \cle{antithèse} : point de vue opposé ; \cle{synthèse} : dépassement des deux.
  \item \cle{Argument} : raison logique qui justifie une idée ; \cle{exemple} : fait précis (texte, histoire, actualité) qui illustre l'argument.
  \item \cle{Stratégies argumentatives} : convaincre (raison), persuader (émotions), délibérer (peser le pour et le contre).
\end{itemize}

\textbf{2. À connaître par c\oe ur}
\begin{center}
\begin{tabularx}{\linewidth}{|l|X|}
\hline
\rowcolor{popblueL} \textbf{Élément} & \textbf{Contenu obligatoire} \\
\hline
Introduction & Phrase d'accroche ; présentation du sujet ; problématique ; annonce du plan. \\
\hline
Paragraphe & Idée directrice ; argument ; exemple ; petite phrase de transition. \\
\hline
Conclusion & Bilan (réponse à la problématique) ; élargissement (ouverture). \\
\hline
Liaisons & Coordination : mais, ou, et, donc, or, ni, car. Subordination : parce que, bien que, si, lorsque. Adverbes : d'abord, ensuite, enfin, cependant, en effet. \\
\hline
\end{tabularx}
\end{center}

\textbf{3. Méthodes express}
\begin{itemize}
  \item \textbf{Analyser le sujet} : souligner les mots-clés, définir chaque terme, repérer les présupposés, formuler la problématique.
  \item \textbf{Mobiliser les lectures} : \emph{Le Vieux nègre et la médaille} de Ferdinand Oyono fournit des exemples (colonialisme, hypocrisie, dignité) réutilisables dans les développements.
  \item \textbf{Construire le plan détaillé} : au brouillon, noter pour chaque partie : idée, arguments, exemples, transitions.
  \item \textbf{Relire et améliorer} : confronter sa copie à celle d'un camarade, corriger orthographe, liaisons et cohérence.
\end{itemize}
\end{bilan}

\begin{aretenir}[Checklist avant le Bac]
Avant de rendre ma copie de dissertation, je vérifie :
\begin{itemize}
  \item[$\square$] J'ai défini tous les mots-clés du sujet au brouillon.
  \item[$\square$] Ma problématique est une vraie question, claire et liée au sujet.
  \item[$\square$] Mon introduction contient les quatre étapes : accroche, sujet, problématique, annonce du plan.
  \item[$\square$] Mon plan comporte au moins deux parties équilibrées et annoncées.
  \item[$\square$] Chaque paragraphe suit le schéma : idée directrice, argument, exemple.
  \item[$\square$] Mes exemples sont précis (œuvres lues, faits historiques, vie courante).
  \item[$\square$] J'ai utilisé des connecteurs variés (d'abord, en effet, cependant, enfin).
  \item[$\square$] Mes phrases sont bien liées par des conjonctions et des pronoms relatifs corrects.
  \item[$\square$] Je ne donne pas mon avis avec « je » de façon systématique ; j'argumente objectivement.
  \item[$\square$] Ma conclusion répond à la problématique et propose une ouverture.
  \item[$\square$] J'ai relu pour corriger l'orthographe, les accords et la ponctuation.
  \item[$\square$] Je n'ai pas recopié le sujet ni écrit hors sujet.
\end{itemize}
\end{aretenir}

\findecours

\end{document}
